% Las conductas de riesgo : sensibilizar para adoptar comportamientos responsables — Espagnol 1ère A — Cours signé OnBuch+
% Programme officiel MINESEC. Compiler avec : tectonic EA04-conductas-de-riesgo-sensibilizar.tex
\newcommand{\DOCMATIERE}{Espagnol}
\newcommand{\DOCNIVEAU}{1ère A}
\newcommand{\DOCMODULE}{Módulo 1 : Vie familiale et intégration sociale}
\newcommand{\DOCLECON}{EA04}
\newcommand{\DOCTITRE}{Las conductas de riesgo : sensibilizar para adoptar comportamientos responsables}
% OnBuch+ — préambule « cours signés » ENRICHI (compilé avec tectonic / XeLaTeX)
% Système visuel « Pop » d'OnBuch+ : fond crème (jamais blanc), bordures encre
% épaisses, ombres « dures » décalées, titres Archivo Black en capitales.
% Version MATHÉMATIQUES : graphes (pgfplots/tikz), boîtes pédagogiques
% (définition, propriété, méthode, attention, le savais-tu, exercice résolu,
% exercice, TP, bilan), page de garde et signature OnBuch+ ancrée en bas.
%
% Macros attendues dans main.tex AVANT \input{preamble} :
%   \DOCMATIERE  \DOCNIVEAU  \DOCMODULE  \DOCLECON (numéro)  \DOCTITRE
\documentclass[11pt]{article}

\usepackage{fontspec}
\usepackage[spanish,french]{babel} % français = langue principale ; l'espagnol via \esp{..} / espagnol
\usepackage[a4paper,margin=2cm,top=3.4cm,bottom=2.8cm,headheight=1.4cm,footskip=1.3cm]{geometry}
\usepackage{amsmath,amssymb,amsfonts}
\usepackage{mathtools} % \xmapsto, \coloneqq... souvent utilisés par les modèles
\usepackage{cancel} % simplifications barrées (bilans de forces)
% \abs{z} (module, valeur absolue) et \norm{u} (norme), utilisés par les modèles
\providecommand{\abs}{}\let\abs\relax\DeclarePairedDelimiter{\abs}{\lvert}{\rvert}
\providecommand{\norm}{}\let\norm\relax\DeclarePairedDelimiter{\norm}{\lVert}{\rVert}
\usepackage[table]{xcolor} % [table] : fournit aussi \rowcolors (alternance de lignes)
\usepackage{pagecolor}
\usepackage{tikz}
\usetikzlibrary{arrows.meta,positioning,calc,shapes.geometric,shapes.misc,shapes.arrows,decorations.pathmorphing,decorations.pathreplacing,decorations.markings,patterns,shadows,mindmap,trees,fit,backgrounds,matrix,intersections,angles,quotes}
\usepackage{pgfplots}
\usepackage[version=4]{mhchem} % équations nucléaires \ce{^{235}_{92}U}
\usepackage[european,siunitx]{circuitikz} % schémas électriques
\pgfplotsset{compat=1.18}
\usepgfplotslibrary{fillbetween,groupplots} % aires entre courbes (calcul intégral)
\usepackage{enumitem}
\usepackage{fancyhdr}
% [most] charge toutes les bibliothèques tcolorbox (listings, minted, raster,
% documentation...) et gonfle inutilement la mémoire TeX sur un document long
% (~300 boîtes) : on ne charge que ce qui est réellement utilisé.
\usepackage[skins,breakable]{tcolorbox}
\usepackage{graphicx}
\usepackage{colortbl}
\usepackage{booktabs}
\usepackage{tabularx}
\usepackage{diagbox} % case d'en-tête diagonale des tableaux à double entrée
% Colonnes extensibles centrée / gauche / droite (C, L, R), souvent utilisées par les modèles
\newcolumntype{C}{>{\centering\arraybackslash}X}
\newcolumntype{L}{>{\raggedright\arraybackslash}X}
\newcolumntype{R}{>{\raggedleft\arraybackslash}X}
\usepackage{array}
\usepackage{multicol}
\usepackage{multirow}
\usepackage{siunitx}
% parse-numbers=false : accepte aussi \SI{2,5 \times 10^{-3}}{...}
\sisetup{locale=FR,output-decimal-marker={,},group-minimum-digits=4,parse-numbers=false}
\DeclareSIUnit\ohm{\ensuremath{\symup{\Omega}}} % U+2126 absent de la police
\DeclareSIUnit\division{div} % réglages d'oscilloscope (ms/div, V/div)
\DeclareSIUnit\tr{tr} % tours (vitesse de rotation, ex. \SI{1500}{\tr\per\minute})
\DeclareSIUnit\molar{M} % concentration molaire (ex. \SI{3}{\molar}, \SI{1}{\micro\molar})
\DeclareSIUnit\an{an} % année (le modèle confond parfois avec \year, un compteur TeX) — ex. \SI{5}{\kilo\meter\per\an}
\DeclareSIUnit\FCFA{FCFA} % franc CFA (ex. \SI{500}{\FCFA\per\kilo\gram})
\DeclareSIUnit\degreeBrix{\textdegree Brix} % teneur en sucre d'un sirop/jus (réfractométrie)
\DeclareSIUnit\month{mois} % le modèle utilise parfois \month, un compteur TeX, comme unité de durée
\usepackage{qrcode}
% \degree : commande fréquemment utilisée par les modèles pour un angle ou une
% température en dehors de \SI{}{} (non fournie par les paquets chargés).
\providecommand{\degree}{^{\circ}}
% \qed (fin de démonstration), utilisé par les modèles sans amsthm
\providecommand{\qed}{\hfill$\square$}
% \barre : double barre des valeurs interdites dans un tableau de variations (tkz-tab)
\providecommand{\barre}{\ensuremath{\|}}
% environnement proof (amsthm non chargé) utilisé par les modèles
\ifdefined\proof\else\newenvironment{proof}[1][Démonstration]{\par\noindent\textit{#1.}\ }{\qed\par}\fi
\usepackage{lastpage}
\usepackage{hyperref}
\hypersetup{hidelinks}

% ============================================================
% Palette « Pop » OnBuch+ — mêmes hex que la classe P (plus_ui.dart)
% ============================================================
\definecolor{poporange}{HTML}{F59321}
\definecolor{popdark}{HTML}{E07A0C}
\definecolor{popcream}{HTML}{FFF6E8}
\definecolor{popink}{HTML}{1C1714}
\definecolor{poppurple}{HTML}{7A5AE0}
\definecolor{popgreen}{HTML}{2E9E6A}
\definecolor{poppink}{HTML}{E0457B}
\definecolor{popblue}{HTML}{2F7DE1}
\definecolor{popgold}{HTML}{C98A12}
\definecolor{popmuted}{HTML}{8A7F76}
\definecolor{popline}{HTML}{EADFD2}
\definecolor{popgray}{HTML}{8A7F76}
\colorlet{popgrey}{popgray}
% Alias tolérants (le modèle nomme parfois les couleurs différemment)
\colorlet{popwhite}{white}
\colorlet{popblack}{popink}
\colorlet{popred}{poppink}
\colorlet{popyellow}{popgold}
% Teintes claires (fonds de tableaux/encadrés) — le modèle nomme parfois
% les couleurs avec un suffixe L (ex. popblueL) sans les avoir définies.
\colorlet{poporangeL}{poporange!20}
\colorlet{popdarkL}{popdark!20}
\colorlet{poppurpleL}{poppurple!20}
\colorlet{popgreenL}{popgreen!20}
\colorlet{poppinkL}{poppink!20}
\colorlet{popblueL}{popblue!20}
\colorlet{popgoldL}{popgold!20}
\colorlet{popredL}{popred!20}
\colorlet{popgrayL}{popgray!20}
% Teintes claires pour fonds de tableaux / zones
\colorlet{popblueL}{popblue!10!white}
\colorlet{poporangeL}{poporange!14!white}
\colorlet{popgreenL}{popgreen!12!white}
\colorlet{poppurpleL}{poppurple!12!white}
\colorlet{poppinkL}{poppink!12!white}
\colorlet{popgoldL}{popgold!14!white}

\pagecolor{popcream}

% ============================================================
% Typographie
% ============================================================
\defaultfontfeatures{Ligatures=TeX}
\setmainfont{PlusJakartaSans-Medium.ttf}[Path=../../../fonts/,BoldFont=PlusJakartaSans-Bold.ttf]
% Symboles, API et emojis absents de la police principale : repli sur DejaVu Sans (caractères actifs)
\newfontface\symfont{DejaVuSans.ttf}[Path=../../../fonts/]
\makeatletter
\newcommand\symactive[1]{\catcode#1=13 \begingroup\lccode`\~=#1 \lowercase{\endgroup\def~}{{\symfont\char#1}}}
\newcommand\symblank[1]{\catcode#1=13 \begingroup\lccode`\~=#1 \lowercase{\endgroup\def~}{}}
\newcommand\symhyph[1]{\catcode#1=13 \begingroup\lccode`\~=#1 \lowercase{\endgroup\def~}{\mbox{-}}}
\newcount\symc
\newcommand\symrange[2]{\symc=#1 \@whilenum\symc<#2 \do{\expandafter\symactive\expandafter{\the\symc}\advance\symc by 1 }}
\newcommand\symblankrange[2]{\symc=#1 \@whilenum\symc<#2 \do{\expandafter\symblank\expandafter{\the\symc}\advance\symc by 1 }}
\makeatother
\symrange{"0250}{"02FF}\symactive{"2713}\symactive{"2714}\symactive{"2717}\symactive{"2718}\symactive{"2610}\symactive{"2611}\symactive{"2612}
\symhyph{"2011}\symhyph{"2E1A}\symblankrange{"1F300}{"1FAFF}\symblank{"FE0F}
% Maths en sans-serif (Fira Math) avec chiffres et lettres droites de Plus
% Jakarta, pour que les formules se fondent dans le texte.
\usepackage[math-style=ISO,bold-style=ISO]{unicode-math}
\setmathfont{FiraMath-Regular.otf}[Path=../../../fonts/]
\setmathfont{PlusJakartaSans-Medium.ttf}[Path=../../../fonts/,range={up/num,up/{Latin,latin},\mathup}]
\usepackage{icomma} % virgule décimale sans espace en mode math
% Caractères mathématiques Unicode absents des polices de texte (Plus Jakarta,
% Archivo Black) : rendus par leur équivalent mathématique, en texte comme en
% formule — sinon ils s'affichent en carré vide (ex. « Arithmétique dans ℤ »).
\usepackage{newunicodechar}
\newunicodechar{ℝ}{\ensuremath{\mathbb{R}}}
\newunicodechar{ℤ}{\ensuremath{\mathbb{Z}}}
\newunicodechar{ℂ}{\ensuremath{\mathbb{C}}}
\newunicodechar{ℕ}{\ensuremath{\mathbb{N}}}
\newunicodechar{ℚ}{\ensuremath{\mathbb{Q}}}
\newunicodechar{𝔻}{\ensuremath{\mathbb{D}}}
\newunicodechar{α}{\ensuremath{\alpha}}
\newunicodechar{β}{\ensuremath{\beta}}
\newunicodechar{γ}{\ensuremath{\gamma}}
\newunicodechar{δ}{\ensuremath{\delta}}
\newunicodechar{ε}{\ensuremath{\varepsilon}}
\newunicodechar{θ}{\ensuremath{\theta}}
\newunicodechar{λ}{\ensuremath{\lambda}}
\newunicodechar{μ}{\ensuremath{\mu}}
\newunicodechar{σ}{\ensuremath{\sigma}}
\newunicodechar{τ}{\ensuremath{\tau}}
\newunicodechar{φ}{\ensuremath{\varphi}}
\newunicodechar{ψ}{\ensuremath{\psi}}
\newunicodechar{ω}{\ensuremath{\omega}}
\newunicodechar{Σ}{\ensuremath{\Sigma}}
% majuscules produites par \MakeUppercase dans les titres (α-aminés -> Α)
\newunicodechar{Α}{\ensuremath{\alpha}}
\newunicodechar{Β}{\ensuremath{\beta}}
\newunicodechar{Γ}{\ensuremath{\Gamma}}
\newunicodechar{Δ}{\ensuremath{\Delta}}
\newunicodechar{Ε}{\ensuremath{\varepsilon}}
\newunicodechar{Θ}{\ensuremath{\Theta}}
\newunicodechar{Λ}{\ensuremath{\Lambda}}
\newunicodechar{Μ}{\ensuremath{\mu}}
\newunicodechar{Τ}{\ensuremath{\tau}}
\newunicodechar{Φ}{\ensuremath{\Phi}}
\newunicodechar{Ψ}{\ensuremath{\Psi}}
\newunicodechar{Ω}{\ensuremath{\Omega}}
\newunicodechar{↦}{\ensuremath{\mapsto}}
\newunicodechar{∈}{\ensuremath{\in}}
\newunicodechar{∉}{\ensuremath{\notin}}
\newunicodechar{∧}{\ensuremath{\wedge}}
\newunicodechar{⇔}{\ensuremath{\Leftrightarrow}}
\newunicodechar{⟺}{\ensuremath{\Longleftrightarrow}}
\newunicodechar{⇒}{\ensuremath{\Rightarrow}}
\newunicodechar{⟹}{\ensuremath{\Longrightarrow}}
\newunicodechar{∘}{\ensuremath{\circ}}
\newunicodechar{∪}{\ensuremath{\cup}}
\newunicodechar{∩}{\ensuremath{\cap}}
\newunicodechar{≡}{\ensuremath{\equiv}}
\newunicodechar{⊂}{\ensuremath{\subset}}
\newunicodechar{⊥}{\ensuremath{\perp}}
\newunicodechar{∃}{\ensuremath{\exists}}
\newunicodechar{∀}{\ensuremath{\forall}}
\newunicodechar{∛}{\ensuremath{\sqrt[3]{\,}}}
\newunicodechar{∜}{\ensuremath{\sqrt[4]{\,}}}
\newunicodechar{⃗}{\ensuremath{{}^{\to}}}
\newunicodechar{ⁿ}{\ifmmode^{n}\else\textsuperscript{\ensuremath{n}}\fi}
\newunicodechar{ˣ}{\ifmmode^{x}\else\textsuperscript{\ensuremath{x}}\fi}
\newunicodechar{ᵇ}{\ifmmode^{b}\else\textsuperscript{\ensuremath{b}}\fi}
\newunicodechar{ʳ}{\ifmmode^{r}\else\textsuperscript{\ensuremath{r}}\fi}
\newunicodechar{ᵘ}{\ifmmode^{u}\else\textsuperscript{\ensuremath{u}}\fi}
\newunicodechar{ʸ}{\ifmmode^{y}\else\textsuperscript{\ensuremath{y}}\fi}
\newunicodechar{ᵃ}{\ifmmode^{a}\else\textsuperscript{\ensuremath{a}}\fi}
\newunicodechar{ᵉ}{\ifmmode^{e}\else\textsuperscript{\ensuremath{e}}\fi}
\newunicodechar{ᵏ}{\ifmmode^{k}\else\textsuperscript{\ensuremath{k}}\fi}
\newunicodechar{ᵖ}{\ifmmode^{p}\else\textsuperscript{\ensuremath{p}}\fi}
\newunicodechar{ᴺ}{\ifmmode^{N}\else\textsuperscript{\ensuremath{N}}\fi}
\newunicodechar{ᶜ}{\ifmmode^{c}\else\textsuperscript{\ensuremath{c}}\fi}
\newunicodechar{ʰ}{\ifmmode^{h}\else\textsuperscript{\ensuremath{h}}\fi}
\newunicodechar{ᵠ}{\ifmmode^{q}\else\textsuperscript{\ensuremath{q}}\fi}
\newunicodechar{ₙ}{\ifmmode_{n}\else\textsubscript{\ensuremath{n}}\fi}
\newunicodechar{ₐ}{\ifmmode_{a}\else\textsubscript{\ensuremath{a}}\fi}
\newunicodechar{ᵢ}{\ifmmode_{i}\else\textsubscript{\ensuremath{i}}\fi}
\newunicodechar{ᵣ}{\ifmmode_{r}\else\textsubscript{\ensuremath{r}}\fi}
\newunicodechar{ₖ}{\ifmmode_{k}\else\textsubscript{\ensuremath{k}}\fi}
\newunicodechar{ᵦ}{\ifmmode_{\beta}\else\textsubscript{\ensuremath{\beta}}\fi}
\newunicodechar{ₓ}{\ifmmode_{x}\else\textsubscript{\ensuremath{x}}\fi}
\newfontfamily\popdisplay{ArchivoBlack-Regular.ttf}[Path=../../../fonts/]
\newfontfamily\poplogo{SpaceGrotesk-Bold.ttf}[Path=../../../fonts/]
\newfontfamily\popsemi{PlusJakartaSans-SemiBold.ttf}[Path=../../../fonts/]
\newfontfamily\popxbold{PlusJakartaSans-ExtraBold.ttf}[Path=../../../fonts/]
\newfontfamily\popxboldls{PlusJakartaSans-ExtraBold.ttf}[Path=../../../fonts/,LetterSpace=3.0]

\setlist[enumerate]{leftmargin=1.7em,itemsep=5pt,topsep=4pt}
\setlist[itemize]{leftmargin=1.7em,itemsep=4pt,topsep=4pt,label={\color{poporange}\textbullet}}
\setlength{\parindent}{0pt}
\setlength{\parskip}{5pt}
\renewcommand{\arraystretch}{1.25}

% pgfplots : style Pop par défaut
\pgfplotsset{
  every axis/.append style={axis line style={popink,line width=1pt},tick style={popink},
    grid style={popline,line width=0.5pt},label style={font=\small},tick label style={font=\footnotesize}},
  popcurve/.style={poporange,line width=1.8pt,no markers,smooth},
}
% Même style disponible dans un \draw TikZ simple (hors pgfplots)
\tikzset{popcurve/.style={poporange,line width=1.8pt}}
\tikzset{popgrid/.style={popline,very thin}}
\colorlet{popcurve}{poporange} % le nom du style est parfois utilisé comme couleur

% ============================================================
% Pill / squiggle
% ============================================================
\newcommand{\poppill}[3]{%
  \tikz[baseline=-0.55ex]\node[draw=popink,line width=1.2pt,rounded corners=2.6pt,%
    fill=#2,inner xsep=6.5pt,inner ysep=3pt]{%
    \popxboldls\fontsize{7}{7}\selectfont\color{#3}\MakeUppercase{#1}};%
}
\newcommand{\popsquiggle}[1]{%
  \tikz[baseline=-0.4ex]{
    \draw[#1,line width=2.6pt,line cap=round]
      plot [smooth] coordinates {(0,0.09) (0.34,0.01) (0.68,0.21) (1.02,0.01) (1.36,0.10)};
  }%
}
% Mot-clé surligné (marqueur orange sous le texte)
\newcommand{\cle}[1]{{\popxbold\color{popdark}#1}}

% ============================================================
% En-tête / pied
% ============================================================
\pagestyle{fancy}
\fancyhf{}
\renewcommand{\headrulewidth}{0pt}
\renewcommand{\footrulewidth}{0pt}
\lhead{\raisebox{-0.15ex}{\poplogo\fontsize{13}{13}\selectfont\color{popink}OnBuch\color{poporange}+}}
\rhead{\poppill{\DOCMATIERE\ \textbullet\ \DOCNIVEAU\ \textbullet\ Leçon \DOCLECON}{white}{popink}}
\lfoot{\popxbold\fontsize{7.6}{7.6}\selectfont\color{popmuted}\MakeUppercase{Cours signé OnBuch+}\ \textbullet\ {\popsemi\DOCTITRE}}
\rfoot{\tikz[baseline=-0.6ex]\node[rounded corners=8.5pt,fill=popink,minimum height=17pt,minimum width=17pt,inner xsep=5pt,inner sep=0]{\popxbold\fontsize{8}{8}\selectfont\color{popcream}\thepage\,/\,\pageref*{LastPage}};}

% ============================================================
% Titres
% ============================================================
\newcommand{\chaptitre}[2]{%
  \begin{tcolorbox}[enhanced,colback=poporange,colframe=popink,boxrule=2.6pt,arc=11pt,
      drop shadow=popink,left=15pt,right=15pt,top=12pt,bottom=11pt,before skip=3pt,after skip=18pt]
    \poppill{#2}{popink}{popcream}\par\vspace{9pt}
    {\raggedright\hyphenpenalty=10000\exhyphenpenalty=10000\popdisplay\fontsize{22}{24}\selectfont\color{popcream}\MakeUppercase{#1}\par}
  \end{tcolorbox}
}
% Section numérotée : pastille orange avec le numéro + titre Archivo
\newcounter{popsec}
\newcommand{\coursec}[1]{%
  \stepcounter{popsec}\setcounter{popsub}{0}%
  \par\vspace{16pt}\noindent
  \begin{minipage}{\linewidth}
  \tikz[baseline=-0.7ex]\node[draw=popink,line width=1.6pt,fill=poporange,rounded corners=4pt,minimum size=22pt,inner sep=0]{\popdisplay\fontsize{12}{12}\selectfont\color{popcream}\thepopsec};%
  \hspace{8pt}{\popdisplay\fontsize{15}{17}\selectfont\color{popink}\MakeUppercase{#1}}\par
  \vspace{4pt}\noindent{\color{poporange}\rule{36pt}{3.6pt}}
  \end{minipage}\par\nopagebreak\vspace{8pt}
}
\newcounter{popsub}
\newcommand{\courssub}[1]{%
  \stepcounter{popsub}%
  \par\vspace{9pt}\noindent{\popxbold\fontsize{11.5}{13}\selectfont\color{popink}{\color{poporange}\thepopsec.\thepopsub}\hspace{6pt}#1}\par\nopagebreak\vspace{4pt}
}

% ============================================================
% Boîtes pédagogiques — même recette : fond blanc, bordure épaisse colorée,
% ombre dure encre, badge plein. Titre optionnel [#1].
% ============================================================
\tcbset{popbase/.style={breakable,enhanced,colback=white,boxrule=2.3pt,arc=9pt,
  shadow={3pt}{-3pt}{0pt}{popink},left=14pt,right=14pt,top=11pt,bottom=11pt,before skip=11pt,after skip=15pt,
  fontupper=\popsemi\fontsize{10.6}{14.5}\selectfont\color{popink},
  fontlower=\popsemi\fontsize{10.6}{14.5}\selectfont\color{popink}}}
% Coche dessinée (le glyphe ✓ n'existe pas dans les polices du document)
\newcommand{\popcheck}[1]{\tikz[baseline=-0.5ex]\draw[#1,line width=1.6pt,line cap=round,line join=round](-0.13,0.0)--(-0.04,-0.09)--(0.13,0.1);}
\providecommand{\coche}{\popcheck{popgreen}}
\providecommand{\resultat}[1]{\par\smallskip\noindent\fcolorbox{popgreen}{popgreenL}{\parbox{\dimexpr\linewidth-2\fboxsep-2\fboxrule\relax}{\textbf{Résultat.} #1}}\par\smallskip}
\newcommand{\popboxhead}[3]{\poppill{#1}{#2}{white}\ifx\relax#3\relax\else\hspace{6pt}{\popxbold\fontsize{10.6}{12}\selectfont\color{popink}#3}\fi\par\vspace{7pt}}

\newtcolorbox{aretenir}[1][]{popbase,colframe=popgreen,before upper={\popboxhead{À retenir}{popgreen}{#1}}}
% Texte en espagnol (document de lecture, dialogue, poème, extrait) : cadre
% d'étude. Un titre optionnel (source, auteur) s'affiche dans l'en-tête.
\newtcolorbox{textebox}[1][]{popbase,colframe=popblue,before upper={\popboxhead{Texto}{popblue}{#1}\begin{otherlanguage}{spanish}},after upper={\end{otherlanguage}}}
% Marques ✓ / ✗ (forme correcte / fautive) souvent utilisées par les modèles
\providecommand{\cross}{\textcolor{poppink}{\ensuremath{\times}}}
\providecommand{\xmark}{\cross}\providecommand{\crossmark}{\cross}\providecommand{\cmark}{\ensuremath{\checkmark}}
% Ligne de réponse / justification à compléter (inventée par les modèles)
\providecommand{\justification}[1]{\par\noindent\textit{Justification~:} \dotfill\par}
% \textipa (package tipa non chargé) : la police ne couvre pas l'API -> simple passage
\providecommand{\textipa}[1]{#1}
% Soulignement (ulem non chargé) : \uline, \ul
\providecommand{\uline}[1]{\underline{#1}}\providecommand{\ul}[1]{\underline{#1}}
% \glossaire{\item ...} inventée par les modèles : liste de glossaire
\providecommand{\glossaire}[1]{\begin{itemize}#1\end{itemize}}
% \alert (beamer) inventée par les modèles : texte mis en évidence
\providecommand{\alert}[1]{\textbf{\textcolor{poppink}{#1}}}
% variantes ASCII d'ordinaux écrites en macro
\providecommand{\iere}{\textsuperscript{re}}\providecommand{\ieme}{\textsuperscript{e}}\providecommand{\ere}{\textsuperscript{re}}\providecommand{\eme}{\textsuperscript{e}}\providecommand{\er}{\textsuperscript{er}}\providecommand{\ie}{\textsuperscript{e}}
% Ordinaux français écrits en macro par les modèles : 1\ière, 3\ième
\newcommand{\ière}{\textsuperscript{re}}\newcommand{\ième}{\textsuperscript{e}}
% \exemple (inventée par les modèles) : étiquette « Exemple : »
\providecommand{\exemple}{\textbf{Exemple~:}\ }
% Mot ou phrase en espagnol à l'intérieur d'un paragraphe français
\newcommand{\esp}[1]{\ifmmode\text{\textit{#1}}\else\begingroup\selectlanguage{spanish}\itshape #1\endgroup\fi}
\newtcolorbox{exemplebox}[1][]{popbase,colframe=poporange,before upper={\popboxhead{Exemple}{poporange}{#1}}}
\newtcolorbox{definition}[1][]{popbase,colframe=popblue,before upper={\popboxhead{Définition}{popblue}{#1}}}
\newtcolorbox{propriete}[1][]{popbase,colframe=poppurple,before upper={\popboxhead{Propriété}{poppurple}{#1}}}
\newtcolorbox{methode}[1][]{popbase,colframe=popgold,before upper={\popboxhead{Méthode}{popgold}{#1}}}
\newtcolorbox{attention}[1][]{popbase,colframe=poppink,before upper={\popboxhead{Attention}{poppink}{#1}}}
\newtcolorbox{experience}[1][]{popbase,colframe=popblue,colback=popblueL,before upper={\popboxhead{Expérience}{popblue}{#1}}}
% « Le savais-tu ? » : carte violette pleine (contexte, culture, Cameroun)
\newtcolorbox{savaistu}[1][]{popbase,colback=poppurple,colframe=popink,
  fontupper=\popsemi\fontsize{10.4}{14.2}\selectfont\color{white},
  before upper={\poppill{Le savais-tu\,?}{white}{poppurple}\ifx\relax#1\relax\else\hspace{6pt}{\popxbold\color{white}#1}\fi\par\vspace{7pt}}}
% Exercice résolu : énoncé en haut, \tcblower puis la solution sur fond crème
\newcounter{popexo}\newcounter{popexr}
\newtcolorbox{exoresolu}[1][]{popbase,colframe=poporange,colbacklower=popcream,
  segmentation style={popink,line width=1.2pt,dashed},
  before upper={\stepcounter{popexr}\popboxhead{Exercice résolu \thepopexr}{poporange}{#1}},
  before lower={\poppill{Solution}{popink}{popcream}\par\vspace{6pt}}}
% Exercice (non résolu en place) — la correction est dans la partie Corrigés
\newtcolorbox{exercice}[1][]{popbase,colframe=popink,
  before upper={\stepcounter{popexo}\popboxhead{Exercice \thepopexo}{popink}{#1}}}
% Corrigé
\newtcolorbox{corrige}[1][]{popbase,colframe=popgreen,colback=popgreenL,
  before upper={\popboxhead{Corrigé}{popgreen}{#1}}}
% Objectifs / prérequis / bilan
\newtcolorbox{objectifs}{popbase,colframe=popink,colback=poporangeL,
  before upper={\popboxhead{Objectifs de la leçon}{popink}{}}}
\newtcolorbox{prerequis}{popbase,colframe=popink,colback=popblueL,
  before upper={\popboxhead{Prérequis}{popblue}{}}}
\newtcolorbox{bilan}{popbase,colframe=popink,colback=white,boxrule=2.8pt,
  before upper={\popboxhead{Fiche bilan}{poporange}{L'essentiel à connaître pour l'examen}}}
% Figure encadrée avec légende
\newtcolorbox{popfigure}[1][]{enhanced,breakable=false,colback=white,colframe=popline,boxrule=1.4pt,arc=8pt,
  left=10pt,right=10pt,top=10pt,bottom=8pt,before skip=10pt,after skip=12pt,halign=center,
  lower separated=false,halign lower=center,
  fontlower=\popsemi\fontsize{9}{11.5}\selectfont\color{popmuted},
  #1}
\newcommand{\legende}[1]{\par\vspace{6pt}{\popsemi\fontsize{9}{11.5}\selectfont\color{popmuted}#1}}

% ============================================================
% Signature OnBuch+
% ============================================================
\newcommand{\poplogoinline}[1]{{\poplogo\fontsize{#1}{#1}\selectfont\color{popink}OnBuch\color{poporange}+}}
\newcommand{\signatureonbuch}{%
  \begin{tcolorbox}[enhanced,colback=popink,colframe=popink,boxrule=2.2pt,arc=10pt,
      drop shadow=poporange,left=14pt,right=14pt,top=10pt,bottom=10pt,before skip=0pt,after skip=0pt]
    \tikz[baseline=-0.7ex]\node[circle,fill=popgreen,draw=popcream,line width=1.3pt,minimum size=22pt,inner sep=0]{\popcheck{white}};%
    \hspace{9pt}\begin{minipage}[c]{0.62\linewidth}
      {\popxbold\fontsize{12}{13}\selectfont\color{popcream}Signé\ \textbullet\ {\poplogo OnBuch\color{poporange}+}}\par\vspace{2pt}
      {\raggedright\popsemi\fontsize{8}{10.5}\selectfont\color{popline}\MakeUppercase{Équipe pédagogique OnBuch+}\\\MakeUppercase{Conforme au programme officiel MINESEC}\par}
    \end{minipage}\hfill
    {\poplogo\fontsize{22}{22}\selectfont\color{popcream}OnBuch\color{poporange}+}
  \end{tcolorbox}%
}

% ============================================================
% Page de garde
% #1 titre · #2 matière · #3 niveau · #4 module · #5 n° leçon · #6 durée ·
% #7 description / objectifs (texte ou itemize)
% ============================================================
\newcommand{\pagegarde}[7]{%
  \thispagestyle{empty}
  \newgeometry{a4paper,margin=1.7cm,top=1.6cm,bottom=1.4cm}%
  \begin{tikzpicture}[remember picture,overlay]
    \fill[poporange!18] ([xshift=-3.2cm,yshift=-3.6cm]current page.north east) circle (4.6cm);
    \fill[poppurple!14] ([xshift=2.4cm,yshift=7.6cm]current page.south west) circle (3.8cm);
  \end{tikzpicture}%
  {\poplogo\fontsize{30}{30}\selectfont\color{popink}OnBuch\color{poporange}+}\hfill
  \raisebox{0.6ex}{\poppill{Cours signé \textbullet\ Premium}{popink}{popcream}}\par
  \vspace{1pt}\popsquiggle{poppurple}\par
  \vspace{8pt}
  \begin{tcolorbox}[enhanced,colback=poporange,colframe=popink,boxrule=2.8pt,arc=13pt,
      drop shadow=popink,left=18pt,right=18pt,top=12pt,bottom=13pt,after skip=10pt]
    \poppill{#2 \textbullet\ #3}{popink}{popcream}\hspace{5pt}\poppill{#4}{white}{popink}\par\vspace{8pt}
    {\popdisplay\fontsize{14}{14}\selectfont\color{popink}LEÇON #5}\par\vspace{4pt}
    {\raggedright\hyphenpenalty=10000\exhyphenpenalty=10000\popdisplay\fontsize{25}{27}\selectfont\color{popcream}\MakeUppercase{#1}\par}
    \vspace{8pt}
    {\popxbold\fontsize{9.5}{9.5}\selectfont\color{popink}\MakeUppercase{Durée conseillée : #6}}
  \end{tcolorbox}
  \begin{tcolorbox}[enhanced,colback=white,colframe=popink,boxrule=2.2pt,arc=10pt,
      drop shadow=popink,left=16pt,right=16pt,top=10pt,bottom=10pt,after skip=8pt]
    \poppill{Dans cette leçon}{poppurple}{white}\par\vspace{5pt}
    \popsemi\fontsize{9.3}{12.6}\selectfont\color{popink}#7
  \end{tcolorbox}
  \noindent\begin{minipage}[t]{0.487\linewidth}
    \begin{tcolorbox}[enhanced,colback=popgreen,colframe=popink,boxrule=2.2pt,arc=10pt,
        drop shadow=popink,left=11pt,right=11pt,top=10pt,bottom=10pt,height=3.1cm,valign=center]
      \tcbox[colback=white,colframe=popink,boxrule=1.3pt,arc=5pt,boxsep=0pt,nobeforeafter,
        left=3pt,right=3pt,top=3pt,bottom=3pt,valign=center]{\qrcode[height=1.9cm]{https://play.google.com/store/apps/details?id=com.luvvix.onbuch&hl=fr}}%
      \hspace{8pt}\begin{minipage}[c]{0.5\linewidth}
        {\popdisplay\fontsize{11}{12.5}\selectfont\color{white}\MakeUppercase{Télécharge OnBuch}}\par\vspace{4pt}
        {\raggedright\popsemi\fontsize{8.4}{11}\selectfont\color{white}Gratuit \textbullet\ Google Play\\Tuteur IA, annales, concours, résultats\par}
      \end{minipage}
    \end{tcolorbox}
  \end{minipage}\hfill
  \begin{minipage}[t]{0.487\linewidth}
    \begin{tcolorbox}[enhanced,colback=poppurple,colframe=popink,boxrule=2.2pt,arc=10pt,
        drop shadow=popink,left=11pt,right=11pt,top=10pt,bottom=10pt,height=3.1cm,valign=center]
      \tikz[baseline=-0.7ex]\node[circle,fill=white,draw=popink,line width=1.3pt,minimum size=1.6cm,inner sep=0]{\popdisplay\fontsize{24}{24}\selectfont\color{poppurple}+};%
      \hspace{8pt}\begin{minipage}[c]{0.6\linewidth}
        {\popdisplay\fontsize{11}{12.5}\selectfont\color{white}\MakeUppercase{Passe à OnBuch+}}\par\vspace{4pt}
        {\raggedright\popsemi\fontsize{8.4}{11}\selectfont\color{white}Tous les cours signés, Tuteur IA illimité, fiches PDF — onglet OnBuch+ de l'app\par}
      \end{minipage}
    \end{tcolorbox}
  \end{minipage}
  \vspace{8pt}
  \popcontenu
  \vfill
  \signatureonbuch
  \restoregeometry
  \newpage
}

% Rangée « Ce cours contient » (page de garde)
\newcommand{\poptile}[3]{% #1 couleur #2 gros texte #3 légende
  \begin{tcolorbox}[enhanced,colback=#1,colframe=popink,boxrule=2pt,arc=9pt,shadow={2.5pt}{-2.5pt}{0pt}{popink},
      left=6pt,right=6pt,top=9pt,bottom=9pt,halign=center,valign=center,height=1.75cm,width=\linewidth,nobeforeafter]
    {\popdisplay\fontsize{12.5}{13}\selectfont\color{white}\MakeUppercase{#2}}\par\vspace{3pt}
    {\centering\popsemi\fontsize{7.8}{9.5}\selectfont\color{white}#3\par}
  \end{tcolorbox}}
\newcommand{\popcontenu}{%
  \par\noindent{\popxboldls\fontsize{7.5}{7.5}\selectfont\color{popmuted}CE COURS SIGNÉ CONTIENT}\par\vspace{7pt}
  \noindent\begin{minipage}[t]{0.235\linewidth}\poptile{popblue}{Cours}{complet, illustré, pas à pas}\end{minipage}\hfill
  \begin{minipage}[t]{0.235\linewidth}\poptile{poporange}{Méthodes}{et exercices résolus}\end{minipage}\hfill
  \begin{minipage}[t]{0.235\linewidth}\poptile{poppink}{Exercices}{type examen + corrigés}\end{minipage}\hfill
  \begin{minipage}[t]{0.235\linewidth}\poptile{popgreen}{Fiche bilan}{l'essentiel à retenir}\end{minipage}\par}

% Sommaire
\newtcolorbox{sommaire}{enhanced,colback=white,colframe=popink,boxrule=2.2pt,arc=10pt,
  drop shadow=popink,left=18pt,right=18pt,top=13pt,bottom=13pt,before skip=6pt,after skip=20pt,
  before upper={\poppill{Sommaire}{poppurple}{white}\par\vspace{9pt}\popsemi\fontsize{10.6}{14.5}\selectfont\color{popink}}}

% Signature de fin de cours — poussée tout en bas de la dernière page
\newcommand{\findecours}{%
  \par\vfill\nopagebreak
  \begin{center}\popsquiggle{poporange}\end{center}\vspace{4pt}
  \signatureonbuch
}
\AtBeginDocument{\def\llbracket{\mathopen{[\mkern-3.5mu[}}\def\rrbracket{\mathclose{]\mkern-3.5mu]}}} % intervalles d'entiers (glyphe ⟦ absent de la police)
% Glyphes absents de Fira Math : redéfinis à partir de glyphes disponibles
\AtBeginDocument{%
  \def\setminus{\mathbin{\backslash}}\let\smallsetminus\setminus
  \def\vdots{\mathord{\vbox{\baselineskip3.2pt\lineskiplimit0pt\kern4pt\hbox{.}\hbox{.}\hbox{.}}}}%
  \def\ddots{\mathinner{\mkern1mu\raise6.5pt\hbox{.}\mkern2mu\raise3.6pt\hbox{.}\mkern2mu\raise0.7pt\hbox{.}\mkern1mu}}%
  \def\triangle{\mathord{\scalebox{0.9}{$\symup{\Delta}$}}}%
  \def\nabla{\mathord{\raisebox{\depth}{\scalebox{1}[-1]{$\symup{\Delta}$}}}}%
  \def\checkmark{\popcheck{popdark}}%
  \def\star{\mathbin{\ast}}\def\bigstar{\mathord{\ast}}%
  \def\diamond{\mathbin{\scalebox{0.6}{$\Diamond$}}}%
  \def\bigcup{\mathop{\mathchoice{\raisebox{-0.3em}{\scalebox{1.7}{$\cup$}}}{\scalebox{1.25}{$\cup$}}{\cup}{\cup}}\displaylimits}%
  \def\bigcap{\mathop{\mathchoice{\raisebox{-0.3em}{\scalebox{1.7}{$\cap$}}}{\scalebox{1.25}{$\cap$}}{\cap}{\cap}}\displaylimits}%
  \def\lesseqqgtr{\mathrel{\lessgtr}}\def\bigoplus{\mathop{\oplus}\displaylimits}%
}
\providecommand{\ssi}{si et seulement si} % abréviation fréquente
\AtBeginDocument{\providecommand{\wideparen}{\overparen}} % arc de cercle

%% FIGFIX-BEGIN v4 — correctifs globaux des figures (docs/onbuch-generation/tools/figfix)
\usepackage{adjustbox}\usepackage{etoolbox}
% toute figure TikZ/pgfplots plus large que la ligne est réduite à la largeur disponible
\BeforeBeginEnvironment{tikzpicture}{\begin{adjustbox}{max width=\linewidth}}
\AfterEndEnvironment{tikzpicture}{\end{adjustbox}}
% légendes pgfplots lisibles par-dessus les courbes
\pgfplotsset{every axis/.append style={legend style={fill=white,fill opacity=0.92,text opacity=1,draw=black!15,inner sep=3pt}}}
% étiquettes d'arêtes (syntaxe edge["..."]) avec fond blanc
\tikzset{every edge quotes/.append style={fill=white,inner sep=1.2pt}}
%% FIGFIX-END
\begin{document}

\pagegarde{Las conductas de riesgo : sensibilizar para adoptar comportamientos responsables}{Espagnol}{1ère A}{Módulo 1 : Vie familiale et intégration sociale}{EA04}{2 h}{Cette leçon mixte du Módulo 1 « Vie familiale et intégration sociale » part d'un texte de sensibilisation sur les comportements responsables à l'école et en société. Elle permet d'acquérir le lexique des conduites à risque, de maîtriser l'impératif affirmatif et négatif ainsi que les structures du conseil et de l'obligation, pour aboutir à la création d'un message de prévention.
\vspace{4pt}
\begin{multicols}{2}\small
\begin{itemize}
\item À la fin de cette leçon, je sais comprendre un texte de sensibilisation sur les comportements responsables
\item À la fin de cette leçon, je sais utiliser le lexique des conduites à risque et des comportements responsables
\item À la fin de cette leçon, je sais former et employer l'impératif affirmatif (tú, vosotros, usted) pour donner des conseils et des consignes
\item À la fin de cette leçon, je sais former l'impératif négatif avec le subjonctif pour interdire
\item À la fin de cette leçon, je sais exprimer le conseil et l'obligation avec deber + infinitif, hay que + infinitif et es importante que + subjonctif
\item À la fin de cette leçon, je sais placer correctement les pronoms avec l'impératif affirmatif et négatif
\end{itemize}\end{multicols}}

\begin{sommaire}
\begin{multicols}{2}
\begin{enumerate}
\item Texte support : « Jóvenes responsables, sociedad fuerte »
\item Lexique thématique : les conduites à risque et les comportements responsables
\item Grammaire 1 : l'impératif affirmatif (tú, vosotros, usted)
\item Grammaire 2 : l'impératif négatif et les structures du conseil et de l'obligation
\item Méthode du commentaire de texte et commentaire modèle
\item Production guidée : créer un message de prévention
\item Activité d'intégration
\item Méthodes et exercices résolus
\item Exercices
\item Corrigés des exercices
\item Fiche bilan
\end{enumerate}
\end{multicols}
\end{sommaire}

\begin{objectifs}
\begin{itemize}
\item À la fin de cette leçon, je sais comprendre un texte de sensibilisation sur les comportements responsables
\item À la fin de cette leçon, je sais utiliser le lexique des conduites à risque et des comportements responsables
\item À la fin de cette leçon, je sais former et employer l'impératif affirmatif (tú, vosotros, usted) pour donner des conseils et des consignes
\item À la fin de cette leçon, je sais former l'impératif négatif avec le subjonctif pour interdire
\item À la fin de cette leçon, je sais exprimer le conseil et l'obligation avec deber + infinitif, hay que + infinitif et es importante que + subjonctif
\item À la fin de cette leçon, je sais placer correctement les pronoms avec l'impératif affirmatif et négatif
\item À la fin de cette leçon, je sais commenter un texte de sensibilisation selon la méthode du commentaire
\item À la fin de cette leçon, je sais rédiger un message de prévention (affiche, slogan, dialogue)
\item À la fin de cette leçon, je sais jouer un court dialogue de sensibilisation
\end{itemize}
\end{objectifs}

\begin{prerequis}
\begin{itemize}
\item Présent de l'indicatif des verbes réguliers et irréguliers (Seconde)
\item Présent du subjonctif : formation de base (début de Première)
\item Les pronoms compléments d'objet direct et indirect (lo, la, le, se...)
\item Le vocabulaire de la vie scolaire et de la famille (Módulo 1)
\item Les verbes modaux simples : querer, poder, tener que
\end{itemize}
\end{prerequis}

\newpage

\coursec{Texte support : « Jóvenes responsables, sociedad fuerte »}

\courssub{Situation d'accroche et problématique}

C'est la rentrée au lycée de Yaoundé. Sur les réseaux sociaux, une vidéo circule : on y voit des élèves qui fument derrière les salles de classe, d'autres qui trichent pendant les évaluations, d'autres encore qui se moquent d'un camarade plus faible. Le proviseur, très inquiet, décide de lancer une grande campagne de sensibilisation. Et comme c'est un lycée bilingue, il demande aux élèves de Première de créer des affiches et des slogans... en espagnol !

Mais voilà : comment dire « Ne fume pas ! », « Il faut respecter les autres », « Tu dois dire non à la drogue » ? Comment rédiger un message de prévention clair et percutant dans la langue de Cervantès ? C'est exactement ce que cette leçon va t'apprendre.

\textbf{Problématique :} comment un texte de sensibilisation en espagnol dénonce-t-il les conduites à risque et conseille-t-il des comportements responsables ? Pour y répondre, nous allons d'abord lire un article de journal scolaire mexicain, enrichir notre vocabulaire, puis repérer dans le texte les outils grammaticaux du conseil et de l'interdiction : l'impératif, \esp{deber} + infinitif, \esp{hay que} et \esp{es importante que} + subjonctif.

\courssub{Lecture globale du texte}

Avant de lire, observe le titre : \esp{Jóvenes responsables, sociedad fuerte}. Que promet-il ? Une idée simple et forte : la responsabilité des jeunes construit la force de la société. Lis maintenant le texte une première fois, sans t'arrêter sur les mots inconnus, pour saisir le sens général.

\begin{textebox}[« Jóvenes responsables, sociedad fuerte » — article de journal scolaire, Mexique]
En nuestra escuela, como en muchas del país, algunos jóvenes adoptan conductas de riesgo: fuman, beben alcohol, hacen trampa en los exámenes o maltratan a sus compañeros. Estas conductas destruyen la salud y la convivencia.

El tabaco y el alcohol parecen una diversión, pero en realidad son una trampa: crean dependencia, dañan el cuerpo y arruinan el futuro. La droga es todavía peor: destruye familias enteras. La trampa en los exámenes es otra forma de engaño: el que hace trampa se engaña a sí mismo y no aprende nada. Finalmente, la violencia y el acoso entre compañeros rompen la amistad y la paz de la escuela. Las redes sociales, mal usadas, también pueden convertirse en un instrumento de burla y de mentiras.

Frente a estos problemas, todos podemos actuar. ¡Di no a las drogas! No fumes y no aceptes la violencia. Debes respetar a los demás y hay que cuidar la salud. Es importante que hables con tus padres y tus profesores si tienes problemas. También es bueno practicar deporte, leer y participar en las actividades de la escuela.

Un joven responsable construye una sociedad fuerte. La escuela no es solamente un lugar para estudiar: es un lugar para aprender a vivir juntos. Sensibilizar a los jóvenes es preparar el futuro de todo el país. ¡Actúa hoy, tu futuro te lo agradecerá!
\end{textebox}

\begin{definition}[Glossaire du texte]
\begin{center}
\begin{tabularx}{\linewidth}{|X|X|}
\hline
\rowcolor{popblueL} \textbf{Español} & \textbf{Français} \\
\hline
las conductas de riesgo & les conduites à risque \\
\hline
sensibilizar & sensibiliser \\
\hline
el tabaco / fumar & le tabac / fumer \\
\hline
el alcohol / beber alcohol & l'alcool / boire de l'alcool \\
\hline
la droga & la drogue \\
\hline
la trampa / hacer trampa & la triche / tricher \\
\hline
maltratar & maltraiter, brutaliser \\
\hline
el acoso & le harcèlement \\
\hline
la convivencia & la vie en communauté, le vivre-ensemble \\
\hline
el respeto / respetar & le respect / respecter \\
\hline
cuidar & prendre soin de \\
\hline
evitar & éviter \\
\hline
la salud & la santé \\
\hline
los demás & les autres \\
\hline
la burla & la moquerie \\
\hline
engañar & tromper, duper \\
\hline
construir & construire \\
\hline
agradecer & remercier, être reconnaissant de \\
\hline
\end{tabularx}
\end{center}
\end{definition}

\begin{attention}[Faux amis et pièges de vocabulaire]
\begin{itemize}
\item \esp{las drogas} = « les drogues » : attention, en espagnol on dit souvent \esp{la droga} au singulier pour parler de la drogue en général ; \esp{las drogas} désigne les différentes substances.
\item \esp{hacer trampa} = « tricher » (à un examen, à un jeu). Ne dis jamais \esp{hacer triche} ! Le nom est \esp{la trampa}, littéralement « le piège ».
\item \esp{la salud} = « la santé » : ne pas confondre avec l'interjection \esp{¡Salud!} = « Santé ! » (trinquons) ou « À tes souhaits ! » (éternuement).
\item \esp{los demás} = « les autres » : \esp{respetar a los demás} = « respecter les autres ». Note la préposition \esp{a} devant les personnes (a personnel).
\item \esp{las redes sociales} = « les réseaux sociaux » : \esp{red} signifie « filet » puis « réseau ».
\item \esp{construir} = « construire » (irrégulier au présent : \esp{construyo}, \esp{construyes}...). Ne pas confondre avec \esp{destruir} (« détruire »).
\end{itemize}
\end{attention}

\courssub{Compréhension globale du texte}

\begin{exoresolu}[Comprendre le sens général]
Énoncé. Réponds aux questions suivantes en français, puis en espagnol : 1) ¿De qué habla el texto? 2) ¿Quién escribe y a quién se dirige? 3) ¿Cuál es el objetivo del texto?
\tcblower
Solution détaillée.
\begin{enumerate}
\item Le texte parle des conduites à risque des adolescents (tabac, alcool, drogue, triche, violence, mauvais usage des réseaux sociaux) et des comportements responsables à adopter. \esp{El texto habla de las conductas de riesgo de los jóvenes y de los comportamientos responsables.}
\item C'est un élève journaliste (ou la rédaction) d'un journal scolaire mexicain qui écrit ; il s'adresse aux élèves de son école, donc à des adolescents. On le voit grâce au tutoiement (\esp{di}, \esp{fumes}, \esp{hables}) et au mot \esp{nuestra escuela}. \esp{Escribe un periodista del periódico escolar y se dirige a los alumnos de la escuela.}
\item L'objectif est de sensibiliser les jeunes et de les convaincre d'adopter des comportements responsables. \esp{El objetivo es sensibilizar a los jóvenes para que adopten comportamientos responsables.}
\end{enumerate}
\end{exoresolu}

\courssub{Compréhension détaillée du texte}

\begin{exercice}[À toi de chercher dans le texte]
\begin{enumerate}
\item Cite trois conduites à risque mentionnées dans le texte.
\item Cite trois conseils donnés aux jeunes.
\item Quelles sont les conséquences des conduites à risque évoquées par l'auteur ?
\item Relève dans le texte deux phrases à l'impératif et deux structures de conseil.
\end{enumerate}
\end{exercice}

\begin{corrige}[Exercice de compréhension détaillée]
\begin{enumerate}
\item Trois conduites à risque : \esp{fumar} (fumer), \esp{beber alcohol} (boire de l'alcool), \esp{hacer trampa en los exámenes} (tricher aux examens). On peut ajouter \esp{la droga}, \esp{la violencia}, \esp{el acoso} et le mauvais usage des \esp{redes sociales}.
\item Trois conseils : \esp{¡Di no a las drogas!} (Dis non à la drogue !), \esp{No fumes} (Ne fume pas), \esp{Debes respetar a los demás} (Tu dois respecter les autres), \esp{habla con tus padres} (parle avec tes parents).
\item Les conséquences : le tabac et l'alcool \esp{crean dependencia, dañan el cuerpo y arruinan el futuro} (créent une dépendance, abîment le corps et ruinent l'avenir) ; la drogue \esp{destruye familias enteras} ; la triche empêche d'apprendre (\esp{no aprende nada}) ; la violence \esp{rompe la amistad y la paz de la escuela}.
\item Impératifs : \esp{¡Di no a las drogas!}, \esp{No fumes}, \esp{¡Actúa hoy!}. Structures de conseil : \esp{Debes respetar a los demás} (\esp{deber} + infinitif), \esp{hay que cuidar la salud} (il faut), \esp{Es importante que hables con tus padres} (il est important que + subjonctif).
\end{enumerate}
\end{corrige}

\courssub{Repérage des impératifs et des structures de conseil}

Le texte de sensibilisation utilise des outils grammaticaux précis pour conseiller et interdire. Observons-les avant de les étudier en détail dans les sections de grammaire qui suivront.

\begin{exemplebox}[Les formes du conseil dans le texte]
\begin{itemize}
\item \esp{¡Di no a las drogas!} « Dis non à la drogue ! » — impératif affirmatif de \esp{decir} (tú).
\item \esp{No fumes.} « Ne fume pas. » — impératif négatif de \esp{fumar} (tú).
\item \esp{No aceptes la violencia.} « N'accepte pas la violence. » — impératif négatif de \esp{aceptar}.
\item \esp{Debes respetar a los demás.} « Tu dois respecter les autres. » — \esp{deber} + infinitif.
\item \esp{Hay que cuidar la salud.} « Il faut prendre soin de sa santé. » — tournure impersonnelle \esp{hay que} + infinitif.
\item \esp{Es importante que hables con tus padres.} « Il est important que tu parles avec tes parents. » — \esp{es importante que} + subjonctif.
\item \esp{¡Actúa hoy!} « Agis aujourd'hui ! » — impératif affirmatif de \esp{actuar}.
\end{itemize}
\end{exemplebox}

\begin{aretenir}[Les quatre outils du conseil]
\begin{center}
\begin{tabularx}{\linewidth}{|X|X|X|}
\hline
\rowcolor{popblueL} \textbf{Structure} & \textbf{Exemple} & \textbf{Français} \\
\hline
Impératif affirmatif & \esp{¡Di no!} & Dis non ! \\
\hline
Impératif négatif & \esp{No fumes.} & Ne fume pas. \\
\hline
\esp{deber} + infinitif & \esp{Debes respetar.} & Tu dois respecter. \\
\hline
\esp{hay que} + infinitif & \esp{Hay que cuidar la salud.} & Il faut prendre soin de la santé. \\
\hline
\esp{es importante que} + subjonctif & \esp{Es importante que hables.} & Il est important que tu parles. \\
\hline
\end{tabularx}
\end{center}
Remarque la nuance : l'impératif est direct et s'adresse à une personne précise ; \esp{hay que} est impersonnel (comme « il faut » en français) et vise tout le monde ; \esp{deber} exprime un devoir moral personnel.
\end{aretenir}

\courssub{Carte mentale du lexique}

Pour mémoriser le vocabulaire du texte, rien de tel qu'une carte mentale : elle organise les mots autour d'une idée centrale et te permet de les réutiliser dans tes propres productions.

\begin{popfigure}
\begin{tikzpicture}[mindmap, grow cyclic, every node/.style={concept, execute at begin node=\hskip0pt}, concept color=popblue!60, text=white,
level 1/.append style={level distance=3.4cm, sibling angle=120},
level 2/.append style={level distance=2.1cm, sibling angle=45}]
\node[concept, concept color=poporange, font=\bfseries] {conductas de riesgo}
child[concept color=popblue] { node[concept] {sustancias}
  child { node[concept, font=\small] {tabaco} }
  child { node[concept, font=\small] {alcohol} }
  child { node[concept, font=\small] {drogas} }
}
child[concept color=poppurple] { node[concept] {comportamiento}
  child { node[concept, font=\small] {violencia} }
  child { node[concept, font=\small] {trampa} }
  child { node[concept, font=\small] {acoso} }
}
child[concept color=popgreen] { node[concept] {soluciones}
  child { node[concept, font=\small] {respeto} }
  child { node[concept, font=\small] {diálogo} }
  child { node[concept, font=\small] {deporte} }
};
\end{tikzpicture}
\legende{Carte mentale : les conduites à risque et les solutions}
\end{popfigure}

\courssub{Méthode : lire un texte de sensibilisation}

\begin{methode}[Lire et comprendre un texte de sensibilisation]
\begin{enumerate}
\item \textbf{Identifier le thème et le destinataire.} Lis le titre et la première phrase : de quoi parle le texte ? À qui s'adresse-t-il ? Ici : les conduites à risque, destinataire = les élèves (tutoiement, \esp{nuestra escuela}).
\item \textbf{Repérer les problèmes dénoncés.} Souligne les noms et les verbes négatifs : \esp{fuman, beben, hacen trampa, maltratan}. Fais une liste.
\item \textbf{Souligner les conseils.} Cherche les impératifs (\esp{di, no fumes, actúa}), \esp{deber} + infinitif, \esp{hay que}, \esp{es importante que}. Ce sont les solutions proposées.
\item \textbf{Dégager le message final.} La dernière phrase contient souvent la morale du texte : \esp{Un joven responsable construye una sociedad fuerte} — un jeune responsable construit une société forte.
\end{enumerate}
\end{methode}

\courssub{Réflexion personnelle et expression orale}

\begin{experience}[Et dans ton lycée ?]
En groupes de trois ou quatre, discutez en espagnol à partir des questions suivantes :
\begin{itemize}
\item \esp{¿Y en tu colegio, existen estas conductas?} (Et dans ton lycée, est-ce que ces conduites existent ?)
\item \esp{¿Cuál es la conducta más peligrosa? ¿Por qué?} (Quelle est la conduite la plus dangereuse ? Pourquoi ?)
\item \esp{¿Qué consejos darías a un amigo que fuma?} (Quels conseils donnerais-tu à un ami qui fume ?)
\end{itemize}
Modèles de réponses : \esp{En mi colegio, algunos alumnos hacen trampa en los exámenes.} « Dans mon lycée, certains élèves trichent aux examens. » — \esp{Yo diría a mi amigo: no fumes, debes cuidar tu salud.} « Je dirais à mon ami : ne fume pas, tu dois prendre soin de ta santé. »
\end{experience}

\begin{savaistu}[Les slogans de prévention dans le monde hispanophone]
Dans le monde hispanophone, les campagnes de prévention utilisent souvent des slogans courts à l'impératif, faciles à retenir. Le plus célèbre est \esp{Si bebes, no conduzcas} « Si tu bois, ne conduis pas », affiché sur les routes d'Espagne et d'Amérique latine. On trouve aussi \esp{Di no a las drogas} « Dis non à la drogue » et \esp{El deporte es vida} « Le sport, c'est la vie ». Au Cameroun, pays voisin de la Guinée équatoriale — seul pays hispanophone d'Afrique —, ces slogans circulent aussi dans les campagnes de santé publique. Retiens la structure : \esp{si} + présent de l'indicatif, puis impératif négatif. Tu pourras l'imiter pour créer tes propres slogans : \esp{Si estudias, no hagas trampa} « Si tu étudies, ne triche pas » !
\end{savaistu}

\begin{attention}[Erreurs fréquentes des francophones]
\begin{itemize}
\item Ne traduis pas « Il faut » par \esp{es necesario fumar} quand tu veux dire « il faut fumer » : la tournure la plus naturelle est \esp{hay que} + infinitif : \esp{Hay que evitar las drogas} « Il faut éviter la drogue ».
\item À l'impératif négatif, le pronom personnel se place \textbf{avant} le verbe : \esp{No lo hagas} « Ne le fais pas », et jamais \esp{no hagas lo}.
\item \esp{Deber} + infinitif exprime le devoir / l'obligation morale : \esp{Debes estudiar} « Tu dois étudier ». Ne confonds pas avec \esp{deber de} + infinitif, qui exprime la probabilité : \esp{Debe de estar enfermo} « Il doit être malade » (il est probablement malade).
\item Attention à l'accent : \esp{salud} ne prend pas d'accent. \esp{di} (impératif de \esp{decir}) ne prend pas d'accent non plus — alors que \esp{dí} (avec accent) est le passé simple de \esp{dar} (« je donnai ») et \esp{sí} (avec accent) signifie « oui ».
\end{itemize}
\end{attention}

\begin{exoresolu}[Premiers pas vers la production]
Énoncé. À partir du texte, complète les slogans suivants avec l'impératif ou la structure de conseil qui convient : 1) ¡... no a la violencia! (decir) 2) No ... trampa en los exámenes. (hacer, tú) 3) ... que respetar a los profesores. (il faut) 4) Es importante que ... deporte. (pratiquer, tú)
\tcblower
Solution détaillée.
\begin{enumerate}
\item \esp{¡Di no a la violencia!} « Dis non à la violence ! » — impératif affirmatif irrégulier de \esp{decir} (tú) : \esp{di}.
\item \esp{No hagas trampa en los exámenes.} « Ne triche pas aux examens. » — impératif négatif de \esp{hacer} : subjonctif présent \esp{hagas} précédé de \esp{no}.
\item \esp{Hay que respetar a los profesores.} « Il faut respecter les professeurs. » — tournure impersonnelle \esp{hay que} + infinitif.
\item \esp{Es importante que practiques deporte.} « Il est important que tu pratiques un sport. » — subjonctif présent de \esp{practicar} après \esp{es importante que}.
\end{enumerate}
\end{exoresolu}

\begin{aretenir}[Bilan de la section]
Le texte \esp{Jóvenes responsables, sociedad fuerte} est un article de sensibilisation : il dénonce les conduites à risque (tabac, alcool, drogue, triche, violence, acoso, mauvais usage des réseaux sociaux) et propose des comportements responsables (respect, dialogue, sport, parole aux adultes). Pour conseiller et interdire, l'espagnol utilise quatre outils que nous étudierons en détail dans les sections suivantes : l'\cle{impératif} affirmatif et négatif, \esp{deber} + infinitif, \esp{hay que} + infinitif et \esp{es importante que} + subjonctif. Garde en mémoire le message final : \esp{Un joven responsable construye una sociedad fuerte} — un jeune responsable construit une société forte.
\end{aretenir}

\coursec{Lexique thématique : les conduites à risque et les comportements responsables}

Dans la section précédente, nous avons lu et commenté le texte \esp{« Jóvenes responsables, sociedad fuerte »}. Ce texte de sensibilisation était riche en vocabulaire : il opposait sans cesse deux mondes, celui des \esp{conductas de riesgo} (fumer, tricher, se droguer) et celui des \esp{comportamientos responsables} (respecter, étudier, dialoguer). Pour bien comprendre ce type de texte — et surtout pour pouvoir en produire un toi-même à l'écrit comme à l'oral — il faut maintenant organiser, mémoriser et manipuler ce lexique. C'est l'objectif de cette section : construire ton « sac à mots » sur le thème des comportements, avec leur prononciation, leurs emplois, leurs pièges et leurs traductions.

\courssub{Observation : le lexique en contexte}

Avant toute liste de mots, observons ce vocabulaire dans un texte authentique. Lis ce court article de presse scolaire et souligne mentalement tous les mots qui désignent un comportement.

\begin{textebox}[Article de presse scolaire — \esp{El Eco del Instituto}]
\textbf{Una campaña de prevención en nuestro instituto}

Desde el lunes, nuestro instituto organiza una campaña de prevención contra las conductas de riesgo. En la entrada, un gran cartel muestra el eslogan de la semana: «Piensa en tu futuro: di no a las drogas».

La directora explica que algunos alumnos fuman en los baños, beben alcohol cerca del campo de fútbol o hacen trampa en los exámenes. Otros sufren el acoso escolar: son insultados o golpeados por sus compañeros. «La violencia y las malas compañías destruyen la convivencia», declara la directora. «Mentir, robar o drogarse puede parecer un juego, pero las consecuencias son graves: la adicción, la enfermedad, el fracaso escolar o la exclusión del grupo».

Para cambiar estas conductas, los profesores aconsejan a los alumnos que respeten las normas, que digan la verdad y que cuiden su salud. El profesor de educación física propone hacer deporte cada día: «El deporte ayuda a los jóvenes a sentirse fuertes y felices». La enfermera del centro advierte a los estudiantes de los peligros del tabaco y del alcohol, y les invita a dialogar con un adulto cuando tienen un problema.

Al final de la semana, cada clase debe crear su propio cartel con un eslogan original. Los mejores carteles serán expuestos en la biblioteca. «Sensibilizar a un joven es proteger a toda la sociedad», concluye la directora con una sonrisa.
\end{textebox}

\textbf{Glossaire :}
\begin{itemize}
\item \esp{el instituto} : le lycée, le collège
\item \esp{la entrada} : l'entrée
\item \esp{los baños} : les toilettes
\item \esp{golpear} : frapper
\item \esp{la convivencia} : la vie en communauté
\item \esp{grave} : grave (attention, se prononce « gra-bé »)
\item \esp{el tabaco} : le tabac
\item \esp{propio} : propre (au sens de « qui lui appartient »)
\item \esp{expuesto} : exposé (participe passé de \esp{exponer})
\item \esp{proteger} : protéger
\end{itemize}

\textbf{Questions de compréhension :}
\begin{enumerate}
\item ¿Qué organiza el instituto esta semana ?
\item ¿Dónde fuman algunos alumnos ?
\item Según la directora, ¿qué destruye la convivencia ?
\item ¿Qué propone el profesor de educación física ?
\item ¿Qué debe crear cada clase al final de la semana ?
\end{enumerate}

\begin{definition}[Conducta de riesgo]
Une \esp{conducta de riesgo} (« con-duic-ta dé riès-go ») est un comportement qui met en danger la santé, la sécurité ou l'intégration sociale d'une personne. Exemple : \esp{Fumar es una conducta de riesgo.} « Fumer est une conduite à risque. » Le contraire est \esp{un comportamiento responsable}, un comportement responsable, qui protège la personne et la communauté.
\end{definition}

\courssub{Champ 1 — Les conduites à risque}

Voici le vocabulaire du « mauvais chemin ». Remarque que beaucoup de ces mots sont des verbes à l'infinitif : en espagnol, comme en français, on nomme souvent un comportement par l'infinitif (\esp{fumar}, \esp{mentir}) ou par un nom (\esp{la violencia}, \esp{el acoso escolar}).

\begin{center}
\begin{tabularx}{\linewidth}{|X|X|X|}\hline
\rowcolor{popblueL} Español & Français & Exemple et traduction \\ \hline
\esp{fumar} & fumer & \esp{Fumar daña los pulmones.} « Fumer abîme les poumons. » \\ \hline
\esp{beber alcohol} & boire de l'alcool & \esp{Los jóvenes no deben beber alcohol.} « Les jeunes ne doivent pas boire d'alcool. » \\ \hline
\esp{drogarse} & se droguer & \esp{Se droga para olvidar sus problemas.} « Il se drogue pour oublier ses problèmes. » \\ \hline
\esp{hacer trampa} & tricher & \esp{Hizo trampa en el examen de matemáticas.} « Il a triché à l'examen de maths. » \\ \hline
\esp{el acoso escolar} & le harcèlement scolaire & \esp{El acoso escolar destruye la convivencia.} « Le harcèlement scolaire détruit la vie en communauté. » \\ \hline
\esp{la violencia} & la violence & \esp{La violencia nunca es la solución.} « La violence n'est jamais la solution. » \\ \hline
\esp{robar} & voler (dérober) & \esp{Robó dinero en la tienda del barrio.} « Il a volé de l'argent dans la boutique du quartier. » \\ \hline
\esp{mentir} & mentir & \esp{No mientas a tus padres.} « Ne mens pas à tes parents. » \\ \hline
\esp{las malas compañías} & les mauvaises fréquentations & \esp{Las malas compañías llevan a las drogas.} « Les mauvaises fréquentations mènent à la drogue. » \\ \hline
\end{tabularx}
\end{center}

\begin{attention}[Pièges de traduction]
\begin{itemize}
\item \esp{Hacer trampa} signifie « tricher » (à un examen, à un jeu). Ne confonds pas avec le sens français de « trampa » = piège : « un piège » se dit \esp{una trampa}, mais \esp{hacer trampa} n'a jamais le sens de « poser un piège ». \esp{Juan hace trampa en los exámenes.} « Juan triche aux examens. »
\item \esp{Drogarse} est un verbe \cle{pronominal} : il se conjugue toujours avec le pronom réfléchi : \esp{me drogo, te drogas, se droga, nos drogamos, os drogáis, se drogan}. Ne dis jamais \esp{droga} seul pour « il se drogue ».
\item \esp{Mentir} est irrégulier : \esp{miento, mientes, miente, mentimos, mentís, mienten} (e $\to$ ie). \esp{Él siempre miente.} « Il ment toujours. »
\item \esp{Robar} = voler un objet ; attention au faux ami : « voler » (un oiseau vole) se dit \esp{volar}. \esp{El pájaro vuela.} « L'oiseau vole. »
\end{itemize}
\end{attention}

\courssub{Champ 2 — Les comportements responsables}

Face aux dangers, voici le vocabulaire positif, celui des solutions et des valeurs. C'est ce lexique que tu utiliseras pour \emph{conseiller} dans tes productions écrites.

\begin{center}
\begin{tabularx}{\linewidth}{|X|X|X|}\hline
\rowcolor{popgreenL} Español & Français & Exemple et traduction \\ \hline
\esp{respetar} & respecter & \esp{Respeta a tus profesores.} « Respecte tes professeurs. » \\ \hline
\esp{ayudar} & aider & \esp{Ayuda a tus compañeros.} « Aide tes camarades. » \\ \hline
\esp{estudiar} & étudier & \esp{Estudia cada día un poco.} « Étudie un peu chaque jour. » \\ \hline
\esp{decir la verdad} & dire la vérité & \esp{Di siempre la verdad.} « Dis toujours la vérité. » \\ \hline
\esp{cuidar la salud} & prendre soin de sa santé & \esp{Cuida tu salud: duerme bien.} « Prends soin de ta santé : dors bien. » \\ \hline
\esp{hacer deporte} & faire du sport & \esp{Hacemos deporte los sábados.} « Nous faisons du sport le samedi. » \\ \hline
\esp{obedecer las normas} & obéir aux règles & \esp{Obedece las normas del instituto.} « Obéis aux règles du lycée. » \\ \hline
\esp{dialogar} & dialoguer & \esp{Dialoga con tus padres.} « Dialogue avec tes parents. » \\ \hline
\end{tabularx}
\end{center}

\begin{propriete}[Petites règles de conjugaison utiles]
\begin{itemize}
\item \esp{Obedecer} est irrégulier à la 1\iere{} personne du présent : \esp{obedezco}, puis régulier : \esp{obedeces, obedece, obedecemos, obedecéis, obedecen}. \esp{Yo obedezco las normas.} « Moi, j'obéis aux règles. »
\item \esp{Decir} est très irrégulier : \esp{digo, dices, dice, decimos, decís, dicen}. \esp{Ella dice la verdad.} « Elle dit la vérité. »
\item \esp{Hacer} : \esp{hago, haces, hace, hacemos, hacéis, hacen}. \esp{Hago deporte los lunes.} « Je fais du sport le lundi. »
\end{itemize}
\end{propriete}

\courssub{Champ 3 — Sensibiliser et conseiller}

Ce troisième champ lexical est celui de l'\emph{action} : c'est le vocabulaire de celui qui prévient, qui conseille, qui interdit. Il est indispensable pour la tâche finale du module (créer un message de prévention).

\begin{itemize}
\item \esp{sensibilizar} : sensibiliser. \esp{La escuela sensibiliza a los alumnos.} « L'école sensibilise les élèves. »
\item \esp{prevenir} : prévenir. Verbe irrégulier conjugué comme \esp{venir} : \esp{prevengo, previenes, previene, prevenimos, prevenís, previenen}. \esp{Es mejor prevenir que curar.} « Mieux vaut prévenir que guérir. »
\item \esp{aconsejar} : conseiller. \esp{Aconsejo a los jóvenes que hagan deporte.} « Je conseille aux jeunes de faire du sport. » (remarque le subjonctif \esp{hagan} après \esp{aconsejar que}).
\item \esp{advertir} : avertir. Irrégulier comme \esp{mentir} : \esp{advierto, adviertes, advierte, advertimos, advertís, advierten}. \esp{Te advierto: es peligroso.} « Je t'avertis : c'est dangereux. »
\item \esp{prohibir} : interdire. \esp{El reglamento prohíbe fumar.} « Le règlement interdit de fumer. »
\item \esp{la campaña de prevención} : la campagne de prévention. \esp{Una campaña de prevención contra el acoso escolar.} « Une campagne de prévention contre le harcèlement scolaire. »
\item \esp{el cartel / el eslogan} : l'affiche / le slogan. \esp{El cartel muestra un eslogan fuerte.} « L'affiche montre un slogan percutant. »
\end{itemize}

\begin{attention}[Faux amis et calques]
\begin{itemize}
\item \esp{Advertir} = avertir, mettre en garde ; ce n'est pas « advertiser » au sens anglais de faire de la publicité.
\item « Sensibiliser » se dit bien \esp{sensibilizar}, mais attention : \esp{sensible} en espagnol veut dire « sensible » (émotif), et « sensible » au sens de « raisonnable » se dit \esp{sensato}.
\item \esp{Prohibir} porte un accent à certaines formes : \esp{prohíbo, prohíbes, prohíbe} (le hiatus se rompt).
\end{itemize}
\end{attention}

\courssub{Champ 4 — Les conséquences}

Toute conduite à risque entraîne des conséquences : ce champ lexical permet de faire le lien logique (\esp{por eso}, \esp{como consecuencia}) dans un texte argumentatif.

\begin{center}
\begin{tabularx}{\linewidth}{|X|X|X|}\hline
\rowcolor{poporangeL} Español & Français & Exemple et traduction \\ \hline
\esp{la enfermedad} & la maladie & \esp{El tabaco causa enfermedades graves.} « Le tabac cause des maladies graves. » \\ \hline
\esp{la adicción} & l'addiction & \esp{La adicción a las drogas es una enfermedad.} « L'addiction aux drogues est une maladie. » \\ \hline
\esp{el fracaso escolar} & l'échec scolaire & \esp{Hacer trampa lleva al fracaso escolar.} « Tricher mène à l'échec scolaire. » \\ \hline
\esp{la exclusión} & l'exclusion & \esp{La violencia provoca la exclusión del grupo.} « La violence provoque l'exclusion du groupe. » \\ \hline
\esp{el castigo / la sanción} & la punition / la sanction & \esp{Recibió una sanción de tres días.} « Il a reçu une sanction de trois jours. » \\ \hline
\end{tabularx}
\end{center}

\begin{exemplebox}[Deux phrases modèles à retenir]
\begin{itemize}
\item \esp{El acoso escolar destruye la convivencia.} « Le harcèlement scolaire détruit la vie en communauté. »
\item \esp{Aconsejo a los jóvenes que hagan deporte.} « Je conseille aux jeunes de faire du sport. »
\end{itemize}
\end{exemplebox}

\courssub{Expressions utiles pour le dialogue de sensibilisation}

Pour le jeu de rôle et l'expression orale, mémorise ces quatre formules toutes faites (nous étudierons la grammaire de l'impératif dans les sections 3 et 4) :

\begin{itemize}
\item \esp{¡Ten cuidado!} « Fais attention ! » (de \esp{tener}, impératif irrégulier \esp{ten})
\item \esp{Es peligroso.} « C'est dangereux. » (féminin : \esp{Es peligrosa.})
\item \esp{Piensa en tu futuro.} « Pense à ton avenir. » (de \esp{pensar}, e $\to$ ie)
\item \esp{Pide ayuda.} « Demande de l'aide. » (de \esp{pedir}, e $\to$ i)
\end{itemize}

\courssub{Carte mentale : les deux chemins}

La figure suivante résume l'opposition centrale de cette leçon : à gauche, les conduites à risque ; à droite, les comportements responsables ; au centre, le rôle de la sensibilisation, qui fait passer de l'un à l'autre.

\begin{popfigure}
\begin{tikzpicture}[font=\small]
\node[draw=popink, fill=poppinkL, rounded corners, align=center, text width=4.2cm, minimum height=3.6cm] (riesgo) at (0,0)
{\textbf{Conductas de riesgo ✗}\\[2pt]
fumar, beber alcohol\\
drogarse, hacer trampa\\
el acoso escolar, la violencia\\
robar, mentir\\
las malas compañías};
\node[draw=popgreen, fill=popgreenL, rounded corners, align=center, text width=4.2cm, minimum height=3.6cm] (resp) at (9.2,0)
{\textbf{Comportamientos responsables ✓}\\[2pt]
respetar, ayudar, estudiar\\
decir la verdad\\
cuidar la salud, hacer deporte\\
obedecer las normas\\
dialogar};
\node[draw=popblue, fill=popblueL, rounded corners, align=center, text width=3.4cm] (sens) at (4.6,2.6)
{\textbf{Sensibilizar}\\ prevenir, aconsejar\\ advertir, prohibir};
\draw[-{Stealth[length=3mm]}, very thick, poporange] (riesgo.east) -- (resp.west)
node[midway, below, align=center, popdark]{\emph{cambiar}};
\draw[-{Stealth[length=2.5mm]}, thick, popblue] (sens.south west) -- (riesgo.north east);
\draw[-{Stealth[length=2.5mm]}, thick, popblue] (sens.south east) -- (resp.north west);
\node[draw=popgold, fill=popgoldL, rounded corners, align=center, text width=3.6cm] (cons) at (4.6,-2.6)
{\textbf{Consecuencias}\\ enfermedad, adicción\\ fracaso escolar, exclusión\\ castigo, sanción};
\draw[-{Stealth[length=2.5mm]}, thick, popink] (riesgo.south east) -- (cons.north west);
\end{tikzpicture}
\legende{Carte mentale lexicale : sensibiliser, c'est transformer les conduites à risque en comportements responsables.}
\end{popfigure}

\courssub{Exercice résolu et entraînement}

\begin{exoresolu}[Complète avec le mot du lexique qui convient]
Complète chaque phrase avec un mot de cette leçon : \esp{acoso escolar, adicción, aconseja, campaña, trampa, verdad}.
\begin{enumerate}
\item Los alumnos no deben hacer \dots{} en los exámenes.
\item El \dots{} es un problema grave: algunos niños tienen miedo de ir a la escuela.
\item Mi madre me \dots{} que cuide mi salud.
\item La \dots{} a las drogas destruye a las familias.
\item Di siempre la \dots{}, no mientas.
\item El instituto organiza una \dots{} de prevención contra el tabaco.
\end{enumerate}
\tcblower
\textbf{Solution détaillée :}
\begin{enumerate}
\item \esp{trampa} : \esp{hacer trampa} = tricher (à un examen). « Les élèves ne doivent pas tricher aux examens. »
\item \esp{acoso escolar} : le contexte (\esp{tener miedo de ir a la escuela}) évoque le harcèlement scolaire.
\item \esp{aconseja} : verbe \esp{aconsejar}, 3\ieme{} personne du singulier (\esp{mi madre}). « Ma mère me conseille de prendre soin de ma santé. »
\item \esp{adicción} : \esp{la adicción a las drogas}. « L'addiction aux drogues détruit les familles. »
\item \esp{verdad} : \esp{decir la verdad}, expression figée. « Dis toujours la vérité, ne mens pas. »
\item \esp{campaña} : \esp{una campaña de prevención}. « Le lycée organise une campagne de prévention contre le tabac. »
\end{enumerate}
\end{exoresolu}

\begin{exercice}[À ton tour]
\begin{enumerate}
\item Traduis en espagnol : a) « Fumer est dangereux pour la santé. » b) « Les mauvaises fréquentations mènent à la violence. » c) « Je conseille à mon ami de faire du sport. » d) « Demande de l'aide à un professeur. »
\item Classe ces mots en deux colonnes (\esp{riesgo} / \esp{responsable}) : \esp{robar, dialogar, mentir, ayudar, drogarse, respetar, hacer trampa, estudiar}.
\item Rédige deux slogans de prévention en espagnol avec le vocabulaire de la leçon (un contre le tabac, un contre le harcèlement scolaire).
\end{enumerate}
\end{exercice}

\begin{corrige}[Exercice « À ton tour »]
\begin{enumerate}
\item a) \esp{Fumar es peligroso para la salud.} b) \esp{Las malas compañías llevan a la violencia.} c) \esp{Aconsejo a mi amigo que haga deporte.} (subjonctif après \esp{aconsejar que}) d) \esp{Pide ayuda a un profesor.}
\item Riesgo : \esp{robar, mentir, drogarse, hacer trampa}. Responsable : \esp{dialogar, ayudar, respetar, estudiar}.
\item Exemples possibles : \esp{«No fumes: tu salud es tu futuro.»} ; \esp{«Contra el acoso escolar, dialoga y respeta.»}
\end{enumerate}
\end{corrige}

\begin{savaistu}[Le lexique de la santé au Bac]
Au baccalauréat, le lexique de la santé et de la vie en société revient très souvent dans les sujets d'expression écrite : lettre de conseil à un ami en difficulté, article de sensibilisation pour le journal du lycée, dialogue de prévention. Les mots de cette leçon (\esp{conductas de riesgo, campaña de prevención, acoso escolar, adicción}) sont donc un investissement direct pour l'examen. Au Cameroun, les sujets aiment aussi les contextes locaux : n'hésite pas à situer tes exemples à Yaoundé ou à Douala, par exemple \esp{En mi barrio de Yaoundé, los jóvenes organizan una campaña contra las drogas.} « Dans mon quartier de Yaoundé, les jeunes organisent une campagne contre la drogue. »
\end{savaistu}

\begin{aretenir}[L'essentiel du lexique]
\begin{itemize}
\item Une \esp{conducta de riesgo} met en danger la santé, la sécurité ou l'intégration sociale ; un \esp{comportamiento responsable} protège la personne et la communauté.
\item Quatre champs à maîtriser : les risques (\esp{fumar, drogarse, hacer trampa, el acoso escolar...}), les comportements positifs (\esp{respetar, ayudar, dialogar...}), l'action de sensibiliser (\esp{prevenir, aconsejar, advertir, prohibir, la campaña, el eslogan}) et les conséquences (\esp{la adicción, el fracaso escolar, la exclusión, la sanción}).
\item Pièges : \esp{hacer trampa} = tricher ; \esp{drogarse} est pronominal ; \esp{mentir} et \esp{advertir} ont un radical en \esp{ie} ; \esp{robar} $\neq$ \esp{volar} ; \esp{prevenir} se conjugue \esp{prevengo} (comme \esp{venir}).
\item Formules du dialogue : \esp{¡Ten cuidado!}, \esp{Es peligroso}, \esp{Piensa en tu futuro}, \esp{Pide ayuda}.
\end{itemize}
\end{aretenir}

Maintenant que ce lexique est bien installé, nous pouvons passer à la grammaire qui permet de le mettre en action : l'impératif affirmatif, pour donner des conseils et des consignes.

\coursec{Grammaire 1 : l'impératif affirmatif (tú, vosotros, usted)}

Dans la section précédente, nous avons construit le lexique des conduites à risque (\esp{el consumo de drogas}, \esp{el acoso escolar}, \esp{la violencia}) et des comportements responsables (\esp{respetar}, \esp{dialogar}, \esp{cuidar la salud}). Or, pour \emph{sensibiliser} et \emph{conseiller}, le vocabulaire ne suffit pas : il faut un mode verbal capable de s'adresser directement à quelqu'un pour lui dire quoi faire. Ce mode, c'est \cle{l'impératif}. C'est le mode des affiches de prévention, des slogans, des consignes du professeur et des conseils des parents. Dans cette section, nous étudions l'\cle{impératif affirmatif} ; la forme négative (\esp{No fumes}) sera traitée dans la section suivante.

\courssub{Observation : l'impératif dans le texte de sensibilisation}

Reprenons les slogans du texte \emph{« Jóvenes responsables, sociedad fuerte »} et observons les formes en gras :

\begin{exemplebox}[Slogans extraits du texte]
\esp{¡Di no a las drogas!} « Dis non à la drogue ! » (verbe \esp{decir})\par
\esp{¡Respeta a los demás!} « Respecte les autres ! » (verbe \esp{respetar})\par
\esp{¡Cuida tu salud!} « Prends soin de ta santé ! » (verbe \esp{cuidar})\par
\esp{¡Habla con tus padres!} « Parle avec tes parents ! » (verbe \esp{hablar})
\end{exemplebox}

Observons attentivement :

\begin{itemize}
\item \esp{respeta}, \esp{cuida}, \esp{habla} ressemblent à des formes déjà connues : ce sont exactement les formes de la \textbf{3\textsuperscript{e} personne du singulier du présent de l'indicatif} (\esp{él habla}, \esp{él respeta}, \esp{él cuida}).
\item \esp{di} (de \esp{decir}) ne correspond à \emph{aucune} forme de l'indicatif : c'est une forme \textbf{irrégulière} qu'il faut mémoriser.
\item Toutes ces phrases commencent par \textbf{¡} et finissent par \textbf{!} : en espagnol, toute phrase exclamative ou impérative est encadrée par les deux signes.
\item Il n'y a \textbf{pas de sujet exprimé} : comme en français (« Dis non ! », « Respecte ! »), la personne à qui l'on parle est sous-entendue.
\end{itemize}

\begin{definition}[L'impératif]
L'\cle{impératif} est le mode qui sert à donner un \textbf{ordre}, un \textbf{conseil}, une \textbf{consigne}, une \textbf{interdiction} ou un \textbf{encouragement}. En espagnol comme en français, il ne s'emploie qu'aux personnes auxquelles on s'adresse directement : \esp{tú} (tu), \esp{vosotros} (vous, familier, Espagne), \esp{usted} / \esp{ustedes} (vous, formel).
\end{definition}

\courssub{Formation régulière de l'impératif affirmatif}

\begin{propriete}[Règle de formation — les trois personnes]
\begin{itemize}
\item \textbf{tú} : on prend la \textbf{3\textsuperscript{e} personne du singulier du présent de l'indicatif} : \esp{hablar → habla}, \esp{comer → come}, \esp{escribir → escribe}.
\item \textbf{vosotros} : on remplace le \textbf{-r} final de l'infinitif par \textbf{-d} : \esp{hablar → hablad}, \esp{comer → comed}, \esp{escribir → escribid}. Cette forme est \textbf{toujours régulière au niveau de la terminaison}, même pour les verbes à radical irrégulier (ex. \esp{ir → id}, \esp{ser → sed}).
\item \textbf{usted} : on prend la \textbf{3\textsuperscript{e} personne du singulier du subjonctif présent} : \esp{hablar → hable}, \esp{comer → coma}, \esp{escribir → escriba} (voyelle « inversée » : -ar donne -e ; -er/-ir donnent -a).
\end{itemize}
\end{propriete}

\begin{center}
\begin{tabularx}{\linewidth}{|X|X|X|X|}\hline
\rowcolor{popblueL} \textbf{Infinitif} & \textbf{tú} & \textbf{vosotros} & \textbf{usted} \\ \hline
hablar (parler) & habla & hablad & hable \\ \hline
comer (manger) & come & comed & coma \\ \hline
escribir (écrire) & escribe & escribid & escriba \\ \hline
respetar (respecter) & respeta & respetad & respete \\ \hline
estudiar (étudier) & estudia & estudiad & estudie \\ \hline
ayudar (aider) & ayuda & ayudad & ayude \\ \hline
\end{tabularx}
\end{center}

\begin{exemplebox}[Exemples traduits]
\esp{Habla con el consejero.} « Parle au conseiller d'orientation. »\par
\esp{Come fruta todos los días.} « Mange des fruits tous les jours. »\par
\esp{Escribe una carta a tu amigo.} « Écris une lettre à ton ami. »\par
\esp{Escuchad al profesor.} « Écoutez le professeur. » (vosotros, un groupe d'élèves)\par
\esp{Respete las normas, por favor.} « Respectez les règles, s'il vous plaît. » (usted, formel)
\end{exemplebox}

\begin{popfigure}
\begin{tikzpicture}[
  node distance=0.35cm,
  head/.style={fill=popblue, text=white, font=\bfseries\small, minimum width=2.6cm, minimum height=0.7cm, align=center},
  cell/.style={fill=popcream, draw=popline, font=\small, minimum width=2.6cm, minimum height=0.7cm, align=center},
  cellb/.style={fill=popblueL, draw=popline, font=\small, minimum width=2.6cm, minimum height=0.7cm, align=center},
  cellg/.style={fill=popgreenL, draw=popline, font=\small, minimum width=2.6cm, minimum height=0.7cm, align=center},
  cello/.style={fill=poporangeL, draw=popline, font=\small, minimum width=2.6cm, minimum height=0.7cm, align=center}
]
\node[head] (h1) at (0,0) {Infinitivo};
\node[head, right=of h1] (h2) {tú};
\node[head, right=of h2] (h3) {vosotros};
\node[head, right=of h3] (h4) {usted};
\node[cell, below=of h1] (a1) {hablar};
\node[cellb, below=of h2] (a2) {habla};
\node[cellg, below=of h3] (a3) {hablad};
\node[cello, below=of h4] (a4) {hable};
\node[cell, below=of a1] (b1) {comer};
\node[cellb, below=of b1.east -| h2] (b2) {come};
\node[cellg, below=of b2] (b3) {comed};
\node[cello, below=of b3] (b4) {coma};
\node[cell, below=of b1] (c1) {escribir};
\node[cellb, below=of c1.east -| h2] (c2) {escribe};
\node[cellg, below=of c2] (c3) {escribid};
\node[cello, below=of c3] (c4) {escriba};
\node[cell, below=of c1] (d1) {decir};
\node[cellb, below=of d1.east -| h2] (d2) {\textbf{di}};
\node[cellg, below=of d2] (d3) {decid};
\node[cello, below=of d3] (d4) {diga};
\node[cell, below=of d1] (e1) {hacer};
\node[cellb, below=of e1.east -| h2] (e2) {\textbf{haz}};
\node[cellg, below=of e2] (e3) {haced};
\node[cello, below=of e3] (e4) {haga};
\node[cell, below=of e1] (f1) {venir};
\node[cellb, below=of f1.east -| h2] (f2) {\textbf{ven}};
\node[cellg, below=of f2] (f3) {venid};
\node[cello, below=of f3] (f4) {venga};
\node[cell, below=of f1] (g1) {ir};
\node[cellb, below=of g1.east -| h2] (g2) {\textbf{ve}};
\node[cellg, below=of g2] (g3) {id};
\node[cello, below=of g3] (g4) {vaya};
\node[cell, below=of g1] (h1b) {ser};
\node[cellb, below=of h1b.east -| h2] (h2b) {\textbf{sé}};
\node[cellg, below=of h2b] (h3b) {sed};
\node[cello, below=of h3b] (h4b) {sea};
\end{tikzpicture}
\legende{Tableau de formation de l'impératif affirmatif : en bleu la colonne \esp{tú} (formes en gras = irrégulières), en vert \esp{vosotros} (terminaison -d toujours régulière), en orange \esp{usted} (subjonctif présent).}
\end{popfigure}

\courssub{Les huit impératifs irréguliers de « tú »}

Huit verbes très fréquents ont une forme d'impératif \esp{tú} \textbf{courte et irrégulière}. Ils doivent être connus \textbf{par cœur}, car on les retrouve dans tous les messages de prévention.

\begin{propriete}[Les 8 impératifs irréguliers de tú]
\begin{center}
\begin{tabular}{|l|l|l|l|}\hline
\rowcolor{popgreenL} \textbf{Infinitif} & \textbf{Impératif tú} & \textbf{Vosotros (régulier)} & \textbf{Exemple (tú) traduit} \\ \hline
decir (dire) & \textbf{di} & decid & \esp{¡Di la verdad!} « Dis la vérité ! » \\ \hline
hacer (faire) & \textbf{haz} & haced & \esp{¡Haz los deberes!} « Fais tes devoirs ! » \\ \hline
ir (aller) & \textbf{ve} & id & \esp{¡Ve a casa!} « Va à la maison ! » \\ \hline
poner (mettre) & \textbf{pon} & poned & \esp{¡Pon la mesa!} « Mets la table ! » \\ \hline
salir (sortir) & \textbf{sal} & salid & \esp{¡Sal de aquí!} « Sors d'ici ! » \\ \hline
ser (être) & \textbf{sé} & sed & \esp{¡Sé responsable!} « Sois responsable ! » \\ \hline
tener (avoir) & \textbf{ten} & tened & \esp{¡Ten cuidado!} « Fais attention ! » \\ \hline
venir (venir) & \textbf{ven} & venid & \esp{¡Ven aquí!} « Viens ici ! » \\ \hline
\end{tabular}
\end{center}
Moyen mnémotechnique : \textbf{D}i, \textbf{H}az, \textbf{V}e, \textbf{P}on, \textbf{S}al, \textbf{S}é, \textbf{T}en, \textbf{V}en — des formes courtes de 2 à 3 lettres. Notez l'accent sur \esp{é} pour \esp{sé} (différencie du pronom \esp{se}).
\end{propriete}

\begin{attention}[L'erreur classique du francophone]
Ne dites jamais \esp{¡Dices no!} ni \esp{¡Haces los deberes!} : ce sont des formes de l'\emph{indicatif} avec \esp{-s} (2\textsuperscript{e} personne). À l'impératif, il faut la forme irrégulière : \esp{¡Di no!}, \esp{¡Haz los deberes!}. De même, n'écrivez pas \esp{¡Venir aquí!} : l'infinitif s'utilise pour des consignes \emph{impersonnelles} sur les panneaux (\esp{No fumar}, \esp{No entrar}), mais pour s'adresser directement à quelqu'un (tú/usted), on utilise l'impératif : \esp{¡Ven aquí!}.
\end{attention}

\courssub{Verbes à changement de voyelle : piensa, vuelve, pide, duerme}

Les verbes à changement de voyelle (\esp{e→ie}, \esp{o→ue}, \esp{e→i}) conservent ce changement à l'impératif \esp{tú}, puisque la forme suit le présent de l'indicatif. En revanche, \esp{vosotros} reste \textbf{toujours régulier} (pas de changement de radical).

\begin{center}
\begin{tabularx}{\linewidth}{|X|X|X|X|}\hline
\rowcolor{poporangeL} \textbf{Infinitif} & \textbf{tú} & \textbf{vosotros} & \textbf{usted} \\ \hline
pensar (e→ie) & \textbf{piensa} & pensad & piense \\ \hline
volver (o→ue) & \textbf{vuelve} & volved & vuelva \\ \hline
pedir (e→i) & \textbf{pide} & pedid & pida \\ \hline
dormir (o→ue) & \textbf{duerme} & dormid & duerma \\ \hline
\end{tabularx}
\end{center}

\begin{exemplebox}[Exemples traduits]
\esp{Piensa antes de actuar.} « Réfléchis avant d'agir. »\par
\esp{Vuelve temprano a casa.} « Rentre tôt à la maison. »\par
\esp{Pide ayuda si la necesitas.} « Demande de l'aide si tu en as besoin. »\par
\esp{Duerme ocho horas cada noche.} « Dors huit heures chaque nuit. »\par
\esp{Pensad en vuestro futuro.} « Pensez à votre avenir. » (vosotros : régulier, sans changement de voyelle)
\end{exemplebox}

\courssub{La place des pronoms : soudés après le verbe}

Grande différence avec le français : à l'impératif \textbf{affirmatif}, les pronoms compléments (COD, COI) et réfléchis se placent \textbf{après le verbe et soudés à lui} en un seul mot. Le français fait l'inverse avec une liaison (« Dis-moi », « Lève-toi »).

\begin{exemplebox}[Verbe + pronom soudé]
\esp{¡Dime!} « Dis-moi ! » (decir + me)\par
\esp{¡Hazlo!} « Fais-le ! » (hacer + lo)\par
\esp{¡Levántate!} « Lève-toi ! » (levantar + te)\par
\esp{¡Ayúdanos!} « Aide-nous ! » (ayudar + nos)\par
\esp{¡Dímelo!} « Dis-le-moi ! » (decir + me + lo)
\end{exemplebox}

Quand on soude un ou deux pronoms, le verbe s'allonge et l'accent tonique tomberait au mauvais endroit : on ajoute donc un \textbf{accent graphique} sur la voyelle qui portait déjà l'accent dans la forme simple : \esp{dime} mais \esp{dímelo} ; \esp{habla} mais \esp{háblame} ; \esp{levanta} mais \esp{levántate}. Pour \esp{vosotros} + pronom réfléchi, le \textbf{-d} tombe : \esp{levantad + os → levantaos} « levez-vous ».

\begin{methode}[Placer correctement les pronoms à l'impératif affirmatif]
\begin{enumerate}
\item Conjugue le verbe à l'impératif : \esp{dí} (decir, tú).
\item Ajoute les pronoms \textbf{après}, soudés, dans l'ordre : pronom indirect (me, te, nos) puis pronom direct (lo, la, los, las) : \esp{di + me + lo}.
\item Compte les syllabes et place l'accent graphique sur la voyelle tonique d'origine : \esp{\textbf{dí}-me-lo}.
\item Vérifie à voix haute : \esp{¡Dímelo!} « Dis-le-moi ! »
\end{enumerate}
\end{methode}

\courssub{Emplois de l'impératif affirmatif}

\begin{aretenir}[Quand utilise-t-on l'impératif ?]
\begin{itemize}
\item \textbf{Conseiller} : \esp{¡Estudia mucho!} « Travaille dur ! », \esp{¡Cuida tu salud!} « Prends soin de ta santé ! »
\item \textbf{Donner un ordre ou une consigne} : \esp{¡Silencio, por favor!} « Silence, s'il vous plaît ! », \esp{¡Abrid los libros!} « Ouvrez vos livres ! »
\item \textbf{Encourager} : \esp{¡Ánimo, tú puedes!} « Courage, tu peux y arriver ! », \esp{¡Sigue así!} « Continue comme ça ! »
\item \textbf{Sensibiliser} (affiches, slogans) : \esp{¡Protege el medio ambiente!} « Protège l'environnement ! », \esp{¡Sé solidario!} « Sois solidaire ! »
\end{itemize}
\end{aretenir}

\begin{savaistu}[« Vosotros » : seulement en Espagne]
La forme \esp{vosotros} (et son impératif en \esp{-d} : \esp{hablad}, \esp{comed}) n'est utilisée qu'en \textbf{Espagne}. Dans toute l'Amérique latine — et en Guinée équatoriale, voisine du Cameroun — on s'adresse à un groupe, même d'amis, avec \esp{ustedes}, suivi de la forme du subjonctif : \esp{¡No fumen, jóvenes!} « Ne fumez pas, les jeunes ! », \esp{¡Escuchen!} « Écoutez ! ». Au Cameroun, les deux formes sont acceptées dans les examens, mais il faut savoir reconnaître les deux.
\end{savaistu}

\begin{exoresolu}[Transforme les consignes à l'impératif]
Énoncé — Le directeur d'un lycée de Yaoundé veut afficher des consignes. Mets les verbes entre parenthèses à l'impératif affirmatif, à la personne indiquée.\par
1. (hablar, tú) \_\_\_\_\_ con tus profesores si tienes problemas.\par
2. (hacer, tú) \_\_\_\_\_ deporte tres veces por semana.\par
3. (respetar, vosotros) \_\_\_\_\_ a vuestros compañeros.\par
4. (venir, tú) \_\_\_\_\_ a la reunión de padres.\par
5. (decir + me, tú) \_\_\_\_\_ la verdad.
\tcblower
Solution détaillée :\par
1. \esp{Habla} : verbe régulier en -ar, tú = 3\textsuperscript{e} pers. sing. du présent. « Parle à tes professeurs si tu as des problèmes. »\par
2. \esp{Haz} : impératif irrégulier de \esp{hacer} (à mémoriser). « Fais du sport trois fois par semaine. »\par
3. \esp{Respetad} : vosotros = infinitif -r → -d, toujours régulier. « Respectez vos camarades. »\par
4. \esp{Ven} : impératif irrégulier de \esp{venir}. « Viens à la réunion des parents. »\par
5. \esp{Dime} : impératif irrégulier \esp{di} + pronom \esp{me} soudé ; pas d'accent graphique ici car \esp{di-me} (deux syllabes, terminé par voyelle) s'accentue naturellement sur le premier son. « Dis-moi la vérité. »
\end{exoresolu}

\begin{exercice}[À toi de jouer : conseils à un camarade]
Mets les verbes à l'impératif affirmatif (tú), sauf indication contraire, puis traduis chaque phrase en français.\par
1. (cuidar) \_\_\_\_\_ tu salud : no fumes.\par
2. (ser) \_\_\_\_\_ puntual en clase.\par
3. (ir) \_\_\_\_\_ al médico si estás enfermo.\par
4. (pensar) \_\_\_\_\_ antes de publicar en las redes sociales.\par
5. (levantar + te) \_\_\_\_\_ temprano los días de examen.\par
6. (escuchar, usted) \_\_\_\_\_ los consejos del médico.
\end{exercice}

\begin{corrige}[Exercice « Conseils à un camarade »]
1. \esp{Cuida tu salud : no fumes.} « Prends soin de ta santé : ne fume pas. »\par
2. \esp{Sé puntual en clase.} « Sois ponctuel en classe. » (irrégulier : ser → sé)\par
3. \esp{Ve al médico si estás enfermo.} « Va chez le médecin si tu es malade. » (irrégulier : ir → ve)\par
4. \esp{Piensa antes de publicar en las redes sociales.} « Réfléchis avant de publier sur les réseaux sociaux. » (changement e→ie)\par
5. \esp{Levántate temprano los días de examen.} « Lève-toi tôt les jours d'examen. » (pronom soudé + accent graphique sur \esp{á})\par
6. \esp{Escuche los consejos del médico.} « Écoutez les conseils du médecin. » (usted = subjonctif présent)
\end{corrige}

\begin{aretenir}[Bilan de la section]
\begin{itemize}
\item \esp{tú} : 3\textsuperscript{e} pers. sing. du présent (\esp{habla, come, escribe}) + 8 irréguliers : \esp{di, haz, ve, pon, sal, sé, ten, ven}.
\item \esp{vosotros} : infinitif -r → -d, toujours régulier (\esp{hablad, pensad, venid, id, sed}).
\item \esp{usted} : subjonctif présent (\esp{hable, coma, escriba}).
\item Les pronoms se soudent \textbf{après} le verbe, avec accent graphique si nécessaire : \esp{¡Dímelo!}, \esp{¡Levántate!}, \esp{¡Levantaos!}.
\item Emplois : conseiller, ordonner, encourager, sensibiliser — le mode roi des messages de prévention.
\end{itemize}
\end{aretenir}

Maintenant que nous savons former et employer l'impératif affirmatif pour conseiller, il reste à apprendre à \emph{interdire} et à nuancer le conseil : ce sera l'objet de la section suivante, consacrée à l'impératif négatif (\esp{No fumes}) et aux structures \esp{deber + infinitivo}, \esp{hay que} et \esp{es importante que + subjuntivo}.

\coursec{Grammaire 2 : l'impératif négatif et les structures du conseil et de l'obligation}

Dans la section précédente, tu as appris à donner des ordres et des conseils \emph{positifs} avec l'impératif affirmatif : \esp{¡Respeta a los demás!}, \esp{¡Haz los deberes!}, \esp{¡Ayuda a tus amigos!}. Mais une campagne de prévention ne dit pas seulement ce qu'il faut \emph{faire} : elle dit surtout ce qu'il ne faut \emph{pas faire}. Regarde les slogans affichés dans les couloirs du lycée : \esp{No fumes}, \esp{No hagas trampa}, \esp{No bebas alcohol}. Comment sont-ils construits ? C'est toute la question de cette section. Nous verrons ensuite trois structures indispensables pour conseiller avec nuance : \esp{deber + infinitif}, \esp{hay que + infinitif} et \esp{es importante que + subjonctif}.

\courssub{Observation : comment dit-on « ne fais pas » en espagnol ?}

\begin{textebox}[Slogans de la campagne de sensibilisation du lycée]
¡No fumes en el patio!\\
¡No hagas trampa en los exámenes!\\
¡No bebas alcohol antes de conducir!\\
¡No te drogues!\\
¡No digas mentiras a tus padres!\\
¡No levantéis la voz en clase!\\
¡No tiren papeles al suelo!
\end{textebox}

\textbf{Glossaire :} \esp{hacer trampa} = tricher ; \esp{conducir} = conduire ; \esp{decir mentiras} = mentir ; \esp{levantar la voz} = élever la voix ; \esp{tirar papeles al suelo} = jeter des papiers par terre.

Observe attentivement ces slogans et réponds aux questions suivantes :

\begin{enumerate}
\item Quel mot revient au début de chaque phrase ?
\item Les verbes \esp{fumes}, \esp{hagas}, \esp{bebas}, \esp{drogues}, \esp{digas} ressemblent-ils à l'impératif affirmatif (\esp{fuma}, \esp{haz}, \esp{bebe}, \esp{di}) ou au subjonctif présent (\esp{fumes}, \esp{hagas}, \esp{bebas}, \esp{digas}) ?
\item Où se place le pronom \esp{te} dans \esp{No te drogues} : avant ou après le verbe ?
\end{enumerate}

Réponses attendues : le mot \esp{no} ouvre chaque phrase ; les formes verbales sont exactement celles du \textbf{subjonctif présent} (et non de l'impératif affirmatif) ; le pronom \esp{te} se place \textbf{avant} le verbe, détaché de lui.

\begin{propriete}[Règle de l'impératif négatif]
L'impératif négatif espagnol se construit avec \cle{no + subjonctif présent}, à \textbf{toutes} les personnes :
\begin{itemize}
\item \textbf{tú} : \esp{no hables}, \esp{no comas}, \esp{no escribas} (ne parle pas, ne mange pas, n'écris pas) ;
\item \textbf{vosotros} : \esp{no habléis}, \esp{no comáis}, \esp{no escribáis} ;
\item \textbf{usted} : \esp{no hable}, \esp{no coma}, \esp{no escriba} ;
\item \textbf{ustedes} : \esp{no hablen}, \esp{no coman}, \esp{no escriban} ;
\item \textbf{nosotros} : \esp{no hablemos}, \esp{no comamos}, \esp{no escribamos} (ne parlons pas, ne mangeons pas, n'écrivons pas).
\end{itemize}
Les \textbf{pronoms} (réfléchis, compléments) se placent \textbf{avant le verbe}, jamais soudés après : \esp{¡No lo hagas!} « Ne le fais pas ! », \esp{No te levantes} « Ne te lève pas ». C'est l'inverse de l'impératif affirmatif, où les pronoms se soudent après le verbe (\esp{¡Hazlo!}, \esp{¡Levántate!}).
\end{propriete}

Rappel de la formation du subjonctif présent : on prend la première personne du présent de l'indicatif (\esp{hago}, \esp{digo}, \esp{salgo}, \esp{tengo}) et on inverse les voyelles thématiques : \textbf{-a} pour les verbes en \esp{-er/-ir}, \textbf{-e} pour les verbes en \esp{-ar}.

\begin{center}
\begin{tabular}{|l|l|l|l|}
\hline
\rowcolor{popblueL} Personne & \esp{fumar} (no\ldots) & \esp{beber} (no\ldots) & \esp{mentir} (no\ldots) \\
\hline
\esp{tú} & \esp{no fumes} & \esp{no bebas} & \esp{no mientas} \\
\hline
\esp{vosotros} & \esp{no fuméis} & \esp{no bebáis} & \esp{no mintáis} \\
\hline
\esp{usted} & \esp{no fume} & \esp{no beba} & \esp{no mienta} \\
\hline
\esp{ustedes} & \esp{no fumen} & \esp{no beban} & \esp{no mientan} \\
\hline
\esp{nosotros} & \esp{no fumemos} & \esp{no bebamos} & \esp{no mintamos} \\
\hline
\end{tabular}
\end{center}

Attention aux verbes irréguliers : le subjonctif garde l'irrégularité de la première personne de l'indicatif. Ainsi \esp{hacer} donne \esp{yo hago} → \esp{no hagas} ; \esp{decir} donne \esp{yo digo} → \esp{no digas} ; \esp{salir} donne \esp{yo salgo} → \esp{no salgas} ; \esp{tener} donne \esp{yo tengo} → \esp{no tengas} ; \esp{ir} donne \esp{no vayas} ; \esp{ser} donne \esp{no seas} (\esp{¡No seas egoísta!} « Ne sois pas égoïste ! »).

\begin{exemplebox}[Affirmatif contre négatif : la transformation complète]
\begin{itemize}
\item \esp{¡Hazlo!} « Fais-le ! » → \esp{¡No lo hagas!} « Ne le fais pas ! »
\item \esp{¡Dile la verdad!} « Dis-lui la vérité ! » → \esp{¡No le digas mentiras!} « Ne lui dis pas de mensonges ! »
\item \esp{¡Levántate temprano!} « Lève-toi tôt ! » → \esp{¡No te levantes tarde!} « Ne te lève pas tard ! »
\item \esp{¡Cómete la sopa!} « Mange ta soupe ! » → \esp{¡No te comas los caramelos!} « Ne mange pas les bonbons ! »
\item \esp{¡Escuchad al profesor!} « Écoutez le professeur ! » → \esp{¡No escuchéis música en clase!} « N'écoutez pas de musique en classe ! »
\end{itemize}
Tu remarques le double mouvement : le verbe change de forme (impératif → subjonctif) \emph{et} le pronom change de place (soudé après → détaché avant).
\end{exemplebox}

\begin{attention}[« No fumes » ou « no fumar » ?]
Ne confonds jamais les deux structures. \esp{No fumes} (subjonctif, à la deuxième personne) s'adresse à quelqu'un : « Ne fume pas ». En revanche, \esp{No fumar} avec l'\textbf{infinitif} est réservé aux \textbf{panneaux et affiches impersonnels} : \esp{No fumar} « Défense de fumer », \esp{No aparcar} « Stationnement interdit », \esp{No tirar basura} « Défense de jeter des ordures ». Dire \esp{no fumas} à un ami serait une faute de grammaire ; écrire \esp{no fumes} sur un panneau officiel serait une maladresse. Retiens : \textbf{personne précise = no + subjonctif ; panneau impersonnel = no + infinitif}.
\end{attention}

\courssub{deber + infinitif : le devoir personnel}

Pour conseiller sans donner un ordre sec, l'espagnol utilise le verbe \esp{deber} (régulier en \esp{-er}) suivi d'un \textbf{infinitif}. Cette structure exprime un devoir qui concerne \textbf{une personne précise} : elle se conjugue donc.

\begin{center}
\begin{tabular}{|l|l|l|}
\hline
\rowcolor{popgreenL} Personne & \esp{deber} & Exemple \\
\hline
\esp{yo} & \esp{debo} & \esp{Debo respetar a mis padres.} « Je dois respecter mes parents. » \\
\hline
\esp{tú} & \esp{debes} & \esp{Debes estudiar más.} « Tu dois travailler davantage. » \\
\hline
\esp{él/ella/usted} & \esp{debe} & \esp{Debe pedir ayuda.} « Il/elle doit demander de l'aide. » \\
\hline
\esp{nosotros} & \esp{debemos} & \esp{Debemos proteger el medio ambiente.} « Nous devons protéger l'environnement. » \\
\hline
\esp{vosotros} & \esp{debéis} & \esp{Debéis llegar a tiempo.} « Vous devez arriver à l'heure. » \\
\hline
\esp{ellos/ustedes} & \esp{deben} & \esp{Deben decir la verdad.} « Ils doivent dire la vérité. » \\
\hline
\end{tabular}
\end{center}

La négation se place devant \esp{deber} : \esp{No debes mentir} « Tu ne dois pas mentir », \esp{No debes juzgar a los demás} « Tu ne dois pas juger les autres ». Remarque la nuance : \esp{no debes salir de noche} signifie « tu ne devrais pas sortir la nuit » (conseil), alors que \esp{no salgas de noche} est une interdiction directe.

\begin{exemplebox}[Exemples traduits]
\begin{itemize}
\item \esp{Debes decir la verdad.} « Tu dois dire la vérité. »
\item \esp{Debes hablar con tus padres de tus problemas.} « Tu dois parler de tes problèmes avec tes parents. »
\item \esp{No debes beber alcohol si vas a conducir.} « Tu ne dois pas boire d'alcool si tu vas conduire. »
\item \esp{Debemos luchar contra el acoso escolar.} « Nous devons lutter contre le harcèlement scolaire. »
\end{itemize}
\end{exemplebox}

\courssub{hay que + infinitif : l'obligation générale}

Quand la règle vaut pour \textbf{tout le monde}, sans viser une personne en particulier, l'espagnol utilise \esp{hay que + infinitif}. C'est l'équivalent exact du français « il faut ».

\begin{exemplebox}[Exemples traduits]
\begin{itemize}
\item \esp{Hay que respetar a los demás.} « Il faut respecter les autres. »
\item \esp{Hay que luchar contra el acoso.} « Il faut lutter contre le harcèlement. »
\item \esp{Hay que proteger a los jóvenes de las drogas.} « Il faut protéger les jeunes de la drogue. »
\item \esp{No hay que juzgar sin conocer.} « Il ne faut pas juger sans connaître. »
\end{itemize}
\end{exemplebox}

\begin{attention}[\esp{hay que} est invariable et impersonnel]
\esp{Hay que} ne se conjugue \textbf{jamais} à la personne : on ne dit pas « \esp{has que} » ni « \esp{tienes que} » pour traduire « il faut ». La forme \esp{hay} est la forme impersonnelle du verbe \esp{haber}. Seul le \textbf{temps} peut changer : au passé, \esp{había que + infinitif} « il fallait » (\esp{Había que llegar temprano} « Il fallait arriver tôt ») ; au futur, \esp{habrá que + infinitif} « il faudra » (\esp{Habrá que sensibilizar a los alumnos} « Il faudra sensibiliser les élèves »). Ces deux formes sont un excellent complément à connaître pour le Bac. Enfin, ne confonds pas \esp{hay que} (il faut, général) avec \esp{tener que} (devoir, personnel) : \esp{Tengo que estudiar} « Je dois travailler ».
\end{attention}

\courssub{es importante que + subjonctif : le conseil insistant}

Pour insister sur un conseil en le formulant avec une subordonnée, l'espagnol emploie une expression impersonnelle suivie de \esp{que} et du \textbf{subjonctif} :

\begin{propriete}[Expression impersonnelle + que + subjonctif]
Après \esp{es importante que} « il est important que », \esp{es necesario que} « il est nécessaire que », \esp{es bueno que} « il est bon que », \esp{es mejor que} « il vaut mieux que » et \esp{te aconsejo que} « je te conseille que », le verbe de la subordonnée se met au \textbf{subjonctif présent}, car on exprime un souhait, un conseil ou une nécessité, pas un fait certain.
\end{propriete}

\begin{exemplebox}[Exemples traduits]
\begin{itemize}
\item \esp{Es importante que pidas ayuda.} « Il est important que tu demandes de l'aide. »
\item \esp{Es importante que hables con tus padres.} « Il est important que tu parles à tes parents. »
\item \esp{Es necesario que los jóvenes eviten las drogas.} « Il est nécessaire que les jeunes évitent la drogue. »
\item \esp{Es bueno que hagas deporte.} « Il est bon que tu fasses du sport. »
\item \esp{Te aconsejo que no salgas solo de noche.} « Je te conseille de ne pas sortir seul la nuit. »
\end{itemize}
\end{exemplebox}

\begin{attention}[N'oublie jamais le subjonctif après « es importante que »]
L'erreur classique du francophone est de mettre un indicatif par calque du français : « \esp{Es importante que pides ayuda} » est \textbf{faux}. En français, « il est important que tu demandes » porte déjà un subjonctif : garde-le en espagnol ! \esp{Es importante que \textbf{pidas} ayuda}. Autre piège : ne mets pas d'infinitif après \esp{que} ; l'infinitif n'est possible que sans \esp{que} et sans sujet propre : \esp{Es importante pedir ayuda} « Il est important de demander de l'aide » (général), mais \esp{Es importante que \textbf{tú} pidas ayuda} (personnel).
\end{attention}

\courssub{Tableau comparatif et organigramme de choix}

\begin{center}
\begin{tabularx}{\linewidth}{|X|X|X|}
\hline
\rowcolor{popblueL} Français & Español & Remarque \\
\hline
Il faut respecter les autres. & \esp{Hay que respetar a los demás.} & Impersonnel, invariable, règle générale. \\
\hline
Tu dois travailler. & \esp{Debes estudiar.} & Personnel, conjugué, s'adresse à quelqu'un. \\
\hline
Ne fume pas ! & \esp{¡No fumes!} & Interdiction directe : no + subjonctif. \\
\hline
Défense de fumer. & \esp{No fumar.} & Panneau impersonnel : no + infinitif. \\
\hline
Il est important que tu demandes de l'aide. & \esp{Es importante que pidas ayuda.} & Subordonnée : subjonctif obligatoire. \\
\hline
Je te conseille de parler à tes parents. & \esp{Te aconsejo que hables con tus padres.} & Conseil personnalisé : subjonctif. \\
\hline
\end{tabularx}
\end{center}

\begin{popfigure}
\begin{tikzpicture}[
  node distance=0.9cm and 1.6cm,
  question/.style={rectangle, rounded corners, draw=popdark, fill=popblueL, align=center, font=\small\bfseries, inner sep=6pt},
  choix/.style={rectangle, rounded corners, draw=poporange, fill=poporangeL, align=center, font=\small, inner sep=5pt},
  fleche/.style={-{Stealth[length=2.5mm]}, thick, draw=popdark},
  etiq/.style={font=\scriptsize\itshape, text=popmuted}
]
\node[question, align=center] (q){Quelle structure\\ pour conseiller ?};
\node[choix, below left=1.1cm and 2.2cm of q, align=center] (pers){Personne précise\\ \esp{Debes estudiar.}\\ \esp{deber + infinitif}};
\node[choix, below=1.1cm of q, align=center] (gen){Règle générale\\ \esp{Hay que respetar.}\\ \esp{hay que + infinitif}};
\node[choix, below right=1.1cm and 2.2cm of q, align=center] (sub){Avec subordonnée\\ \esp{Es importante que pidas ayuda.}\\ expression + \esp{que} + subjonctif};
\node[choix, below=2.6cm of q, fill=poppinkL, draw=poppink, align=center] (inter){Interdiction\\ \esp{¡No fumes!}\\ \esp{no} + subjonctif};
\draw[fleche] (q) -- node[etiq, above left, align=center]{devoir\\ personnel} (pers);
\draw[fleche] (q) -- node[etiq, right, align=center]{obligation\\ générale} (gen);
\draw[fleche] (q) -- node[etiq, above right, align=center]{conseil\\ insistant} (sub);
\draw[fleche] (q.south) .. controls +(0,-0.6) and +(0,0.8) .. node[etiq, left]{interdire} (inter.north);
\end{tikzpicture}
\legende{Organigramme de choix : quelle structure employer pour conseiller, obliger ou interdire ?}
\end{popfigure}

\begin{methode}[Choisir la bonne structure de conseil]
\begin{enumerate}
\item \textbf{Identifie ta cible.} Je m'adresse à une personne précise (un ami, un élève) → j'utilise \esp{deber + infinitif} conjugué : \esp{Debes descansar más}.
\item \textbf{Formule une règle générale.} Le conseil vaut pour tout le monde, sans destinataire précis → j'utilise \esp{hay que + infinitif} : \esp{Hay que respetar las normas}.
\item \textbf{Insiste avec une subordonnée.} Je veux marquer l'importance ou donner un conseil appuyé → j'utilise \esp{es importante que / te aconsejo que + subjonctif} : \esp{Te aconsejo que hables con un adulto}.
\item \textbf{Interdis clairement.} Je veux défendre une conduite à risque → \esp{no + subjonctif} pour une personne (\esp{¡No te drogues!}) ou \esp{no + infinitif} pour une affiche (\esp{No fumar}).
\item \textbf{Vérifie les pronoms.} À la forme négative, ils passent \emph{avant} le verbe : \esp{¡No lo hagas!} et jamais « \esp{no hágalo} » soudé.
\end{enumerate}
\end{methode}

\begin{exoresolu}[Transforme les slogans]
\textbf{Énoncé.} Transforme chaque conseil affirmatif en interdiction ou en conseil négatif, selon le modèle : \esp{¡Fuma menos!} → \esp{¡No fumes!}
\begin{enumerate}
\item \esp{¡Bebe alcohol en la fiesta!} (tú)
\item \esp{¡Haz trampa en el examen!} (tú)
\item \esp{¡Levántate tarde todos los días!} (tú, pronominal)
\item \esp{¡Dile mentiras a tu madre!} (tú, pronom \esp{le})
\item \esp{¡Tirad papeles al suelo!} (vosotros)
\end{enumerate}
\tcblower
\textbf{Solution détaillée.}
\begin{enumerate}
\item \esp{¡No bebas alcohol en la fiesta!} « Ne bois pas d'alcool à la fête ! » — \esp{beber} : subjonctif \esp{bebas} (voyelle \esp{-a} pour un verbe en \esp{-er}).
\item \esp{¡No hagas trampa en el examen!} « Ne triche pas à l'examen ! » — \esp{hacer} : première personne \esp{hago} → subjonctif \esp{hagas}.
\item \esp{¡No te levantes tarde todos los días!} « Ne te lève pas tard tous les jours ! » — verbe pronominal : le pronom \esp{te} se détache et passe \textbf{avant} le verbe.
\item \esp{¡No le digas mentiras a tu madre!} « Ne dis pas de mensonges à ta mère ! » — \esp{decir} : \esp{digo} → \esp{digas} ; le pronom \esp{le} passe avant le verbe.
\item \esp{¡No tiréis papeles al suelo!} « Ne jetez pas de papiers par terre ! » — à \esp{vosotros}, le subjonctif prend \esp{-éis} pour les verbes en \esp{-ar}.
\end{enumerate}
\end{exoresolu}

\begin{exercice}[À ton tour : conseille un camarade]
Ton ami Carlos veut abandonner los estudios, sortir chaque nuit et ne parle plus à ses parents. Rédige cinq phrases en espagnol en utilisant chacune des structures suivantes une fois : \esp{deber + infinitif}, \esp{hay que + infinitif}, \esp{es importante que + subjonctif}, \esp{te aconsejo que + subjonctif}, \esp{no + subjonctif}.
\end{exercice}

\begin{corrige}[Exercice « À ton tour : conseille un camarade »]
Proposition de corrigé (d'autres réponses sont possibles) :
\begin{enumerate}
\item \esp{Debes seguir estudiando, Carlos.} « Tu dois continuer à étudier, Carlos. »
\item \esp{Hay que pensar en el futuro.} « Il faut penser à l'avenir. »
\item \esp{Es importante que hables con tus padres.} « Il est important que tu parles avec tes parents. »
\item \esp{Te aconsejo que no salgas todas las noches.} « Je te conseille de ne pas sortir tous les soirs. »
\item \esp{¡No abandones los estudios!} « N'abandonne pas tes études ! »
\end{enumerate}
Vérifie toujours : le verbe après \esp{no} et après \esp{que} est-il bien au subjonctif ? Les pronoms sont-ils placés avant le verbe à la forme négative ?
\end{corrige}

\begin{aretenir}[L'essentiel de la section]
\begin{itemize}
\item Impératif négatif = \cle{no + subjonctif présent} : \esp{no fumes}, \esp{no hagas}, \esp{no bebáis}, \esp{no fume} (usted), \esp{no fumemos} (nosotros).
\item Les pronoms passent \textbf{avant} le verbe à la forme négative : \esp{¡No lo hagas!} contre \esp{¡Hazlo!}.
\item \esp{deber + infinitif} : devoir \textbf{personnel}, conjugué (\esp{Debes decir la verdad}).
\item \esp{hay que + infinitif} : obligation \textbf{générale}, invariable (\esp{Hay que respetar}) ; passé \esp{había que}, futur \esp{habrá que}.
\item \esp{es importante que / es necesario que / te aconsejo que + subjonctif} : conseil insistant avec subordonnée (\esp{Es importante que pidas ayuda}).
\item Panneau impersonnel : \esp{no + infinitif} (\esp{No fumar} = « Défense de fumer »).
\end{itemize}
\end{aretenir}

\coursec{Méthode du commentaire de texte et commentaire modèle}

Dans la section précédente, tu as appris à repérer et à employer l'impératif négatif (\esp{no fumes}, \esp{no bebas}) ainsi que les structures du conseil et de l'obligation (\esp{deber + infinitivo}, \esp{hay que}, \esp{es importante que + subjuntivo}). Ces outils grammaticaux ne servent pas seulement à parler : ils sont aussi des \textbf{indices de lecture}. Quand tu commentes un texte de sensibilisation comme \esp{Jóvenes responsables, sociedad fuerte}, reconnaître ces procédés et expliquer leur effet sur le lecteur, c'est exactement ce qu'attend l'examinateur. Cette section t'apprend donc à transformer tes connaissances de langue en \cle{commentaire de texte} organisé, cité et analysé.

\courssub{Le texte support : rappel et relecture méthodique}

Avant de commenter, relisons le texte de la leçon avec un regard d'analyste : souligne mentalement le lexique négatif, les impératifs et les structures de conseil.

\begin{textebox}[Jóvenes responsables, sociedad fuerte (texte original)]
En muchas ciudades, los jóvenes se enfrentan cada día a tentaciones peligrosas: el alcohol en las fiestas, las drogas en la calle, la violencia en las redes sociales. Estas conductas de riesgo destruyen la salud, rompen las familias y debilitan a toda la sociedad.

Sin embargo, cada joven puede elegir otro camino. ¡Di no a las drogas! ¡Respeta a tus profesores y a tus padres! No tires la basura en la calle y cuida tu barrio. Estudia con esfuerzo, porque la educación es la llave del futuro.

Además, debemos ayudar a los demás: hay que proteger a los más pequeños y es importante que todos participemos en la vida de nuestra comunidad. Un joven responsable no piensa solo en sí mismo: piensa en su familia, en su escuela y en su país.

Queridos jóvenes, el futuro está en vuestras manos. Sed valientes, sed solidarios, sed responsables. Si cada uno de vosotros adopta comportamientos responsables, construiremos juntos una sociedad más fuerte, más justa y más humana. La responsabilidad de hoy es la libertad de mañana.
\end{textebox}

\textbf{Glossaire :} \esp{se enfrentan a} : sont confrontés à ; \esp{las redes sociales} : les réseaux sociaux ; \esp{rompen} : brisent ; \esp{tirar la basura} (forme du texte : \esp{no tires}) : jeter les ordures ; \esp{el barrio} : le quartier ; \esp{la llave} : la clé ; \esp{los demás} : les autres ; \esp{sed} : soyez (impératif de \esp{ser}, 2e pers. plur.) ; \esp{solidarios} : solidaires ; \esp{justa} : juste, équitable.

\textbf{Questions de compréhension :} 1) ¿Qué conductas de riesgo menciona el texto? 2) ¿Qué consejos da el autor a los jóvenes? 3) ¿Cuál es la visión de sociedad que propone el texto?

\courssub{La méthode du commentaire : le plan en entonnoir}

Le commentaire de texte au lycée suit une logique d'\textbf{entonnoir} : on part du texte (introduction), on l'analyse (développement), et on s'élargit vers une réflexion personnelle (conclusion).

\begin{popfigure}
\begin{tikzpicture}[
  etape/.style={rectangle, rounded corners=4pt, align=center, text width=6.2cm, minimum height=1.05cm, font=\small},
  fleche/.style={-{Stealth[length=3mm]}, thick, popmuted}
]
\node[etape, fill=popblueL, draw=popblue, align=center] (intro){\textbf{Introduction}\\ présenter le texte : type, thème, source};
\node[etape, fill=poporangeL, draw=poporange, below=0.75cm of intro, align=center] (dev){\textbf{Analyse (2-3 axes)}\\ constat / procédés / message};
\node[etape, fill=popgreenL, draw=popgreen, below=0.75cm of dev, align=center] (concl){\textbf{Conclusion}\\ bilan + ouverture personnelle};
\draw[fleche] (intro) -- (dev);
\draw[fleche] (dev) -- (concl);
\node[right=0.25cm of intro, font=\footnotesize\itshape, text=popmuted] {du texte...};
\node[right=0.25cm of concl, font=\footnotesize\itshape, text=popmuted] {...vers ta réflexion};
\end{tikzpicture}
\legende{Le plan en entonnoir du commentaire de texte : du texte vers ta réflexion.}
\end{popfigure}

\begin{methode}[Les trois étapes du commentaire]
\begin{enumerate}
\item \textbf{Introduction (3-4 lignes)} : présente le texte en trois informations — le \textbf{type} de texte (article, discours, texte de sensibilisation), le \textbf{thème} (ici : les conduites à risque et la responsabilité des jeunes) et la \textbf{source ou la situation} (une campagne de prévention, une revue scolaire...). Termine par l'annonce du plan.
\item \textbf{Développement (2-3 axes)} : chaque axe = un paragraphe = une idée + des citations analysées. Ne mélange jamais deux idées dans le même paragraphe.
\item \textbf{Conclusion (2-3 lignes)} : fais le bilan (que nous apprend le texte ?) puis ouvre par ton avis personnel, en réutilisant les structures de la leçon : \esp{En mi opinión, hay que educar a los jóvenes...} (« À mon avis, il faut éduquer les jeunes... »).
\end{enumerate}
\end{methode}

\begin{aretenir}[Le vocabulaire espagnol du commentaire]
\begin{center}
\begin{tabularx}{\linewidth}{|X|X|}
\hline
\rowcolor{popblueL} \textbf{Español} & \textbf{Français} \\
\hline
El texto trata de... & Le texte traite de... \\
\hline
El autor denuncia / aconseja / sensibiliza & L'auteur dénonce / conseille / sensibilise \\
\hline
En el primer párrafo... & Dans le premier paragraphe... \\
\hline
El autor utiliza el imperativo para... & L'auteur utilise l'impératif pour... \\
\hline
Esta cita muestra que... & Cette citation montre que... \\
\hline
El objetivo del texto es... & L'objectif du texte est... \\
\hline
En conclusión / En mi opinión & En conclusion / À mon avis \\
\hline
\end{tabularx}
\end{center}
\end{aretenir}

\courssub{Les trois axes d'analyse appliqués au texte}

\textbf{Axe 1 : le constat — les conduites à risque dénoncées.} Relève le \textbf{lexique négatif} et les exemples concrets : \esp{el alcohol}, \esp{las drogas}, \esp{la violencia en las redes sociales}. Observe aussi les verbes forts : \esp{destruyen la salud, rompen las familias, debilitan a toda la sociedad} — une gradation (santé → famille → société) qui montre que le danger dépasse l'individu.

\textbf{Axe 2 : les procédés de sensibilisation.} C'est ici que ta grammaire devient un outil d'analyse :
\begin{itemize}
\item les \textbf{impératifs affirmatifs} : \esp{¡Di no a las drogas!}, \esp{¡Respeta a tus profesores!} — le tutoiement (\esp{di}, \esp{respeta}) crée un dialogue direct avec le jeune lecteur ;
\item l'\textbf{impératif négatif} : \esp{No tires la basura} — l'interdiction marque la limite à ne pas franchir ;
\item les \textbf{structures de conseil et d'obligation} : \esp{debemos ayudar}, \esp{hay que proteger}, \esp{es importante que todos participemos} — elles transforment le conseil en devoir collectif ;
\item l'\textbf{opposition risque / responsabilité} : le texte est construit en deux blocs (danger au §1, solutions aux §2-4), articulés par \esp{Sin embargo} (« cependant »).
\end{itemize}

\textbf{Axe 3 : le message et la visée.} Le texte veut \textbf{responsabiliser} : \esp{el futuro está en vuestras manos} (« l'avenir est entre vos mains »). La formule finale \esp{La responsabilidad de hoy es la libertad de mañana} fonctionne comme un slogan : responsabiliser les jeunes, c'est construire une société \esp{más fuerte, más justa y más humana}.

\begin{methode}[Comment citer dans un commentaire : la règle CTA]
Pour chaque idée, applique la règle \textbf{C-T-A} : \textbf{C}iter — \textbf{T}raduire — \textbf{A}nalyser.
\begin{enumerate}
\item \textbf{Citer} en espagnol, entre guillemets : « \esp{¡Di no a las drogas!} »
\item \textbf{Traduire} entre parenthèses : (« Dis non à la drogue ! »)
\item \textbf{Analyser} le procédé et son effet : l'impératif à la deuxième personne interpelle directement le jeune lecteur et lui donne un ordre simple et mémorable, comme un slogan de campagne.
\end{enumerate}
Une citation sans analyse ne rapporte aucun point ; une analyse sans citation est du paraphrase.
\end{methode}

\begin{attention}[Les deux erreurs fatales au commentaire]
\begin{itemize}
\item \textbf{La paraphrase} : raconter le texte avec tes propres mots (« el autor dice que las drogas son malas... ») n'est pas un commentaire. Le correcteur veut savoir \emph{comment} le texte agit sur le lecteur, pas \emph{ce qu'il raconte}.
\item \textbf{La citation abandonnée} : chaque citation doit être suivie d'une analyse (procédé + effet). Interdit : aligner trois citations sans commentaire.
\item Piège de langue : « commenter » ne se dit pas \esp{comentar} au sens scolaire sans complément ; dis \esp{analizar el texto}, \esp{hacer un comentario de texto}.
\end{itemize}
\end{attention}

\courssub{Commentaire modèle rédigé}

\begin{exemplebox}[Commentaire modèle (environ 10 lignes)]
\esp{Este texto de sensibilización, titulado «Jóvenes responsables, sociedad fuerte», trata de las conductas de riesgo de los jóvenes y de la importancia de adoptar comportamientos responsables. El autor empieza con un diagnóstico negativo: denuncia «el alcohol en las fiestas, las drogas en la calle, la violencia en las redes sociales» (« l'alcool dans les fêtes, la drogue dans la rue, la violence sur les réseaux sociaux »), y utiliza verbos fuertes como «destruyen» y «rompen» para mostrar la gravedad del problema. Después, el autor utiliza el imperativo para aconsejar directamente al lector: «¡Di no a las drogas!» (« Dis non à la drogue ! »). Este procedimiento crea un diálogo directo con el joven. También emplea estructuras de obligación como «hay que proteger» (« il faut protéger ») para transformar el consejo en deber colectivo. Finalmente, el texto transmite un mensaje de esperanza: «el futuro está en vuestras manos» (« l'avenir est entre vos mains »). En conclusión, es un texto eficaz porque mezcla denuncia y consejos. En mi opinión, hay que educar a los jóvenes desde la escuela para construir una sociedad más responsable.}
\end{exemplebox}

\textbf{Traduction du commentaire modèle :} « Ce texte de sensibilisation, intitulé \emph{Des jeunes responsables, une société forte}, traite des conduites à risque des jeunes et de l'importance d'adopter des comportements responsables. L'auteur commence par un diagnostic négatif : il dénonce "l'alcool dans les fêtes, la drogue dans la rue, la violence sur les réseaux sociaux", et utilise des verbes forts comme "détruisent" et "brisent" pour montrer la gravité du problème. Ensuite, l'auteur utilise l'impératif pour conseiller directement le lecteur : "Dis non à la drogue !". Ce procédé crée un dialogue direct avec le jeune. Il emploie aussi des structures d'obligation comme "il faut protéger" pour transformer le conseil en devoir collectif. Enfin, le texte transmet un message d'espérance : "l'avenir est entre vos mains". En conclusion, c'est un texte efficace parce qu'il mêle dénonciation et conseils. À mon avis, il faut éduquer les jeunes dès l'école pour construire une société plus responsable. »

Observe la phrase-clé du commentaire, à réutiliser dans tes copies : \esp{El autor utiliza el imperativo para aconsejar directamente al lector} = « L'auteur utilise l'impératif pour conseiller directement le lecteur ». Autres modèles utiles : \esp{El autor denuncia... para sensibilizar al lector} (« L'auteur dénonce... pour sensibiliser le lecteur ») ; \esp{Esta oposición entre riesgo y responsabilidad estructura todo el texto} (« Cette opposition entre risque et responsabilité structure tout le texte »).

\begin{exoresolu}[Analyser une citation avec la règle CTA]
Énoncé : analyse la citation suivante extraite du texte en appliquant la règle CTA (citer, traduire, analyser) : \esp{Sed valientes, sed solidarios, sed responsables.}
\tcblower
Solution détaillée :
\begin{enumerate}
\item \textbf{Citation} : « \esp{Sed valientes, sed solidarios, sed responsables} ».
\item \textbf{Traduction} : « Soyez courageux, soyez solidaires, soyez responsables ».
\item \textbf{Analyse} : il s'agit de l'impératif de \esp{ser} à la deuxième personne du pluriel (\esp{sed}), forme irrégulière. La répétition trois fois du même verbe crée une \textbf{anaphore} qui donne un rythme de slogan, facile à mémoriser. La gradation des adjectifs (du courage individuel à la responsabilité collective) élargit progressivement le message : le lecteur passe de sa personne à la société entière. Le passage du \esp{tú} au \esp{vosotros} montre que la responsabilité est un projet commun.
\end{enumerate}
\end{exoresolu}

\begin{savaistu}[Le commentaire au Bac : trois critères de notation]
Au baccalauréat camerounais, le commentaire de texte est noté sur trois critères : la \textbf{compréhension} du texte, la \textbf{pertinence des citations} (courtes, bien choisies, traduites) et la \textbf{qualité de ton espagnol}. Astuce de stratège : dans ta conclusion personnelle, réutilise les structures de la leçon — \esp{hay que}, \esp{deber + infinitivo}, \esp{es importante que + subjuntivo} — tu montres ainsi au correcteur que tu maîtrises la grammaire du programme. Exemple : \esp{En mi opinión, es importante que los jóvenes digan no a las drogas y hay que ayudarlos con campañas de prevención.}
\end{savaistu}

\courssub{Entraînement guidé}

\begin{exercice}[Rédiger l'introduction et un axe d'un commentaire]
Voici un extrait d'un texte similaire : \esp{« En nuestra escuela, algunos alumnos fuman detrás de las aulas y otros abandonan los estudios. ¡No copies en los exámenes! ¡Ayuda a tus compañeros! Debemos construir una escuela donde todos respeten las normas. »}
\begin{enumerate}
\item Rédige l'\textbf{introduction} d'un commentaire (3-4 lignes) : type de texte, thème, annonce du plan.
\item Rédige \textbf{un axe} (5-6 lignes) sur les procédés de sensibilisation, avec au moins deux citations traduites et analysées (règle CTA).
\item Termine par une phrase de conclusion utilisant \esp{hay que} ou \esp{es importante que + subjuntivo}.
\end{enumerate}
\end{exercice}

\begin{corrige}[Exercice : éléments de réponse]
\textbf{Introduction possible :} \esp{Este texto de sensibilización trata de las malas conductas de algunos alumnos y de la necesidad de respetar las normas en la escuela. Vamos a analizar primero el diagnóstico negativo, después los recursos de sensibilización y finalmente el mensaje del autor.}

\textbf{Axe 2 possible :} \esp{El autor utiliza el imperativo negativo para prohibir: «¡No copies en los exámenes!» (« Ne triche pas aux examens ! »). Este imperativo marca un límite claro para el alumno. También emplea el imperativo afirmativo «¡Ayuda a tus compañeros!» (« Aide tes camarades ! ») para proponer un comportamiento positivo. Finalmente, la estructura «debemos construir» (« nous devons construire ») transforma el consejo individual en un proyecto colectivo de toda la escuela.}

\textbf{Conclusion possible :} \esp{En conclusión, el texto mezcla prohibiciones y consejos positivos. En mi opinión, es importante que todos los alumnos respeten las normas y hay que ayudar a los que abandonan los estudios.}
\end{corrige}

\begin{aretenir}[L'essentiel de la section]
\begin{itemize}
\item Un commentaire = introduction (type, thème, source) + 2-3 axes + conclusion (bilan + ouverture).
\item Chaque citation suit la règle \textbf{CTA} : citer en espagnol, traduire entre parenthèses, analyser le procédé et son effet.
\item Jamais de paraphrase ; jamais de citation sans analyse.
\item Réutilise la grammaire de la leçon (\esp{hay que}, \esp{deber}, subjonctif) dans ta conclusion pour gagner des points de langue.
\end{itemize}
\end{aretenir}

\coursec{Texte support : « Jóvenes responsables, sociedad fuerte »}

\begin{textebox}[Texte original — Campagne de prévention du Ministère de l'Éducation (adapté)]
En los institutos de Yaoundé, Douala o Malabo, los jóvenes se enfrentan cada día a \cle{conductas de riesgo} : el consumo de alcohol y tabaco, la violencia en los recreos, la adicción al móvil y la trampa en los exámenes. Estas conductas ponen en peligro su \cle{salud}, su \cle{futuro} y la \cle{convivencia} en el centro.

Ante esta realidad, la comunidad educativa lanza un mensaje claro: \textbf{la prevención empieza por uno mismo}. \esp{¡No esperes a que otros decidan por ti!} \esp{Respeta las normas,} \esp{dialoga} con tus profesores y \esp{busca ayuda} si la necesitas. \esp{Recuerda:} un joven responsable construye una sociedad fuerte.

\esp{Hay que} decir \esp{no} a la presión del grupo. \esp{Debes} pensar en las consecuencias de tus actos. \esp{Es importante que} valores tu esfuerzo y el de los demás. Solo así transformarás los retos en oportunidades.
\end{textebox}

\textbf{Traduction et glossaire} : Dans les lycées de Yaoundé, Douala ou Malabo, les jeunes font face chaque jour à des \textbf{conduites à risque} : consommation d'alcool et de tabac, violence pendant les récréations, addiction au portable et tricherie aux examens. Ces conduites mettent en danger leur \textbf{santé}, leur \textbf{avenir} et la \textbf{vie ensemble} dans l'établissement. Face à cette réalité, la communauté éducative lance un message clair : \textbf{la prévention commence par soi-même}. « N'attends pas que les autres décident pour toi ! » « Respecte les règles, » « dialogue » avec tes professeurs et « cherche de l'aide » si tu en as besoin. « Souviens-toi : » un jeune responsable construit une société forte. « Il faut » dire « non » à la pression du groupe. « Tu dois » penser aux conséquences de tes actes. « Il est important que » tu valeurs ton effort et celui des autres. Ce n'est qu'ainsi que tu transformeras les défis en opportunités.

\begin{definition}[Mots-clés du texte]
\begin{itemize}
\item \esp{conducta de riesgo} (nf) : conduite à risque / comportement à risque.
\item \esp{convivencia} (nf) : vie ensemble, cohabitation harmonieuse (plus large que « coexistence »).
\item \esp{prevención} (nf) : prévention. \esp{prevenir} = prévenir, éviter.
\item \esp{presión del grupo} (nf) : pression du groupe (pair pressure).
\item \esp{valorar} (v) : valoriser, apprécier à sa juste valeur.
\item \esp{reto} (nm) : défi, challenge.
\end{itemize}
\end{definition}

\coursec{Lexique thématique : les conduites à risque et les comportements responsables}

\begin{center}
\begin{tabularx}{\linewidth}{|X|X|X|}
\hline
\rowcolor{popblueL} \textbf{Français} & \textbf{Español} & \textbf{Remarque / Contexte} \\
\hline
la conduite à risque & \esp{la conducta de riesgo} & \emph{Locution nominale fixe.} \\
la violence & \esp{la violencia} & \esp{violento/a} (adj). \\
le harcèlement (scolaire) & \esp{el acoso (escolar)} & \esp{acosar} (v). \emph{Faux ami :} \esp{acoso} $\neq$ « acajou ». \\
la tricherie & \esp{el fraude} / \esp{la trampa} & \esp{copiar en un examen} (tricher à un examen). \\
l'addiction, la dépendance & \esp{la adicción} & \esp{adicto/a} (adj/n). \esp{ser adicto a} (+ nom). \\
l'alcool & \esp{el alcohol} & \esp{beber alcohol}. \emph{Attention :} pas d'article contracté \emph{del} ici (alcohol commence par a tonique $\rightarrow$ \emph{el alcohol}, mais \emph{beber alcohol}). \\
le tabac / fumer & \esp{el tabaco} / \esp{fumar} & \esp{fumador/a} (fumeur). \\
le portable / les réseaux sociaux & \esp{el móvil} / \esp{las redes sociales} & \esp{estar enganchado al móvil} (être accro). \\
la santé & \esp{la salud} & \esp{cuida tu salud} (prends soin de ta santé). \\
le respect / respecter & \esp{el respeto} / \esp{respetar} & \esp{respetar las normas / a los demás}. \\
la responsabilité / responsable & \esp{la responsabilidad} / \esp{responsable} & adjectif invariable en genre : \esp{un chico responsable}, \esp{una chica responsable}. \\
le dialogue / dialoguer & \esp{el diálogo} / \esp{dialogar} & \emph{Verbe en -ar.} \\
l'effort & \esp{el esfuerzo} & \esp{valer el esfuerzo} (valoir l'effort). \\
les normes, les règles & \esp{las normas} & \esp{cumplir / respetar las normas}. \\
conseiller & \esp{aconsejar} / \esp{consejar} & \esp{te aconsejo que} + subjonctif. \\
interdire & \esp{prohibir} & \esp{Está prohibido} + infinitif (impersonnel). \\
\hline
\end{tabularx}
\end{center}

\begin{savaistu}[Cameroun et Guinée équatoriale : contextes partagés]
La Guinée équatoriale (\esp{Guinea Ecuatorial}) est le seul pays africain où l'espagnol est langue officielle. Les lycées de Malabo ou Bata connaissent les mêmes enjeux que ceux de Yaoundé ou Douala : prévention de la violence, éducation à la santé, réussite au baccalauréat. Les campagnes du \esp{MINEDUC} (Ministère de l'Éducation) y utilisent souvent des slogans bilingues (espagnol/français) ou en \esp{fang} / \esp{bubi}. Exemple réel : \esp{« Tu futuro está en tus manos, ¡cuídalo! »} (« Ton avenir est entre tes mains, prends-en soin ! »).
\end{savaistu}

\coursec{Grammaire 1 : L'impératif affirmatif (tú, vosotros, usted)}

\begin{propriete}[Formation de l'impératif affirmatif]
L'impératif affirmatif n'a que trois personnes : \textbf{tú}, \textbf{vosotros} (Espagne), \textbf{usted} (formel, Amérique + contexte formel partout). Il n'existe pas de 1\up{re} personne du singulier ni de 3\up{e} personne (on utilise le subjonctif pour \emph{él/ella/ellos}).

\begin{center}
\begin{tabular}{|l|c|c|c|}
\hline
\rowcolor{popblueL} \textbf{Verbe} & \textbf{tú} & \textbf{vosotros} & \textbf{usted} \\
\hline
\esp{hablar} (parler) & \esp{habla} & \esp{hablad} & \esp{hable} \\
\hline
\esp{comer} (manger) & \esp{come} & \esp{comed} & \esp{coma} \\
\hline
\esp{vivir} (vivre) & \esp{vive} & \esp{vivid} & \esp{viva} \\
\hline
\esp{respetar} & \esp{respeta} & \esp{respetad} & \esp{respete} \\
\hline
\esp{dialogar} & \esp{dialoga} & \esp{dialogad} & \esp{dialogue} \\
\hline
\end{tabular}
\end{center}

\textbf{Règle générale :}
\begin{itemize}
\item \textbf{tú} = 3\up{e} personne du singulier de l'indicatif présent (\esp{habla, come, vive}).
\item \textbf{vosotros} = radical + \esp{-ad / -ed / -id} (accent sur la voyelle thématique).
\item \textbf{usted} = 1\up{re} personne du singulier du \textbf{subjonctif présent} (\esp{hable, coma, viva}).
\end{itemize}
\end{propriete}

\begin{attention}[Impératifs affirmatifs irréguliers (tú) — À apprendre par cœur]
Ces 8 verbes sont \textbf{très fréquents} dans les conseils et consignes. Le francophone a tendance à utiliser la 3\up{e} personne de l'indicatif (ex. \esp{*hace}, \esp{*pon}) : \textbf{erreur grave}.

\begin{center}
\begin{tabular}{|l|l|l|l|}
\hline
\rowcolor{poporangeL} \textbf{Infinitif} & \textbf{tú} & \textbf{vosotros} & \textbf{usted} \\
\hline
\esp{ser} (être) & \esp{sé} & \esp{sed} & \esp{sea} \\
\hline
\esp{ir} (aller) & \esp{ve} & \esp{id} & \esp{vaya} \\
\hline
\esp{haber} (il y a / avoir) & \esp{—} & \esp{—} & \esp{haya} \\
\hline
\esp{hacer} (faire) & \esp{haz} & \esp{haced} & \esp{haga} \\
\hline
\esp{poner} (mettre) & \esp{pon} & \esp{poned} & \esp{ponga} \\
\hline
\esp{tener} (avoir) & \esp{ten} & \esp{tened} & \esp{tenga} \\
\hline
\esp{decir} (dire) & \esp{di} & \esp{decid} & \esp{diga} \\
\hline
\esp{salir} (sortir) & \esp{sal} & \esp{salid} & \esp{salga} \\
\hline
\esp{venir} (venir) & \esp{ven} & \esp{venid} & \esp{venga} \\
\hline
\end{tabular}
\end{center}

\emph{Mnémotechnique :} \textbf{« VIN DIESEL HAS TEN PONS »} (\esp{Ven, Di, Sal, Haz, Ten, Pon} + \esp{Ve, Sé} pour \esp{Ir, Ser}). \esp{Vosotros} est toujours régulier (\esp{-d} final).
\end{attention}

\begin{exemplebox}[Exemples d'impératifs affirmatifs en contexte]
\begin{itemize}
\item \esp{¡Respeta a tus profesores!} (tú) « Respecte tes professeurs ! »
\item \esp{¡Sed responsables, chicos!} (vosotros) « Soyez responsables, les gars ! »
\item \esp{¡Hable más despacio, por favor!} (usted) « Parlez plus lentement, s'il vous plaît ! »
\item \esp{¡Haz la tarea ahora!} (tú, irrégulier) « Fais les devoirs maintenant ! »
\item \esp{¡No olvides: ten paciencia!} « N'oublie pas : aie de la patience ! »
\end{itemize}
\end{exemplebox}

\coursec{Grammaire 2 : L'impératif négatif et les structures du conseil et de l'obligation}

\courssub{A. L'impératif négatif (tú, vosotros, usted)}

\begin{propriete}[Formation : \esp{NO} + SUBJONCTIF PRÉSENT]
Pour \textbf{toutes les personnes} (tú, vosotros, usted, ustedes), l'impératif négatif se forme avec \esp{no} + \textbf{subjonctif présent}. Les irrégularités de l'affirmatif (tú) disparaissent : on applique la règle générale du subjonctif.

\begin{center}
\begin{tabular}{|l|c|c|c|}
\hline
\rowcolor{popgreenL} \textbf{Verbe} & \textbf{no + tú} & \textbf{no + vosotros} & \textbf{no + usted} \\
\hline
\esp{hablar} & \esp{no hables} & \esp{no habléis} & \esp{no hable} \\
\hline
\esp{comer} & \esp{no comas} & \esp{no comáis} & \esp{no coma} \\
\hline
\esp{vivir} & \esp{no vivas} & \esp{no viváis} & \esp{no viva} \\
\hline
\esp{hacer} & \esp{no hagas} & \esp{no hagáis} & \esp{no haga} \\
\hline
\esp{ir} & \esp{no vayas} & \esp{no vayáis} & \esp{no vaya} \\
\hline
\esp{poner} & \esp{no pongas} & \esp{no pongáis} & \esp{no ponga} \\
\hline
\end{tabular}
\end{center}

\textbf{Rappel subjonctif présent (réguliers) :} radical + voyelle opposée (\emph{-ar} $\rightarrow$ \emph{-e, -es, -e, -emos, -éis, -en} ; \emph{-er/-ir} $\rightarrow$ \emph{-a, -as, -a, -amos, -áis, -an}).
\end{propriete}

\begin{attention}[Piège majeur : \esp{No fumas} vs \esp{No fumes}]
\begin{itemize}
\item \esp{No fumas} = Indicatif présent. Signifie « Tu ne fumes pas » (constat). \textbf{Pas un ordre.}
\item \esp{No fumes} = Subjonctif présent = Impératif négatif. Signifie « Ne fume pas ! » (ordre/interdiction).
\end{itemize}
Le francophone calque « Ne fume pas » $\rightarrow$ \esp{*No fumas}. \textbf{Interdit.} Toujours \esp{No + subjonctif}.
\end{attention}

\courssub{B. Structures du conseil, de l'obligation et de la nécessité}

\begin{propriete}[Tableau récapitulatif des nuances]
\begin{center}
\begin{tabularx}{\linewidth}{|X|X|X|X|}
\hline
\rowcolor{poppurpleL} \textbf{Structure} & \textbf{Valeur} & \textbf{Registre / Nuance} & \textbf{Exemple traduit} \\
\hline
\esp{Deber + infinitif} & Devoir moral, conseil fort, obligation personnelle & Subjectif : « Je pense que tu dois » & \esp{Debes estudiar más.} (Tu dois étudier plus.) \\
\hline
\esp{Tener que + infinitif} & Nécessité, contrainte externe, obligation factuelle & Plus factuel : « Tu es obligé de » & \esp{Tienes que entregar el trabajo mañana.} (Tu dois rendre le travail demain.) \\
\hline
\esp{Hay que + infinitif} & Obligation générale, impersonnelle, règle sociale & « Il faut », « On doit » (valeur générale) & \esp{Hay que respetar las normas.} (Il faut respecter les règles.) \\
\hline
\esp{Es importante que} + \textbf{subjonctif} & Conseil insistant, valeur éducative, avis de l'énonciateur & Subjectif + évaluatif. \textbf{Subjonctif obligatoire.} & \esp{Es importante que duermas bien.} (Il est important que tu dormes bien.) \\
\hline
\esp{Es necesario que} + \textbf{subjonctif} & Nécessité perçue comme objective & Subjonctif obligatoire. & \esp{Es necesario que vengas.} (Il faut que tu viennes.) \\
\hline
\esp{Recomiendo que} / \esp{Aconsejo que} + \textbf{subjonctif} & Conseil direct, avis personnel & Verbe de volonté $\rightarrow$ subjonctif. & \esp{Te aconsejo que lo pienses.} (Je te conseille d'y réfléchir.) \\
\hline
\end{tabularx}
\end{center}
\end{propriete}

\begin{attention}[Faux amis et pièges grammaticaux fréquents]
\begin{enumerate}
\item \esp{Deber} vs \esp{Deber de} : \esp{Debes estudiar} (Tu dois étudier \textbf{obligation/conseil}). \esp{Debe de estar cansado} (Il \textbf{doit être} fatigué \textbf{probabilité}). \textbf{Jamais} \esp{debes de estudiar} pour un conseil.
\item \esp{Hay que} est \textbf{impersonnel} (invariable). Jamais \esp{*Hay que estudias}, \esp{*He que estudiar}. Toujours \esp{Hay que + infinitif}.
\item \esp{Es importante que} appelle \textbf{le subjonctif}, pas l'indicatif. \esp{*Es importante que estudias} $\rightarrow$ \esp{Es importante que estudies}.
\item \esp{Prohibido} (participe/adjectif) s'accorde : \esp{Está prohibido fumar} / \esp{Están prohibidas las armas}. Ne pas confondre avec l'impératif \esp{Prohibid} (vosotros).
\end{enumerate}
\end{attention}

\begin{exemplebox}[Exemples variés pour la prévention]
\begin{itemize}
\item \esp{No bebas alcohol si vas a conducir.} (Impératif négatif \esp{beber} $\rightarrow$ \esp{bebas})
\item \esp{Debes decir la verdad siempre.} (Conseil moral fort)
\item \esp{Hay que denunciar el acoso escolar.} (Règle générale, impersonnel)
\item \esp{Es importante que pidas ayuda si te sientes mal.} (Subjonctif \esp{pidas} < \esp{pedir})
\item \esp{Tienes que apagar el móvil en clase.} (Contrainte scolaire concrète)
\item \esp{Te recomiendo que hables con el orientador.} (Conseil personnel, subjonctif \esp{hables})
\end{itemize}
\end{exemplebox}

\coursec{Méthode du commentaire de texte et commentaire modèle}

\begin{methode}[Commenter un texte de sensibilisation (4 étapes)]
\begin{enumerate}
\item \textbf{Identification et contextualisation} : Nature du document (article, affiche, discours, extrait de manuel), source, date, auteur, public visé, thème général (conduites à risque, prévention).
\item \textbf{Analyse du contenu (le \emph{quoi})} : Dégager la thèse principale (ex. « La responsabilité individuelle est la clé d'une société forte »). Repérer les arguments (constats, dangers, solutions). Noter le ton (alarmiste, bienveillant, autoritaire, pédagogique).
\item \textbf{Analyse de la forme (le \emph{comment})} : Relever les \textbf{procédés linguistiques} au service de la persuasion :
\begin{itemize}
\item \textbf{Impératifs} (affirmatifs/négatifs) : ordre direct, interpellation.
\item \textbf{Structures modales} (\esp{deber, hay que, es importante que}) : gradation dans l'obligation.
\item \textbf{Connecteurs logiques} : \esp{ante esta realidad, solo así, por tanto, en conclusión}.
\item \textbf{Champ lexical} : risque (\esp{peligro, riesgo, daño}) vs responsabilité (\esp{responsable, esfuerzo, futuro, convivencia}).
\item \textbf{Figures de style} : anaphore (\esp{¡No esperes! ¡Respeta! ¡Dialoga!}), antithèse (risque / avenir), rimes (slogans).
\end{itemize}
\item \textbf{Synthèse et évaluation critique} : L'auteur atteint-il son but ? Le texte est-il adapté au public jeune (registre \esp{tú}, exemples concrets) ? Quelle portée pour le Cameroun / la Guinée équatoriale ?
\end{enumerate}
\end{methode}

\begin{exoresolu}[Commentaire modèle du texte « Jóvenes responsables, sociedad fuerte »]
Énoncé. Commentez le texte support de la leçon en suivant la méthode.
\tcblower
\textbf{1. Identification.} Ce texte est un extrait de campagne de sensibilisation (type affiche ou discours d'ouverture d'année scolaire), émanant du Ministère de l'Éducation (\esp{MINEDUC}), destiné aux élèves de lycées hispanophones (Cameroun, Guinée équatoriale, Espagne). Thème : la prévention des conduites à risque chez les adolescents.

\textbf{2. Contenu.} Le texte dresse un constat alarmant (énumération : alcool, tabac, violence, mobile, triche) $\rightarrow$ thèse : « La prévention commence par soi-même » $\rightarrow$ solutions concrètes (respect, dialogue, aide) $\rightarrow$ conclusion optimiste : « Un jeune responsable construit une société forte ». Ton pédagogique, bienveillant mais ferme.

\textbf{3. Forme / Procédés.}
\begin{itemize}
\item \textbf{Accroche} : ancrage géographique (\esp{Yaoundé, Douala, Malabo}) $\rightarrow$ identification immédiate du lecteur camerounais/équatoguinéen.
\item \textbf{Champ lexical du risque} : \esp{conductas de riesgo, ponen en peligro, salud, futuro, presión del grupo}.
\item \textbf{Champ lexical de la solution} : \esp{prevención, respetar, dialogar, buscar ayuda, responsable, sociedad fuerte, esfuerzo, oportunidades}.
\item \textbf{Impératifs affirmatifs (tú)} : \esp{¡No esperes!, Respeta, dialoga, busca, Recuerda.} Interpellation directe, urgence.
\item \textbf{Structures modales variées} : \esp{Hay que decir no} (règle générale), \esp{Debes pensar} (devoir personnel), \esp{Es importante que valores} (conseil évaluatif + subjonctif). Gradation : contrainte sociale $\rightarrow$ devoir moral $\rightarrow$ valeur personnelle.
\item \textbf{Antithèse} : \esp{conductas de riesgo} vs \esp{sociedad fuerte} ; \esp{presión del grupo} vs \esp{piensa en las consecuencias}.
\item \textbf{Anaphore} (début des phrases solutions) : \esp{Hay que... Debes... Es importante que... Solo así...}
\end{itemize}

\textbf{4. Synthèse.} Texte efficace : ancré localement, langue accessible (A2/B1), progression logique (constat $\rightarrow$ injonction $\rightarrow$ espoir). Utilise parfaitement la grammaire du module (impératifs, modaux). Adapté au programme de 1ère A.
\end{exoresolu}

\coursec{Production guidée : créer un message de prévention}

Dans la section précédente, nous avons appris à commenter un texte de sensibilisation : identifier le thème, relever les procédés de langue (impératifs, structures du conseil), dégager l'intention de l'auteur. Il est temps maintenant de passer du rôle de lecteur au rôle d'\textbf{auteur} : à votre tour de produire un message de prévention. C'est la compétence finale du module : \emph{sensibiliser et conseiller pour des comportements responsables}. Trois productions sont au programme : une \cle{affiche} (\esp{cartel}), un \cle{slogan} (\esp{eslogan}) et un \cle{dialogue} (\esp{diálogo}).

\courssub{Consigne 1 : créer une affiche de prévention pour le lycée}

\textbf{La tâche.} Vous devez créer une affiche de prévention destinée aux élèves de votre lycée (par exemple contre la violence, le tabac, la triche aux examens ou l'usage excessif du téléphone). L'affiche doit contenir obligatoirement :

\begin{itemize}
\item un \textbf{slogan} à l'impératif affirmatif : \esp{¡Di no a la violencia!} « Dis non à la violence ! », \esp{¡Cuida tu salud, cuida tu futuro!} « Prends soin de ta santé, prends soin de ton avenir ! » ;
\item \textbf{trois conseils} avec \esp{deber + infinitivo} ou \esp{hay que + infinitivo} : \esp{Debes respetar a tus compañeros.} « Tu dois respecter tes camarades. » ;
\item \textbf{une interdiction} à l'impératif négatif : \esp{No copies en los exámenes.} « Ne triche pas aux examens. » ;
\item une \textbf{image ou un symbole} (dessiné ou décrit).
\end{itemize}

\begin{methode}[Rédiger une affiche de prévention en quatre étapes]
\begin{enumerate}
\item \textbf{Choisir le thème} : une seule conduite à risque par affiche (la violence, le tabac, l'alcool, la triche, les réseaux sociaux...).
\item \textbf{Trouver le slogan} : verbe à l'impératif + complément court + idée positive. Exemple : \esp{Respeta y serás respetado.} « Respecte et tu seras respecté. »
\item \textbf{Rédiger le corps} : 3 conseils variés (un avec \esp{deber}, un avec \esp{hay que}, un avec \esp{es importante que} + subjonctif) et 1 interdiction (\esp{no} + subjonctif présent).
\item \textbf{Soigner la présentation} : slogan en gros en haut, conseils au milieu, interdiction dans un cadre, symbole visible (une main levée, un cœur, un panneau d'interdiction).
\end{enumerate}
\end{methode}

\begin{popfigure}
\begin{tikzpicture}
\draw[fill=popcream, draw=popdark, line width=1.2pt, rounded corners=4pt] (0,0) rectangle (10,13);
\draw[fill=popblueL, draw=popblue, rounded corners=3pt] (0.5,10.6) rectangle (9.5,12.5);
\node[align=center, text=popdark, font=\bfseries] at (5,11.55) {TITRE / SLOGAN\\ \small ¡Di no a la violencia! (impératif)};
\draw[fill=popgreenL, draw=popgreen, rounded corners=3pt] (0.5,6.2) rectangle (6.2,10.1);
\node[align=center, text=popdark] at (3.35,8.15) {\textbf{3 CONSEILS}\\ \small Debes respetar...\\ \small Hay que dialogar...\\ \small Es importante que ayudes...};
\draw[fill=poporangeL, draw=poporange, rounded corners=3pt] (6.7,6.2) rectangle (9.5,10.1);
\node[align=center, text=popdark] at (8.1,8.15) {\textbf{IMAGE}\\ \textbf{SYMBOLE}\\ \small (main levée,\\ \small cœur, panneau)};
\draw[fill=poppinkL, draw=poppink, rounded corners=3pt] (0.5,2.2) rectangle (9.5,5.7);
\node[align=center, text=popdark] at (5,3.95) {\textbf{INTERDICTION}\\ \small No copies en los exámenes.\\ \small (no + subjonctif présent)};
\draw[fill=popgoldL, draw=popgold, rounded corners=3pt] (0.5,0.4) rectangle (9.5,1.7);
\node[align=center, text=popdark] at (5,1.05) {\small Signature : Campagne du club d'espagnol — Lycée de Yaoundé};
\end{tikzpicture}
\legende{Maquette d'affiche de prévention : chaque zone correspond à une consigne grammaticale précise.}
\end{popfigure}

\courssub{Consigne 2 : les règles d'un bon slogan}

Observez d'abord ces slogans authentiques de campagnes hispanophones :

\begin{exemplebox}[Slogans de campagnes de prévention]
\begin{itemize}
\item \esp{Si bebes, no conduzcas.} « Si tu bois, ne conduis pas. » (campagne routière espagnole DGT)
\item \esp{No fumes: tu salud es tu futuro.} « Ne fume pas : ta santé, c'est ton avenir. »
\item \esp{Con drogas no hay futuro.} « Avec la drogue, il n'y a pas d'avenir. » (classique prévention)
\item \esp{Respeta y serás respetado.} « Respecte et tu seras respecté. »
\item \esp{Debes pensar antes de actuar.} « Tu dois réfléchir avant d'agir. »
\end{itemize}
\end{exemplebox}

Qu'observe-t-on ? Tous ces slogans sont \textbf{courts} (moins de huit mots), \textbf{rythmés}, utilisent l'\textbf{impératif} ou le présent de vérité générale, et sont faciles à \textbf{mémoriser}.

\begin{propriete}[Les quatre règles d'un bon slogan]
\begin{enumerate}
\item \textbf{Court} : 3 à 8 mots maximum.
\item \textbf{Rythmé} : répétition, rime ou structure parallèle (\esp{Cuida tu salud, cuida tu futuro}).
\item \textbf{Impératif} : le verbe d'action à l'impératif interpelle directement le lecteur (\esp{¡Di no!}, \esp{¡Respeta!}, \esp{¡Cuida!}).
\item \textbf{Positif ou frappant} : une idée qui donne envie d'agir ou qui marque les esprits.
\end{enumerate}
Formule magique : \textbf{verbe à l'impératif + complément court + idée positive}.
\end{propriete}

\begin{savaistu}[La rime, arme secrète des slogans espagnols]
Les slogans espagnols jouent souvent sur la rime ou l'assonance : \esp{Con drogas no hay futuro} (« drogas » rime presque avec « futuro » par la sonorité en « o »). Tu peux faire pareil en français et en espagnol pour ton affiche : \esp{No a la violencia, sí a la convivencia} (« Non à la violence, oui à la vie ensemble ») : la rime en \emph{-encia} rend le slogan inoubliable. À toi d'inventer : \esp{Estudia hoy, triunfa mañana} « Étudie aujourd'hui, triomphe demain » !
\end{savaistu}

\courssub{Consigne 3 : écrire et jouer un dialogue de sensibilisation}

\textbf{La tâche.} Par groupes de deux, écrivez un dialogue de 8 à 10 répliques entre un élève tenté par une mauvaise conduite (tricher, fumer, se battre, sécher les cours) et un ami qui le conseille. Utilisez : des impératifs affirmatifs et négatifs, \esp{deber + infinitivo}, \esp{hay que + infinitivo} et \esp{es importante que} + subjonctif.

Lisez d'abord ce dialogue modèle :

\begin{textebox}[Diálogo : « El examen de mañana »]
— Oye, Pablo, mañana hay examen de matemáticas y no he estudiado nada. Voy a copiar del compañero de al lado.\\
— ¡No hagas eso, Carlos! Copiar es una conducta de riesgo. Si te descubren, te expulsan del examen.\\
— Pero es que no sé nada. ¿Qué hago?\\
— Primero, no te pongas nervioso. Debes estudiar esta tarde: todavía hay tiempo.\\
— ¿Tú crees?\\
— Claro. Hay que repasar los ejercicios del cuaderno y las fórmulas importantes.\\
— Es demasiado trabajo para una tarde...\\
— Escucha: es importante que lo intentes. Copiar es fácil, pero no aprendes nada.\\
— Tienes razón. ¿Me ayudas a repasar?\\
— Por supuesto. Ven a mi casa a las cuatro. Y recuerda: ¡no copies nunca! Tu esfuerzo vale más que una nota falsa.\\
— Gracias, amigo. Eres un buen consejero.
\end{textebox}

\textbf{Glossaire} : \esp{copiar} = tricher (copier sur qqn) ; \esp{te expulsan} = on t'exclut (impersonnel 3\up{e} pl) ; \esp{ponerse nervioso} = stresser, s'angoisser ; \esp{repasar} = réviser ; \esp{el cuaderno} = le cahier ; \esp{el esfuerzo} = l'effort ; \esp{una nota falsa} = une note non méritée (trichée) ; \esp{consejero} = conseiller.

\textbf{Questions de compréhension} :
\begin{enumerate}
\item ¿Qué quiere hacer Carlos? (Il veut tricher / copier sur son voisin).
\item ¿Qué consejos le da Pablo? (Ne fais pas ça, ne stresse pas, tu dois étudier, il faut réviser, il est important que tu essaies, viens chez moi, ne triche jamais).
\item ¿Qué estructuras usa Pablo para convencer a su amigo? (Impératifs négatifs \esp{no hagas, no te pongas, no copies}, \esp{deber + inf}, \esp{hay que + inf}, \esp{es importante que + subj}, impératif affirmatif \esp{ven, recuerda}).
\end{enumerate}

\begin{attention}[Adapter le registre : \esp{tú} ou \esp{usted} ?]
Dans un dialogue oral, le choix de l'impératif dépend de la personne à qui tu parles. \textbf{Entre amis / pairs} : impératif en \esp{tú} : \esp{¡No fumes!} « Ne fume pas ! ». \textbf{Avec un adulte, un professeur, une autorité} : impératif en \esp{usted} (= subjonctif présent, 3\up{e} personne du singulier) : \esp{¡Respete las normas, señor!} « Respectez les règles, monsieur ! ». Erreur fréquente du francophone : utiliser \esp{tú} avec un professeur ou un policier — en espagnol, cela peut paraître impoli, familier déplacé. Et attention aux impératifs affirmatifs irréguliers (tú) : \esp{hacer} $\rightarrow$ \esp{haz} (pas \esp{*hace} !), \esp{ir} $\rightarrow$ \esp{ve}, \esp{decir} $\rightarrow$ \esp{di}, \esp{poner} $\rightarrow$ \esp{pon}, \esp{salir} $\rightarrow$ \esp{sal}, \esp{tener} $\rightarrow$ \esp{ten}, \esp{venir} $\rightarrow$ \esp{ven}, \esp{ser} $\rightarrow$ \esp{sé}.
\end{attention}

\begin{center}
\begin{tabularx}{\linewidth}{|X|X|X|}
\hline
\rowcolor{popblueL} \textbf{Situation} & \textbf{Forme} & \textbf{Exemple} \\
\hline
Conseil à un ami (\esp{tú}) & impératif affirmatif & \esp{¡Estudia esta tarde!} \\
\hline
Interdiction à un ami & \esp{no} + subjonctif & \esp{¡No copies!} \\
\hline
Conseil à un adulte (\esp{usted}) & subjonctif, 3\up{e} pers. & \esp{¡Siéntese, por favor!} \\
\hline
Obligation générale & \esp{hay que} + infinitif & \esp{Hay que respetar las normas.} \\
\hline
Devoir personnel & \esp{deber} + infinitif & \esp{Debes pensar antes de actuar.} \\
\hline
Conseil insistant & \esp{es importante que} + subj. & \esp{Es importante que lo intentes.} \\
\hline
\end{tabularx}
\end{center}

\courssub{Grille d'auto-évaluation}

Avant de rendre votre production (affiche ou dialogue), vérifiez chaque point de cette grille. Cochez mentalement : chaque élément absent doit être ajouté.

\begin{aretenir}[Grille d'auto-évaluation de ma production]
\begin{itemize}
\item[$\square$] Au moins \textbf{2 impératifs affirmatifs} (\esp{¡Cuida!}, \esp{¡Respeta!}, \esp{¡Ven!}, \esp{¡Haz!})
\item[$\square$] Au moins \textbf{2 impératifs négatifs} (\esp{No fumes}, \esp{No copies}, \esp{No hagas}, \esp{No te pongas})
\item[$\square$] Au moins \textbf{1 \esp{hay que} + infinitif}
\item[$\square$] Au moins \textbf{1 \esp{deber} + infinitif}
\item[$\square$] Au moins \textbf{1 \esp{es importante que} + subjonctif}
\item[$\square$] Au moins \textbf{5 mots du lexique} de la leçon (\esp{conducta de riesgo}, \esp{salud}, \esp{normas}, \esp{respetar}, \esp{prevenir}, \esp{violencia}, \esp{responsable}, \esp{adicción}, \esp{acoso}, \esp{esfuerzo}...)
\item[$\square$] Registre adapté (\esp{tú} entre amis, \esp{usted} avec un adulte)
\item[$\square$] Accents et ponctuation espagnole (¡ ! ¿ ?)
\end{itemize}
\end{aretenir}

\courssub{Correction collective d'une production type}

Voici une affiche rédigée par un élève de Première. Elle contient cinq erreurs typiques des francophones. Cherchez-les avant de lire la correction.

\begin{exoresolu}[Corrige l'affiche de Mamadou]
Énoncé. Voici le texte de l'affiche de Mamadou : \esp{¡No fumas! El tabaco destruye tu salud. \textbf{Debes de} respetar tu cuerpo. Hay que \textbf{hacer} deporte y comer bien. Es importante que \textbf{hablas} con tus padres si tienes problemas. \textbf{¡Hace} ejercicio cada día!} Retrouve et corrige les cinq erreurs.
\tcblower
Solution détaillée :
\begin{enumerate}
\item \esp{No fumas} $\rightarrow$ \esp{No fumes} : l'impératif négatif exige le \textbf{subjonctif présent} (\esp{fumes} < \esp{fumar}), pas l'indicatif. Calque du français « ne fume pas ».
\item \esp{Debes de respetar} $\rightarrow$ \esp{Debes respetar} : \esp{deber + infinitivo} (obligation/conseil) se construit \textbf{sans} la préposition \esp{de}. (\esp{Deber de} + infinitif exprime la \textbf{probabilité} : \esp{Debe de estar enfermo} = « Il doit être malade ».)
\item \esp{Es importante que hablas} $\rightarrow$ \esp{Es importante que hables} : après \esp{es importante que} (expression de valeur/jugement), toujours le \textbf{subjonctif présent} (\esp{hables} < \esp{hablar}).
\item \esp{¡Hace ejercicio!} $\rightarrow$ \esp{¡Haz ejercicio!} : l'impératif affirmatif de \esp{hacer} à \esp{tú} est \textbf{irrégulier} : \esp{haz}. (Règle : \esp{hacer} $\rightarrow$ \esp{haz}).
\item \esp{Hay que hacer deporte} est correct (infinitif après \esp{hay que}), mais attention à ne jamais conjuguer : \esp{*hay que haces} est faux.
\end{enumerate}
\textbf{Texte corrigé :} \esp{¡No fumes! El tabaco destruye tu salud. Debes respetar tu cuerpo. Hay que hacer deporte y comer bien. Es importante que hables con tus padres si tienes problemas. ¡Haz ejercicio cada día!} « Ne fume pas ! Le tabac détruit ta santé. Tu dois respecter ton corps. Il faut faire du sport et bien manger. Il est important que tu parles avec tes parents si tu as des problèmes. Fais de l'exercice chaque jour ! »
\end{exoresolu}

\begin{exercice}[À toi de jouer !]
\begin{enumerate}
\item Rédige le slogan d'une affiche contre la triche aux examens, avec un impératif et une rime si possible.
\item Écris trois conseils pour un camarade qui passe trop de temps sur son téléphone : un avec \esp{deber}, un avec \esp{hay que}, un avec \esp{es importante que} + subjonctif.
\item Avec un voisin, écris un dialogue de 8 répliques : un élève veut sécher le cours d'espagnol, son ami le conseille. Jouez-le devant la classe.
\end{enumerate}
\end{exercice}

\begin{corrige}[Exercice « À toi de jouer ! » — éléments de réponse]
1. Slogans possibles : \esp{¡No copies: estudia y triunfa!} « Ne triche pas : étudie et triomphe ! » ; \esp{Copiar hoy, fracasar mañana} « Tricher aujourd'hui, échouer demain » (rime en \emph{-ar}) ; \esp{¡Di no a la trampa, sí a tu talento!} « Dis non à la triche, oui à ton talent ! »\\
2. \esp{Debes apagar el teléfono durante las clases.} « Tu dois éteindre le téléphone pendant les cours. » \esp{Hay que limitar el tiempo en las redes sociales.} « Il faut limiter le temps sur les réseaux sociaux. » \esp{Es importante que duermas ocho horas cada noche.} « Il est important que tu dormes huit heures chaque nuit » (subjonctif \esp{duermas} < \esp{dormir}, diphtongue).\\
3. Modèle de dialogue :\\
— \esp{Mañana no voy a la clase de español.} \\
— \esp{¡No hagas eso! Faltar a clase es una mala conducta.} \\
— \esp{Pero no entiendo nada...} \\
— \esp{Debes hablar con el profesor: él puede ayudarte.} \\
— \esp{¿Tú crees?} \\
— \esp{Sí. Hay que trabajar un poco cada día.} \\
— \esp{Es difícil.} \\
— \esp{Es importante que no abandones. ¡Ven conmigo mañana!} \\
— \esp{Está bien, gracias por el consejo.}
\end{corrige}

\begin{experience}[Exposition des affiches de la classe]
Chaque groupe présente son affiche à la classe en deux minutes : lecture du slogan à voix haute, explication des conseils en espagnol. La classe vote pour le meilleur slogan selon les quatre critères (court, rythmé, impératif, mémorable). Les affiches gagnantes sont accrochées dans la salle de classe ou à la vie scolaire pour une vraie campagne de prévention du lycée.
\end{experience}

\begin{aretenir}[L'essentiel de la section]
Un message de prévention efficace repose sur trois piliers : un \textbf{slogan} court à l'impératif, des \textbf{conseils} variés (\esp{deber}, \esp{hay que}, \esp{es importante que} + subjonctif) et une \textbf{interdiction} claire (\esp{no} + subjonctif). Adapte toujours le registre (\esp{tú} / \esp{usted}) et vérifie ta production avec la grille d'auto-évaluation. Tu sais désormais sensibiliser les autres en espagnol : \esp{¡Sé responsable y sensibiliza a los demás!} « Sois responsable et sensibilise les autres ! »
\end{aretenir}

\newpage

\coursec{Activité d'intégration}

\courssub{Objectifs de l'activité}

À l'issue de cette activité d'intégration, tu seras capable de :
\begin{itemize}
\item mobiliser, dans une tâche de communication authentique, l'\cle{impératif} affirmatif et négatif pour conseiller et interdire ;
\item utiliser correctement les structures du \cle{conseil} et de l'\cle{obligation} : \esp{deber + infinitivo}, \esp{hay que + infinitivo}, \esp{es importante que + subjuntivo} ;
\item employer le lexique des \cle{conduites à risque} et des \cle{comportements responsables} étudié dans la leçon ;
\item rédiger une \cle{affiche de prévention} avec un slogan percutant et des conseils clairs ;
\item jouer un court \cle{dialogue de sensibilisation} devant la classe, avec une prononciation soignée ;
\item évaluer ta propre production à l'aide de la grille d'auto-évaluation de la section 6.
\end{itemize}

\courssub{Situation}

Tu es élève en classe de Première A au Lycée Bilingue de Yaoundé. À l'approche de la \emph{Semaine de la vie citoyenne}, la conseillère d'orientation et le bureau des élèves ont constaté des comportements inquiétants : certains élèves fument derrière les salles de classe, d'autres trichent pendant les évaluations, d'autres encore passent la nuit sur les réseaux sociaux et s'endorment en cours. Le proviseur lance donc un concours : chaque classe doit concevoir une campagne de sensibilisation intitulée \esp{« Jóvenes responsables, sociedad fuerte »}. La meilleure affiche et le meilleur slogan seront affichés à l'entrée du lycée et publiés dans le journal des élèves.

Voici l'appel à participation affiché au tableau d'information du lycée :

\begin{textebox}[Appel à participation — Lycée Bilingue de Yaoundé]
\textbf{CONVOCATORIA: CAMPAÑA DE SENSIBILIZACIÓN «JÓVENES RESPONSABLES, SOCIEDAD FUERTE»}

Queridos alumnos:

Con motivo de la Semana de la Vida Ciudadana, el liceo organiza un concurso de mensajes de prevención. Por grupos de cuatro, debéis crear:

\begin{enumerate}
\item Un cartel con un eslogan en imperativo y tres consejos para los jóvenes.
\item Un diálogo corto de sensibilización para representar delante de la clase.
\end{enumerate}

Temas posibles: el tabaco y el alcohol, la trampa en los exámenes, el uso excesivo del teléfono móvil, la violencia en el recreo, el respeto de los demás.

El mejor eslogan será elegido por votación y se colocará en la entrada del liceo. ¡Participad! Vuestras ideas pueden cambiar las conductas de vuestros compañeros.

El director del liceo
\end{textebox}

\textbf{Glossaire :} \esp{convocatoria} = appel à participation ; \esp{cartel} = affiche ; \esp{eslogan} = slogan ; \esp{la trampa} = la triche ; \esp{el recreo} = la récréation ; \esp{colocar} = afficher, placer ; \esp{compañeros} = camarades.

\courssub{Consignes}

Travaillez par groupes de quatre. Suivez les étapes dans l'ordre.

\begin{enumerate}
\item \textbf{Compréhension et choix du thème.} Lisez l'appel à participation. En groupe, choisissez \emph{un seul} thème parmi ceux proposés et justifiez votre choix en deux phrases en espagnol : \esp{Elegimos este tema porque...}
\item \textbf{Remue-méninges lexical.} Sur une feuille, dressez deux colonnes : les \cle{conduites à risque} liées à votre thème (minimum 4 mots ou expressions) et les \cle{comportements responsables} à promouvoir (minimum 4 mots ou expressions). Utilisez le lexique de la section 2 de la leçon.
\item \textbf{Création de l'affiche.} Rédigez :
\begin{enumerate}
\item un \cle{slogan} court à l'impératif (affirmatif ou négatif), par exemple avec \esp{tú} ou \esp{vosotros} ;
\item trois \cle{conseils} : un avec \esp{deber + infinitivo}, un avec \esp{hay que + infinitivo}, un avec \esp{es importante que + subjuntivo}.
\end{enumerate}
Ajoutez un titre et, si possible, un dessin simple. Soignez l'orthographe et les accents.
\item \textbf{Écriture du dialogue.} Écrivez un dialogue de 8 à 10 répliques : un élève hésite à adopter une conduite à risque (tricher, fumer, sortir sans prévenir...), son ami(e) le conseille avec l'impératif et les structures du conseil. Le dialogue doit se terminer par une décision responsable.
\item \textbf{Répétition et représentation.} Répartissez les rôles, répétez en veillant à la prononciation (le \esp{h} est muet, le \esp{j} se prononce « r » guttural, le \esp{c} devant \esp{e/i} se prononce « s » en Amérique latine), puis jouez la scène devant la classe.
\item \textbf{Vote et auto-évaluation.} Chaque groupe vote pour le meilleur slogan de la classe (on ne peut pas voter pour soi). Le slogan gagnant est écrit au tableau. Enfin, relisez votre production écrite avec la grille d'auto-évaluation de la section 6 et corrigez-la avant de la rendre au professeur.
\end{enumerate}

\courssub{Ressources à mobiliser}

\begin{aretenir}[Boîte à outils de la campagne]
\textbf{Impératif affirmatif (tú)} : 3\up{e} personne du singulier du présent : \esp{respeta}, \esp{estudia}, \esp{evita}. Irréguliers : \esp{decir → di}, \esp{hacer → haz}, \esp{ir → ve}, \esp{poner → pon}, \esp{salir → sal}, \esp{ser → sé}, \esp{tener → ten}, \esp{venir → ven}.

\textbf{Impératif négatif} : \esp{no} + subjonctif présent : \esp{no fumes}, \esp{no hagas trampa}, \esp{no pierdas el tiempo}.

\textbf{Conseil et obligation} :
\begin{itemize}
\item \esp{Debes respetar a tus profesores.} « Tu dois respecter tes professeurs. »
\item \esp{Hay que dormir ocho horas.} « Il faut dormir huit heures. »
\item \esp{Es importante que cuides tu salud.} « Il est important que tu prennes soin de ta santé. »
\end{itemize}

\textbf{Lexique utile} : \esp{fumar}, \esp{el tabaco}, \esp{el alcohol}, \esp{hacer trampa}, \esp{la violencia}, \esp{respetar}, \esp{ser solidario}, \esp{estudiar con regularidad}, \esp{proteger la salud}, \esp{decir no}.
\end{aretenir}

\begin{attention}[Pièges à éviter]
\begin{itemize}
\item Ne dis pas \esp{no fumes} avec un infinitif : l'impératif négatif exige le subjonctif (\esp{no fumes}, jamais \esp{no fumar} dans un conseil direct à \esp{tú}).
\item Après \esp{es importante que}, le verbe est au \cle{subjonctif} : \esp{es importante que vengas}, et non \esp{que vienes}.
\item Faux ami : \esp{asistir a} signifie « assister à » (être présent), pas « aider » ; pour « aider », utilise \esp{ayudar}.
\item N'oublie pas les accents : \esp{salud}, \esp{educación}, \esp{prevención}.
\end{itemize}
\end{attention}

\courssub{Proposition de production et corrigé commenté}

\begin{exoresolu}[Modèle de campagne : « Di no a la trampa »]
Énoncé : produis l'affiche (slogan + 3 conseils) et le dialogue de sensibilisation sur le thème de la triche aux examens.

\tcblower

\textbf{1. L'affiche (production écrite modèle)}

\textbf{CAMPAÑA DEL LICEO BILINGÜE DE YAUNDÉ}

\textbf{Eslogan : ¡Di no a la trampa, sé honesto contigo mismo!}

\begin{itemize}
\item \esp{Debes estudiar con regularidad durante todo el año.}
\item \esp{Hay que preparar los exámenes sin copiar a los demás.}
\item \esp{Es importante que confíes en tus propias capacidades.}
\end{itemize}

\textbf{2. Le dialogue (production orale modèle)}

\textbf{Mamadou :} Oye, Amina, mañana tenemos el examen de matemáticas y no he estudiado nada. Voy a copiar a mi vecino.

\textbf{Amina :} ¡No hagas eso, Mamadou! Hacer trampa es una conducta de riesgo.

\textbf{Mamadou :} Pero si suspendo, mis padres van a castigarme.

\textbf{Amina :} Escúchame: debes ser honesto. Si copias y te descubren, el castigo será peor.

\textbf{Mamadou :} Tienes razón, pero es demasiado tarde para estudiar.

\textbf{Amina :} No, no es tarde. Hay que aprovechar esta tarde. Ven a mi casa y repasamos juntos.

\textbf{Mamadou :} ¿De verdad me ayudas?

\textbf{Amina :} Claro. Es importante que seamos solidarios entre amigos. Y recuerda: no copies nunca más.

\textbf{Mamadou :} Gracias, Amina. A partir de hoy, voy a estudiar con regularidad. ¡La trampa no es una solución!

\textbf{Amina :} ¡Muy bien! Un joven responsable construye una sociedad fuerte.

\textbf{3. Commentaire des choix de langue}

\begin{itemize}
\item \textbf{Slogan} : deux impératifs de \esp{tú} : \esp{di} (irrégulier, de \esp{decir}) et \esp{sé} (irrégulier, de \esp{ser}) — formes courtes et percutantes, idéales pour une affiche.
\item \textbf{Conseils} : les trois structures de la leçon sont présentes : \esp{debes estudiar} (\esp{deber} + infinitif), \esp{hay que preparar} (obligation impersonnelle), \esp{es importante que confíes} (subjonctif après \esp{es importante que}).
\item \textbf{Dialogue} : \esp{¡No hagas eso!} emploie l'impératif négatif (subjonctif de \esp{hacer}) ; \esp{Escúchame} montre le pronom \esp{me} accolé à l'impératif affirmatif avec accent d'intensité ; \esp{no copies nunca más} renforce l'interdiction.
\item \textbf{Lexique} : \esp{hacer trampa}, \esp{conducta de riesgo}, \esp{ser solidario}, \esp{estudiar con regularidad} reprennent le vocabulaire de la section 2.
\item \textbf{Contexte} : les prénoms et le cadre (examen de mathématiques, sanction familiale) rendent la scène crédible pour des élèves camerounais.
\end{itemize}
\end{exoresolu}

\courssub{Bilan de l'activité}

Cette activité t'a permis de réinvestir l'ensemble des acquis de la leçon dans une tâche de communication réelle : convaincre et conseiller. Vérifie que tu sais désormais :

\begin{itemize}
\item former l'impératif affirmatif de \esp{tú}, y compris les huit formes irrégulières, et l'impératif négatif avec \esp{no} + subjonctif ;
\item construire un conseil avec \esp{deber + infinitivo}, une obligation générale avec \esp{hay que} et une recommandation avec \esp{es importante que + subjuntivo} ;
\item rédiger un slogan court, mémorable et correctement accentué ;
\item tenir un rôle dans un dialogue de sensibilisation avec une prononciation claire ;
\item évaluer et corriger ta propre production grâce à la grille d'auto-évaluation.
\end{itemize}

\begin{savaistu}[Et après ?]
Les campagnes de prévention sont une tradition vivante dans le monde hispanophone : en Espagne, le ministère de la Santé lance chaque année des slogans célèbres contre le tabac et l'alcool, et en Guinée équatoriale, pays hispanophone voisin du Cameroun, les écoles organisent des journées de sensibilisation où les élèves créent leurs propres messages. Ton affiche s'inscrit dans cette grande famille de la communication citoyenne en espagnol.
\end{savaistu}

\newpage

\coursec{Méthodes et exercices résolus}

\courssub{Savoir-faire 1 : Donner des conseils et des consignes à l'impératif}

\begin{methode}[Transformer un conseil en consigne à l'impératif]
Pour donner un conseil ou une consigne en espagnol, suis ces étapes :
\begin{enumerate}
\item \cle{Identifier le destinataire} : à qui parles-tu ? À un camarade (\esp{tú}), à un groupe d'élèves (\esp{vosotros}), à un adulte ou au public d'une affiche (\esp{usted}) ?
\item \cle{Choisir le registre} : entre amis, tutoiement ; dans un message officiel de prévention, \esp{usted} est plus poli et plus solennel.
\item \cle{Former l'impératif affirmatif} :
\begin{itemize}
\item \esp{tú} : 3\up{e} personne du singulier du présent (\esp{habla}, \esp{come}, \esp{escribe}) ; irréguliers : \esp{di} (decir), \esp{haz} (hacer), \esp{ve} (ir), \esp{pon} (poner), \esp{sal} (salir), \esp{sé} (ser), \esp{ten} (tener), \esp{ven} (venir).
\item \esp{vosotros} : infinitif dont le \esp{-r} final devient \esp{-d} : \esp{hablad, comed, escribid}.
\item \esp{usted} : subjonctif présent, 3\up{e} personne du singulier : \esp{hable, coma, escriba}.
\end{itemize}
\item \cle{Former l'impératif négatif} : toujours \esp{no} + subjonctif présent : \esp{no fumes}, \esp{no bebáis}, \esp{no conduzca}.
\item \cle{Placer les pronoms} : à l'affirmatif, ils se collent après le verbe (accent sur la voyelle tonique : \esp{levántate}, \esp{dímelo}) ; au négatif, ils se placent avant : \esp{no te levantes}, \esp{no me lo digas}.
\item \cle{Vérifier l'accent} écrit quand on ajoute un pronom : \esp{hazlo}, \esp{respétalos}, \esp{cuídate}.
\end{enumerate}
\end{methode}

\begin{aretenir}[Tableau récapitulatif de l'impératif]
\begin{center}
\begin{tabularx}{\linewidth}{|l|X|X|X|}
\hline
\rowcolor{popblueL} Personne & \esp{hablar} & \esp{beber} & \esp{salir} (irrégulier à \esp{tú}) \\
\hline
\esp{tú} (affirmatif) & \esp{habla} & \esp{bebe} & \esp{sal} \\
\hline
\esp{tú} (négatif) & \esp{no hables} & \esp{no bebas} & \esp{no salgas} \\
\hline
\esp{vosotros} (affirmatif) & \esp{hablad} & \esp{bebed} & \esp{salid} \\
\hline
\esp{vosotros} (négatif) & \esp{no habléis} & \esp{no bebáis} & \esp{no salgáis} \\
\hline
\esp{usted} (affirmatif) & \esp{hable} & \esp{beba} & \esp{salga} \\
\hline
\esp{usted} (négatif) & \esp{no hable} & \esp{no beba} & \esp{no salga} \\
\hline
\end{tabularx}
\end{center}
Observation essentielle : seules les formes affirmatives de \esp{tú} et de \esp{vosotros} ont des formes propres ; toutes les autres (et tous les négatifs) empruntent le subjonctif présent.
\end{aretenir}

\begin{exoresolu}[Des conseils pour la santé]
Énoncé. Transforme les phrases suivantes en consignes à l'impératif, selon le destinataire indiqué entre parenthèses.
\begin{enumerate}
\item Es bueno dormir ocho horas. (un amigo)
\item Es malo fumar en el patio del colegio. (los alumnos)
\item Es necesario respetar a los profesores. (usted)
\item Es peligroso enviar mensajes mientras se conduce una moto. (un amigo)
\item Hay que lavarse las manos antes de comer. (usted)
\end{enumerate}
\tcblower
Solution détaillée.
\begin{enumerate}
\item Destinataire : un ami, donc \esp{tú}. Verbe \esp{dormir} : 3\up{e} pers. du singulier du présent \esp{duerme} ; l'impératif affirmatif de \esp{tú} reprend cette forme. Réponse : \esp{Duerme ocho horas.} « Dors huit heures. »
\item Destinataire : les élèves, donc \esp{vosotros}, et c'est une interdiction : impératif négatif = \esp{no} + subjonctif. Subjonctif de \esp{fumar} à la 2\up{e} pers. du pluriel : \esp{fuméis}. Réponse : \esp{No fuméis en el patio del colegio.} « Ne fumez pas dans la cour du collège. »
\item Destinataire : \esp{usted}, conseil positif : subjonctif 3\up{e} pers. du singulier de \esp{respetar} : \esp{respete}. Réponse : \esp{Respete a los profesores.} « Respectez les professeurs. » Attention à la préposition \esp{a} devant les personnes (complément d'objet direct de personne).
\item Interdiction à un ami : \esp{no} + subjonctif de \esp{enviar} (2\up{e} pers. sing. : \esp{envíes}, accent obligatoire sur le \esp{i}). Réponse : \esp{No envíes mensajes mientras conduces la moto.} « N'envoie pas de messages pendant que tu conduis la moto. »
\item Conseil à \esp{usted} avec un verbe pronominal : à l'affirmatif, le pronom \esp{se} se colle après le verbe, ce qui impose un accent écrit. Subjonctif de \esp{lavar} : \esp{lave} ; avec pronom : \esp{lávese}. Réponse : \esp{Lávese las manos antes de comer.} « Lavez-vous les mains avant de manger. »
\end{enumerate}
Conclusion : le choix de la personne (\esp{tú}, \esp{vosotros}, \esp{usted}) détermine toute la forme ; le négatif exige toujours le subjonctif ; les pronoms collés à l'affirmatif provoquent un accent écrit.
\end{exoresolu}

\courssub{Savoir-faire 2 : Exprimer le conseil et l'obligation}

\begin{methode}[Conseiller avec deber, hay que, es importante que]
Pour varier les formulations d'un message de prévention, mobilise trois structures :
\begin{enumerate}
\item \esp{deber + infinitif} : devoir moral, conseil personnel. \esp{Debes estudiar más.} « Tu dois travailler davantage. » Se conjugue selon la personne : \esp{debo, debes, debe, debemos, debéis, deben}.
\item \esp{hay que + infinitif} : obligation générale, impersonnelle, idéale pour les règles collectives. \esp{Hay que respetar las normas.} « Il faut respecter les règles. » Ne s'emploie qu'à la 3\up{e} personne du singulier ; la négation porte sur \esp{hay} : \esp{no hay que fumar} « il ne faut pas fumer ».
\item \esp{es importante que + subjonctif} : conseil nuancé, insistant. \esp{Es importante que cuides tu salud.} « Il est important que tu prennes soin de ta santé. » Après \esp{que}, toujours le subjonctif.
\item \cle{Choisir selon l'effet recherché} : \esp{deber} s'adresse à une personne précise ; \esp{hay que} concerne tout le monde ; \esp{es importante que} permet de nommer le destinataire dans la subordonnée.
\item \cle{Vérifier le mode} : infinitif après \esp{deber} et \esp{hay que} ; subjonctif après \esp{es importante que}, \esp{es necesario que}, \esp{conviene que}.
\end{enumerate}
\end{methode}

\begin{exoresolu}[Thème : conseils à un jeune conducteur de moto-taxi]
Énoncé. Traduis en espagnol.
\begin{enumerate}
\item Tu dois porter un casque quand tu conduis.
\item Il ne faut pas boire d'alcool avant de conduire.
\item Il est important que tu respectes les feux tricolores.
\item Nous devons protéger notre vie et celle des passagers.
\item Il est nécessaire que les jeunes évitent les conduites à risque.
\end{enumerate}
\tcblower
Solution détaillée.
\begin{enumerate}
\item « Tu dois » = \esp{debes + infinitif}. « Porter un casque » : \esp{llevar un casco} (faux ami : \esp{portar} existe mais \esp{llevar} est l'usage courant pour les vêtements et accessoires). « Quand tu conduis » : \esp{cuando conduces} (indicatif, habitude). Réponse : \esp{Debes llevar un casco cuando conduces.}
\item « Il ne faut pas » = \esp{no hay que + infinitif}. « Boire de l'alcool » : \esp{beber alcohol} (pas de partitif en espagnol). « Avant de conduire » : \esp{antes de conducir} (infinitif après préposition, comme en français). Réponse : \esp{No hay que beber alcohol antes de conducir.}
\item « Il est important que » + subjonctif : \esp{es importante que respetes}. « Les feux tricolores » : \esp{los semáforos}. Réponse : \esp{Es importante que respetes los semáforos.}
\item « Nous devons » : \esp{debemos + infinitif}. « Celle des passagers » : on évite la répétition avec l'article : \esp{la de los pasajeros}. Réponse : \esp{Debemos proteger nuestra vida y la de los pasajeros.}
\item « Il est nécessaire que » + subjonctif : \esp{es necesario que los jóvenes eviten}. « Les conduites à risque » : \esp{las conductas de riesgo}. Réponse : \esp{Es necesario que los jóvenes eviten las conductas de riesgo.}
\end{enumerate}
Conclusion : trois outils suffisent pour tout conseil : \esp{deber} (personnel), \esp{hay que} (général), \esp{es importante/necesario que} + subjonctif (insistant). Le piège majeur est d'oublier le subjonctif dans la troisième structure.
\end{exoresolu}

\courssub{Savoir-faire 3 : Commenter un texte de sensibilisation}

\begin{methode}[La fiche de commentaire de texte]
\begin{enumerate}
\item \cle{Lire deux fois} le texte : une fois pour le sens global, une fois en soulignant le thème, le destinataire, les mots-clés et les procédés (impératifs, chiffres, exemples).
\item \cle{Présenter le texte} en 3 ou 4 lignes : nature du document (article, campagne, discours), thème, auteur ou source, destinataire.
\item \cle{Dégager l'idée principale} en une phrase : que veut obtenir l'auteur ? Ici : sensibiliser les jeunes aux comportements responsables.
\item \cle{Analyser la structure} : problème exposé, conséquences, solutions ou conseils. Repérer les articulations logiques (\esp{sin embargo}, \esp{por eso}, \esp{además}).
\item \cle{Étudier la langue} : impératifs et structures de conseil, champ lexical du risque et de la responsabilité, ton (grave, encourageant).
\item \cle{Donner son avis} argumenté en deux ou trois phrases, avec \esp{en mi opinión}, \esp{me parece que}, \esp{estoy de acuerdo porque}.
\item \cle{Conclure} : bilan et ouverture (que peut faire le lecteur, l'école, la famille ?).
\end{enumerate}
\end{methode}

\begin{exoresolu}[Commentaire guidé : « Jóvenes responsables, sociedad fuerte »]
Énoncé. Lis le texte suivant, puis réponds aux questions en espagnol.

\esp{En muchos barrios de Yaoundé, los jóvenes se enfrentan cada día a tentaciones peligrosas: el alcohol, el tabaco, las carreras de motos sin casco. Sin embargo, ser joven no significa ser imprudente. Un joven responsable cuida su salud, respeta a los demás y piensa en su futuro. Por eso, hay que decir no a las conductas de riesgo y sí a la vida. La escuela y la familia deben trabajar juntas: es importante que los padres hablen con sus hijos y que los profesores den el ejemplo. Jóvenes: vuestra responsabilidad de hoy construye la sociedad de mañana.}

Questions : 1. ¿Cuál es el tema del texto? 2. ¿A quién se dirige el autor? 3. Cita dos estructuras de consejo del texto. 4. ¿Qué papel tienen la escuela y la familia? 5. ¿Estás de acuerdo con el autor? Justifica.
\tcblower
Solution détaillée.
\begin{enumerate}
\item Thème : les comportements à risque des jeunes et l'appel à la responsabilité. Réponse rédigée : \esp{El tema del texto es la sensibilización de los jóvenes contra las conductas de riesgo, como el alcohol, el tabaco y las carreras de motos.}
\item Destinataire : la dernière phrase s'adresse directement aux jeunes (\esp{vuestra responsabilidad}, possessif de \esp{vosotros}). Réponse : \esp{El autor se dirige sobre todo a los jóvenes, pero también a los padres y a los profesores.}
\item On relève \esp{hay que decir no a las conductas de riesgo} (obligation générale) et \esp{es importante que los padres hablen con sus hijos} (subjonctif \esp{hablen} après \esp{es importante que}). On peut citer aussi \esp{deben trabajar juntas} (\esp{deber + infinitif}).
\item Réponse : \esp{La escuela y la familia deben trabajar juntas: los padres tienen que hablar con sus hijos y los profesores deben dar el ejemplo.}
\item Avis personnel, par exemple : \esp{Sí, estoy de acuerdo con el autor porque un joven responsable protege su vida y su futuro. En mi opinión, las campañas de prevención en los colegios son muy útiles.}
\end{enumerate}
Conclusion : un bon commentaire suit toujours le même plan : présentation, idée principale, analyse des procédés, avis argumenté. Citer le texte (ici les structures de conseil) prouve qu'on l'a réellement lu.
\end{exoresolu}

\courssub{Savoir-faire 4 : Rédiger un message de prévention (affiche, slogan)}

\begin{methode}[Construire une affiche de prévention efficace]
\begin{enumerate}
\item \cle{Définir la cible et le danger} : un seul risque par affiche (tabac, alcool au volant, moto sans casque, réseaux sociaux).
\item \cle{Rédiger un slogan court} : 5 à 10 mots, mémorable, souvent à l'impératif ou avec \esp{no} : \esp{Di no al tabaco, di sí a la vida.}
\item \cle{Ajouter 2 ou 3 consignes} à l'impératif (\esp{usted} pour le grand public) : \esp{Protéjase. Informe a sus hijos.}
\item \cle{Ajouter une phrase d'obligation} avec \esp{hay que} ou \esp{deber} : \esp{Hay que hablar de los peligros en familia.}
\item \cle{Terminer par une formule d'espoir} : \esp{Tu futuro está en tus manos.}
\item \cle{Relire} : accents, subjonctif après \esp{es importante que}, pas de calque du français.
\end{enumerate}
\end{methode}

\begin{exoresolu}[Création d'une affiche contre l'alcool au volant]
Énoncé. Pour la Semaine de la Sécurité Routière de ton lycée à Douala, rédige en espagnol une affiche de prévention contre l'alcool au volant : un slogan, trois consignes et une phrase de conclusion (6 à 10 lignes).
\tcblower
Solution détaillée. Raisonnement : la cible est le grand public (lycéens, parents, conducteurs de motos-taxis) ; on choisit donc \esp{usted} pour les consignes et \esp{hay que} pour la règle générale. Le slogan doit être bref et rythmé.

Modèle de réponse rédigée :

\esp{¡ALCOHOL Y VOLANTE, PELIGRO CONSTANTE!}

\esp{Si bebe, no conduzca. Respete los límites de velocidad y póngase el cinturón de seguridad. En la moto, lleve siempre el casco. Hay que proteger la vida: es importante que cada conductor piense en su familia antes de beber.}

\esp{Su vida vale más que una copa. ¡Conduzca con responsabilidad!}

Traduction : « Alcool et volant, danger constant ! Si vous buvez, ne conduisez pas. Respectez les limites de vitesse et attachez votre ceinture de sécurité. À moto, portez toujours le casque. Il faut protéger la vie : il est important que chaque conducteur pense à sa famille avant de boire. Votre vie vaut plus qu'un verre. Conduisez avec responsabilité ! »

Justifications grammaticales : \esp{no conduzca} = impératif négatif de \esp{usted} (subjonctif de \esp{conducir}, verbe en \esp{-cir} dont la racine prend \esp{zc} : \esp{conduzco}, d'où \esp{conduzca}) ; \esp{póngase} = impératif affirmatif irrégulier \esp{ponga} + pronom \esp{se} collé, d'où l'accent ; \esp{piense} = subjonctif après \esp{es importante que} ; \esp{vale} = présent de \esp{valer}. Conclusion : l'affiche respecte le plan slogan, consignes, obligation, espoir.
\end{exoresolu}

\courssub{Savoir-faire 5 : Jouer un dialogue de sensibilisation}

\begin{methode}[Construire un dialogue de conseil entre amis]
\begin{enumerate}
\item \cle{Situer la scène} : deux amis, un problème (un camarade fume, boit, séche les cours), un conseiller et un conseillé.
\item \cle{Ouvrir le dialogue} par une constatation : \esp{¿Qué te pasa? Te veo triste.}
\item \cle{Faire parler le problème} : le conseillé explique la tentation ou la mauvaise habitude.
\item \cle{Donner les conseils} avec l'impératif de \esp{tú} et les structures étudiées : \esp{No fumes más. Debes hablar con tus padres. Es importante que pidas ayuda.}
\item \cle{Prévoir une réaction} : hésitation (\esp{Es difícil...}), puis acceptation (\esp{Tienes razón, voy a intentarlo}).
\item \cle{Clore positivement} : \esp{Cuenta conmigo. Juntos podemos.}
\item \cle{À l'oral} : soigner l'intonation montante des questions, la prononciation de \esp{j} (« r » grasseyé) et de \esp{ñ} (« gn »).
\end{enumerate}
\end{methode}

\begin{exoresolu}[Dialogue : convaincre un ami d'arrêter de fumer]
Énoncé. Ton ami Carlos fume derrière le gymnase du lycée. Écris en espagnol un dialogue de 10 à 12 répliques dans lequel tu le sensibilises aux dangers du tabac et tu lui donnes des conseils. Utilise au moins deux impératifs, une structure avec \esp{deber} et une avec \esp{es importante que}.
\tcblower
Solution détaillée. Raisonnement : dialogue entre amis, donc \esp{tú} partout ; on alterne constatation, explication du risque, conseils variés, réaction et conclusion positive.

Modèle de réponse rédigée :

\esp{— Carlos, ¿qué haces aquí detrás del gimnasio? ¿Estás fumando?\\
— Sí... todos mis amigos fuman, no es tan grave.\\
— ¡Claro que es grave! El tabaco destruye los pulmones y cuesta mucho dinero. No fumes más, por favor.\\
— Pero no puedo parar, es muy difícil.\\
— Debes pedir ayuda: habla con tus padres o con el médico del colegio.\\
— ¿Tú crees que pueden ayudarme?\\
— Seguro. Es importante que hables con alguien de confianza. Y haz deporte conmigo: el fútbol es mejor que el tabaco.\\
— Tienes razón. Voy a intentarlo desde mañana.\\
— Perfecto. Cuenta conmigo, amigo. ¡Di no al tabaco y sí a la vida!}

Traduction : « — Carlos, qu'est-ce que tu fais derrière le gymnase ? Tu fumes ? — Oui... tous mes amis fument, ce n'est pas si grave. — Bien sûr que c'est grave ! Le tabac détruit les poumons et coûte cher. Ne fume plus, je t'en prie. — Mais je ne peux pas arrêter, c'est très difficile. — Tu dois demander de l'aide : parle à tes parents ou au médecin du lycée. — Tu crois qu'ils peuvent m'aider ? — Bien sûr. Il est important que tu parles à quelqu'un de confiance. Et fais du sport avec moi : le football vaut mieux que le tabac. — Tu as raison. Je vais essayer dès demain. — Parfait. Compte sur moi, mon ami. Dis non au tabac et oui à la vie ! »

Justifications : \esp{no fumes} (impératif négatif, subjonctif), \esp{habla} et \esp{haz} (impératifs affirmatifs de \esp{tú}, \esp{haz} irrégulier), \esp{debes pedir} (\esp{deber + infinitif}), \esp{es importante que hables} (subjonctif), \esp{voy a intentarlo} (futur proche + pronom collé à l'infinitif). Conclusion : le dialogue respecte la progression constatation, risque, conseils, acceptation, espoir.
\end{exoresolu}

\begin{attention}[Les 5 erreurs qui coûtent des points au Bac]
\begin{enumerate}
\item \cle{Oublier le subjonctif} après \esp{es importante que} : on écrit \esp{es importante que hables}, jamais \esp{es importante que hablas}. Toute expression impersonnelle de conseil ou de nécessité suivie de \esp{que} exige le subjonctif.
\item \cle{Confondre impératif affirmatif et négatif de tú} : \esp{fuma} (affirmatif) mais \esp{no fumes} (négatif = subjonctif). Le négatif ne reprend jamais la forme affirmative.
\item \cle{Oublier l'accent} quand on colle un pronom à l'impératif : \esp{levántate}, \esp{póngase}, \esp{dímelo}. Sans accent, la forme est fautive.
\item \cle{Calquer le français} : « prendre soin de » se dit \esp{cuidar} ou \esp{cuidarse}, pas \esp{tomar cuidado} ; « il faut » se traduit \esp{hay que}, pas \esp{es necesario} suivi d'un infinitif ; « conduire » est \esp{conducir}, pas \esp{conduir}.
\item \cle{Faux amis et accords} : \esp{asistir} signifie « assister à » (être présent), pas « aider » (\esp{ayudar}) ; \esp{una conducta responsable}, \esp{los jóvenes responsables} : l'adjectif s'accorde toujours avec son nom.
\end{enumerate}
\end{attention}

\newpage

\coursec{Exercices}

\courssub{Je vérifie mes connaissances}

\begin{exercice}[QCM : la bonne forme de l'impératif]
Pour chaque phrase, entoure la bonne réponse.
\begin{enumerate}
\item ¡\_\_\_ la verdad siempre ! \quad a) \esp{Dices} \quad b) \esp{Di} \quad c) \esp{Decí}
\item ¡No \_\_\_ trampa en el examen ! \quad a) \esp{haces} \quad b) \esp{hagas} \quad c) \esp{haces no}
\item ¡\_\_\_ aquí, por favor ! (usted) \quad a) \esp{Ven} \quad b) \esp{Venga} \quad c) \esp{Vengas}
\item ¡\_\_\_ deporte todos los días ! (vosotros) \quad a) \esp{Haced} \quad b) \esp{Hacéis} \quad c) \esp{Hagan}
\item ¡\_\_\_ lo ahora mismo ! \quad a) \esp{Lo haz} \quad b) \esp{Hazlo} \quad c) \esp{Haz lo}
\end{enumerate}
\end{exercice}

\begin{exercice}[Vrai ou faux ?]
Réponds par VRAI ou FAUX et justifie ta réponse en une phrase.
\begin{enumerate}
\item À l'impératif négatif, le pronom se place après le verbe : \esp{No lo hagas} est incorrect, il faut dire \esp{No hagas lo}.
\item \esp{Hay que} exprime une obligation générale et impersonnelle.
\item Après \esp{es importante que}, on utilise l'indicatif présent.
\item L'impératif affirmatif de \esp{ir} à la personne \esp{tú} est \esp{ve}.
\item \esp{Deber + infinitivo} sert à exprimer un conseil ou une obligation morale.
\end{enumerate}
\end{exercice}

\begin{exercice}[Vocabulaire : conduites à risque]
Relie chaque mot espagnol à sa traduction française.
\begin{center}
\begin{tabularx}{\linewidth}{|X|X|}
\hline
\rowcolor{popblueL} \textbf{Español} & \textbf{Français} \\
\hline
1. \esp{el acoso escolar} & a) la triche \\
\hline
2. \esp{el consumo de drogas} & b) le harcèlement scolaire \\
\hline
3. \esp{la trampa} & c) la consommation de drogues \\
\hline
4. \esp{el abandono de los estudios} & d) le respect des autres \\
\hline
5. \esp{el respeto a los demás} & e) l'abandon des études \\
\hline
\end{tabularx}
\end{center}
\end{exercice}

\begin{exercice}[Conjugaison : à l'impératif (tú)]
Mets les verbes entre parenthèses à l'impératif affirmatif ou négatif selon le sens de la phrase.
\begin{enumerate}
\item (decir) \_\_\_\_\_\_\_\_ la verdad a tus padres.
\item No (hacer) \_\_\_\_\_\_\_\_ trampa en los exámenes.
\item (venir) \_\_\_\_\_\_\_\_ aquí, quiero hablar contigo.
\item No (drogarse) \_\_\_\_\_\_\_\_, es muy peligroso.
\item (ser) \_\_\_\_\_\_\_\_ responsable en la escuela.
\item No (salir) \_\_\_\_\_\_\_\_ por la noche sin permiso.
\item (poner) \_\_\_\_\_\_\_\_ atención en clase.
\item No (irse) \_\_\_\_\_\_\_\_ sin despedirte.
\end{enumerate}
\end{exercice}

\courssub{Je m'entraîne}

\begin{exercice}[Thème : traduis en espagnol]
Traduis les phrases suivantes en espagnol. Attention à la place des pronoms et aux structures du conseil.
\begin{enumerate}
\item Ne fume pas, c'est dangereux pour ta santé.
\item Tu dois faire du sport et il faut respecter les autres.
\item Il est important que tu parles à tes parents.
\item Ne quitte pas l'école, travaille dur et écoute tes professeurs.
\end{enumerate}
\end{exercice}

\begin{exercice}[Transformation : du conseil à l'interdiction]
Transforme chaque phrase : remplace la structure du conseil par un impératif négatif à la personne \esp{tú}, en plaçant correctement les pronoms si nécessaire.
\begin{enumerate}
\item \esp{No debes fumar.} $\rightarrow$ ¡No \_\_\_\_\_\_\_\_ !
\item \esp{No debes hacer trampa.} $\rightarrow$ ¡No \_\_\_\_\_\_\_\_ !
\item \esp{No debes decir mentiras.} $\rightarrow$ ¡No \_\_\_\_\_\_\_\_ !
\item \esp{No debes beber alcohol.} $\rightarrow$ ¡No \_\_\_\_\_\_\_\_ !
\item \esp{No debes abandonar los estudios.} $\rightarrow$ ¡No \_\_\_\_\_\_\_\_ !
\end{enumerate}
\end{exercice}

\begin{exercice}[Texte à trous : une campagne de prévention]
Complète le texte suivant avec les mots de la liste : \esp{hay que} — \esp{hagas} — \esp{es importante que} — \esp{debes} — \esp{respeta} — \esp{participes}.

\begin{textebox}[Campaña contra el acoso escolar]
Queridos alumnos: el acoso escolar es un problema grave. \_\_\_\_\_\_\_\_ luchar contra él todos juntos. Si ves una injusticia, no \_\_\_\_\_\_\_\_ como si no pasara nada: habla con un profesor. \_\_\_\_\_\_\_\_ a tus compañeros y acepta las diferencias. \_\_\_\_\_\_\_\_ denunciar los malos tratos. Además, \_\_\_\_\_\_\_\_ tú también \_\_\_\_\_\_\_\_ en las actividades de sensibilización del colegio. ¡Juntos podemos construir una escuela segura !
\end{textebox}
\end{exercice}

\begin{exercice}[Appariement : conseils et situations]
Associe chaque situation (colonne A) au conseil qui convient (colonne B).
\begin{center}
\begin{tabularx}{\linewidth}{|X|X|}
\hline
\rowcolor{popblueL} \textbf{A. Situación} & \textbf{B. Consejo} \\
\hline
1. Tu amigo fuma detrás del colegio. & a) \esp{Estudia un poco cada día y no copies.} \\
\hline
2. Tu hermana quiere dejar los estudios. & b) \esp{No fumes : es malo para tu salud.} \\
\hline
3. Un alumno copia en todos los exámenes. & c) \esp{Habla con ella y explícale la importancia de la educación.} \\
\hline
4. Tu vecino bebe demasiado alcohol. & d) \esp{Cuéntaselo a un profesor o al director.} \\
\hline
5. Un compañero sufre acoso escolar. & e) \esp{Debe beber menos y hacer deporte.} \\
\hline
\end{tabularx}
\end{center}
\end{exercice}

\courssub{Je me prépare au Bac}

\begin{exercice}[Compréhension et questions de langue]
Lis le texte suivant, puis réponds aux questions \textbf{en espagnol}, avec des phrases complètes.

\begin{textebox}[Un colegio contra las drogas]
En el colegio San José de Yaoundé, los profesores han organizado una semana de sensibilización contra las conductas de riesgo. Cada mañana, los alumnos asisten a charlas sobre los peligros del alcohol, del tabaco y de las drogas. « No consumas drogas : destruyen tu cuerpo y tu futuro », repite el doctor Mbarga a los jóvenes. Los alumnos también preparan carteles con eslóganes como « Di no a las drogas, di sí a la vida ».

La directora del colegio explica : « Hay que educar a los jóvenes desde muy pronto. Es importante que los padres hablen con sus hijos y que los escuchen con atención. No basta con castigar : hay que comprender y aconsejar ».

Al final de la semana, los estudiantes representan pequeñas obras de teatro. En una de ellas, un joven quiere abandonar los estudios, pero sus amigos le convencen : « No dejes el colegio, piensa en tu futuro. Debes ser fuerte y responsable ». Todos los alumnos aplauden y prometen adoptar comportamientos responsables.
\end{textebox}

\textbf{I. Comprensión del texto}
\begin{enumerate}
\item ¿Qué ha organizado el colegio San José ? ¿Con qué objetivo ?
\item ¿Qué consejo repite el doctor Mbarga a los jóvenes ?
\item Según la directora, ¿qué deben hacer los padres ?
\item ¿Qué ocurre en la obra de teatro de los estudiantes ?
\end{enumerate}

\textbf{II. Questions de langue}
\begin{enumerate}
\item Relève dans le texte deux impératifs négatifs et un impératif affirmatif, puis conjugue ces trois verbes à toutes les personnes de l'impératif (\esp{tú}, \esp{vosotros}, \esp{usted}).
\item Explique pourquoi on dit \esp{« Es importante que los padres hablen »} et non \esp{hablan}.
\item Trouve dans le texte deux mots de la même famille que \esp{responsable} et \esp{educar}.
\end{enumerate}
\end{exercice}

\begin{exercice}[Thème et version]
\textbf{A. Thème.} Traduis en espagnol (10 points) :

« Chers jeunes, il faut dire non aux drogues. Ne sortez pas le soir avec de mauvais amis et ne buvez pas d'alcool. Vous devez travailler à l'école et respecter vos professeurs. Il est important que vous écoutiez les conseils de vos parents. Soyez responsables : votre avenir est entre vos mains. »

\textbf{B. Version.} Traduis en français (10 points) :

\begin{textebox}[Consejos para los jóvenes]
Jóvenes de Camerún: no abandonéis los estudios, porque la educación es la llave del futuro. No os dejéis influenciar por los malos amigos. Hay que respetar a los demás y ayudar a los más débiles. Es importante que participéis en la vida de vuestra comunidad y que habléis con vuestros padres de vuestros problemas. Sed fuertes y responsables: el país os necesita.
\end{textebox}
\end{exercice}

\begin{exercice}[Expression écrite type Bac]
\textbf{Sujet :} \esp{Tu hermano menor quiere dejar los estudios y fuma con sus amigos. Escríbele una carta para aconsejarle.}

Rédige en espagnol une lettre de \textbf{80 à 100 mots}. Tu dois :
\begin{enumerate}
\item utiliser la forme de la lettre (en-tête, salutations, formule finale) ;
\item employer au moins \textbf{deux impératifs affirmatifs} et \textbf{deux impératifs négatifs} ;
\item utiliser les trois structures suivantes : \esp{deber + infinitivo}, \esp{hay que + infinitivo}, \esp{es importante que + subjuntivo} ;
\item expliquer les dangers de la conduite de ton frère et lui proposer des solutions.
\end{enumerate}
\end{exercice}

\newpage

\coursec{Corrigés des exercices}

\begin{corrige}[Exercice 1 — QCM : la bonne forme de l'impératif]
\begin{enumerate}
\item \textbf{b)} \esp{Di} : impératif irrégulier de \esp{decir} à la personne \esp{tú} (forme courte). \esp{Dices} est un présent de l'indicatif, \esp{Decí} n'existe pas (l'impératif \esp{vosotros} est \esp{decid}).
\item \textbf{b)} \esp{hagas} : à l'impératif négatif, on emploie le subjonctif présent : \esp{no hagas}. \esp{Haces} est un indicatif, interdit après \esp{no} à l'impératif.
\item \textbf{b)} \esp{Venga} : pour \esp{usted}, l'impératif affirmatif est identique au subjonctif présent : \esp{venga}. \esp{Ven} correspond à \esp{tú}, \esp{vengas} serait un impératif négatif (\esp{no vengas}).
\item \textbf{a)} \esp{Haced} : impératif \esp{vosotros} de \esp{hacer} (infinitif \esp{hacer} $\rightarrow$ remplacer -r par -d). \esp{Hacéis} est un présent de l'indicatif, \esp{hagan} correspond à \esp{ustedes}.
\item \textbf{b)} \esp{Hazlo} : à l'impératif affirmatif, le pronom se place \emph{après} le verbe et s'y attache ; l'impératif de \esp{hacer} à \esp{tú} est \esp{haz} (et non \esp{haz} avec z ailleurs : \esp{lo haz} est incorrect).
\end{enumerate}
\end{corrige}

\begin{corrige}[Exercice 2 — Vrai ou faux ?]
\begin{enumerate}
\item \textbf{FAUX.} À l'impératif négatif, le pronom se place \emph{avant} le verbe : \esp{No lo hagas} est la seule forme correcte. C'est à l'impératif \emph{affirmatif} que le pronom suit le verbe (\esp{Hazlo}).
\item \textbf{VRAI.} \esp{Hay que + infinitivo} exprime une obligation générale, valable pour tout le monde, sans sujet précis : \esp{Hay que respetar a los demás} « Il faut respecter les autres ».
\item \textbf{FAUX.} Après \esp{es importante que}, on utilise le \emph{subjonctif} présent : \esp{Es importante que hables con tus padres}, et non \esp{hablas}.
\item \textbf{VRAI.} L'impératif de \esp{ir} à la personne \esp{tú} est la forme irrégulière \esp{ve} : \esp{¡Ve a casa !} « Va à la maison ! ».
\item \textbf{VRAI.} \esp{Deber + infinitivo} exprime le devoir, le conseil ou l'obligation morale : \esp{Debes estudiar más} « Tu dois étudier davantage ».
\end{enumerate}
\end{corrige}

\begin{corrige}[Exercice 3 — Vocabulaire : conduites à risque]
\begin{center}
\begin{tabularx}{\linewidth}{|X|X|}
\hline
\rowcolor{popblueL} \textbf{Español} & \textbf{Français} \\
\hline
1. \esp{el acoso escolar} & b) le harcèlement scolaire \\
\hline
2. \esp{el consumo de drogas} & c) la consommation de drogues \\
\hline
3. \esp{la trampa} & a) la triche \\
\hline
4. \esp{el abandono de los estudios} & e) l'abandon des études \\
\hline
5. \esp{el respeto a los demás} & d) le respect des autres \\
\hline
\end{tabularx}
\end{center}
\textbf{Points de vigilance :} attention au faux ami \esp{la trampa}, qui signifie « la triche » (et aussi « le piège »), et non « le tremplin » ; \esp{los demás} signifie « les autres », sans article partitif en espagnol.
\end{corrige}

\begin{corrige}[Exercice 4 — Conjugaison : à l'impératif (tú)]
\begin{enumerate}
\item \esp{Di} la verdad a tus padres. « Dis la vérité à tes parents. » (impératif irrégulier de \esp{decir})
\item No \esp{hagas} trampa en los exámenes. « Ne triche pas aux examens. » (impératif négatif = subjonctif présent)
\item \esp{Ven} aquí, quiero hablar contigo. « Viens ici, je veux parler avec toi. » (irrégulier : \esp{venir} $\rightarrow$ \esp{ven})
\item No \esp{te drogues}, es muy peligroso. « Ne te drogue pas, c'est très dangereux. » (verbe pronominal : pronom \esp{te} placé avant le verbe à la forme négative)
\item \esp{Sé} responsable en la escuela. « Sois responsable à l'école. » (irrégulier : \esp{ser} $\rightarrow$ \esp{sé})
\item No \esp{salgas} por la noche sin permiso. « Ne sors pas le soir sans permission. » (subjonctif de \esp{salir})
\item \esp{Pon} atención en clase. « Fais attention en classe. » (irrégulier : \esp{poner} $\rightarrow$ \esp{pon})
\item No \esp{te vayas} sin despedirte. « Ne t'en va pas sans dire au revoir. » (\esp{irse} : subjonctif \esp{vayas} + pronom \esp{te} devant)
\end{enumerate}
\end{corrige}

\begin{corrige}[Exercice 5 — Thème : traduis en espagnol]
\begin{enumerate}
\item \esp{No fumes, es peligroso para tu salud.} (impératif négatif = subjonctif : \esp{fumes} ; attention : « pour ta santé » = \esp{para tu salud}.)
\item \esp{Debes hacer deporte y hay que respetar a los demás.} (\esp{deber + infinitivo} pour « tu dois » ; \esp{hay que} pour « il faut » ; \esp{hacer deporte} sans article.)
\item \esp{Es importante que hables con tus padres.} (subjonctif obligatoire après \esp{es importante que} ; « parler à quelqu'un » = \esp{hablar con alguien}.)
\item \esp{No dejes la escuela, trabaja duro y escucha a tus profesores.} (impératif négatif \esp{no dejes} ; impératifs affirmatifs réguliers \esp{trabaja} et \esp{escucha} ; \esp{duro} est ici un adverbe invariable.)
\end{enumerate}
\textbf{Points de vigilance :} ne pas calquer le français « ne quitte pas » par \esp{no quites} : on dit \esp{dejar} ou \esp{abandonar} pour « quitter l'école ». Ne jamais mettre d'indicatif après \esp{es importante que}.
\end{corrige}

\begin{corrige}[Exercice 6 — Transformation : du conseil à l'interdiction]
\begin{enumerate}
\item ¡No \esp{fumes} ! (subjonctif présent de \esp{fumar})
\item ¡No \esp{hagas} trampa ! (subjonctif de \esp{hacer}, radical irrégulier \esp{hag-})
\item ¡No \esp{digas} mentiras ! (subjonctif de \esp{decir} : \esp{digas})
\item ¡No \esp{bebas} alcohol ! (subjonctif de \esp{beber})
\item ¡No \esp{abandones} los estudios ! (subjonctif de \esp{abandonar})
\end{enumerate}
\textbf{Rappel de méthode :} pour passer de \esp{no debes + infinitivo} à l'impératif négatif, on supprime \esp{debes} et on met le verbe au subjonctif présent, 2\up{e} personne du singulier : \esp{no + subjonctif}.
\end{corrige}

\begin{corrige}[Exercice 7 — Texte à trous : une campagne de prévention]
Texte complété :

\esp{Queridos alumnos: el acoso escolar es un problema grave. \textbf{Hay que} luchar contra él todos juntos. Si ves una injusticia, no \textbf{hagas} como si no pasara nada: habla con un profesor. \textbf{Respeta} a tus compañeros y acepta las diferencias. \textbf{Debes} denunciar los malos tratos. Además, \textbf{es importante que} tú también \textbf{participes} en las actividades de sensibilización del colegio. ¡Juntos podemos construir una escuela segura !}

\textbf{Traduction :} « Chers élèves : le harcèlement scolaire est un problème grave. Il faut lutter contre lui tous ensemble. Si tu vois une injustice, ne fais pas comme si de rien n'était : parle à un professeur. Respecte tes camarades et accepte les différences. Tu dois dénoncer les mauvais traitements. De plus, il est important que toi aussi tu participes aux activités de sensibilisation du collège. Ensemble, nous pouvons construire une école sûre ! »

\textbf{Justifications :} \esp{hay que} introduit l'obligation générale ; \esp{no hagas} est l'impératif négatif ; \esp{respeta} est l'impératif affirmatif régulier ; \esp{debes} exprime le conseil personnel ; \esp{es importante que} appelle le subjonctif \esp{participes}.
\end{corrige}

\begin{corrige}[Exercice 8 — Appariement : conseils et situations]
\begin{center}
\begin{tabularx}{\linewidth}{|X|X|}
\hline
\rowcolor{popblueL} \textbf{A. Situación} & \textbf{B. Consejo} \\
\hline
1. Tu amigo fuma detrás del colegio. & b) \esp{No fumes : es malo para tu salud.} \\
\hline
2. Tu hermana quiere dejar los estudios. & c) \esp{Habla con ella y explícale la importancia de la educación.} \\
\hline
3. Un alumno copia en todos los exámenes. & a) \esp{Estudia un poco cada día y no copies.} \\
\hline
4. Tu vecino bebe demasiado alcohol. & e) \esp{Debe beber menos y hacer deporte.} \\
\hline
5. Un compañero sufre acoso escolar. & d) \esp{Cuéntaselo a un profesor o al director.} \\
\hline
\end{tabularx}
\end{center}
\textbf{Remarques grammaticales :} en b), impératif négatif \esp{no fumes} ; en c), \esp{explícale} : impératif affirmatif + pronom attaché, avec accent écrit pour conserver la prononciation ; en d), \esp{cuéntaselo} : deux pronoms attachés à l'impératif (\esp{se} + \esp{lo} $\rightarrow$ \esp{selo}), d'où l'accent sur \esp{cuéntaselo}.
\end{corrige}

\begin{corrige}[Exercice 9 — Compréhension et questions de langue]
\textbf{I. Comprensión del texto} (réponses modèles)
\begin{enumerate}
\item \esp{El colegio San José ha organizado una semana de sensibilización contra las conductas de riesgo, con el objetivo de informar a los alumnos sobre los peligros del alcohol, del tabaco y de las drogas.}
\item \esp{El doctor Mbarga repite a los jóvenes que no consuman drogas, porque destruyen su cuerpo y su futuro.}
\item \esp{Según la directora, los padres deben hablar con sus hijos y escucharlos con atención, porque no basta con castigar: hay que comprender y aconsejar.}
\item \esp{En la obra de teatro, un joven quiere abandonar los estudios, pero sus amigos le convencen de no dejar el colegio y de pensar en su futuro.}
\end{enumerate}

\textbf{II. Questions de langue}
\begin{enumerate}
\item Impératifs négatifs : \esp{no consumas} (\esp{consumir}), \esp{no dejes} (\esp{dejar}). Impératif affirmatif : \esp{di} (\esp{decir}), ou \esp{piensa} (\esp{pensar}).

Conjugaison à l'impératif :
\begin{center}
\begin{tabular}{|l|l|l|l|}
\hline
\rowcolor{popblueL} \textbf{Verbe} & \textbf{tú} & \textbf{vosotros} & \textbf{usted} \\
\hline
\esp{consumir} & \esp{consume / no consumas} & \esp{consumid / no consumáis} & \esp{consuma / no consuma} \\
\hline
\esp{dejar} & \esp{deja / no dejes} & \esp{dejad / no dejéis} & \esp{deje / no deje} \\
\hline
\esp{decir} & \esp{di / no digas} & \esp{decid / no digáis} & \esp{diga / no diga} \\
\hline
\end{tabular}
\end{center}

\item On dit \esp{« Es importante que los padres hablen »} parce que le tour impersonnel \esp{es importante que} exprime un souhait, une nécessité : il exige le \emph{subjonctif} présent (\esp{hablen}), et non l'indicatif (\esp{hablan}), qui exprimerait un fait constaté.

\item Famille de \esp{responsable} : \esp{responsables} / \esp{comportamientos responsables} (on accepte \esp{responsable} employé comme adjectif et comme nom). Famille d'\esp{educar} : \esp{educación} n'apparaît pas, mais on trouve \esp{educar} dans \esp{« Hay que educar a los jóvenes »} ; autre mot de la famille possible selon le texte : \esp{sensibilización} n'appartient pas à cette famille. Réponses attendues : \esp{responsables} (famille de \esp{responsable}) et \esp{educar} $\rightarrow$ \esp{educación} n'étant pas présent, on retient \esp{educar} et le nom \esp{educación} n'est pas dans le texte ; la réponse correcte est donc \esp{responsables} et, pour \esp{educar}, le verbe lui-même conjugué \esp{educar} n'ayant pas de dérivé présent, on accepte \esp{educación} si l'élève le justifie par la famille, sinon \esp{sensibilización / sensibilizar} n'est pas valable. \textbf{Corrigé retenu :} \esp{responsables} (famille de \esp{responsable}) et \esp{educar} $\rightarrow$ \esp{educación} (famille du verbe \esp{educar}, mot admis comme dérivé connu) ; dans le texte strictement : \esp{responsables} et \esp{educar}.
\end{enumerate}
\textbf{Barème indicatif (20 points) :} compréhension 8 pts (2 pts par réponse complète et exacte) ; questions de langue 12 pts (relevé 3 pts, conjugaison 6 pts, explication du subjonctif 2 pts, familles de mots 1 pt).
\end{corrige}

\begin{corrige}[Exercice 10 — Thème et version]
\textbf{A. Thème} — proposition de traduction :

\esp{Queridos jóvenes: hay que decir no a las drogas. No salgáis por la noche con malos amigos y no bebáis alcohol. Debéis trabajar en la escuela y respetar a vuestros profesores. Es importante que escuchéis los consejos de vuestros padres. Sed responsables: vuestro futuro está en vuestras manos.}

\textbf{Points de vigilance :} « il faut » = \esp{hay que} (et non \esp{es necesario de}) ; « ne sortez pas » : impératif négatif \esp{vosotros} = subjonctif \esp{no salgáis} ; « soyez » : impératif irrégulier \esp{sed} ; « entre vos mains » se dit \esp{en vuestras manos}. Barème indicatif : 2 pts par phrase (forme verbale 1 pt, vocabulaire et syntaxe 1 pt).

\textbf{B. Version} — proposition de traduction :

« Jeunes du Cameroun : n'abandonnez pas vos études, car l'éducation est la clé de l'avenir. Ne vous laissez pas influencer par les mauvais amis. Il faut respecter les autres et aider les plus faibles. Il est important que vous participiez à la vie de votre communauté et que vous parliez de vos problèmes avec vos parents. Soyez forts et responsables : le pays a besoin de vous. »

\textbf{Points de vigilance :} \esp{no os dejéis influenciar} : tournure réflexive, « ne vous laissez pas influencer » ; \esp{la llave} = « la clé » (faux ami : pas « la lave ») ; \esp{el país os necesita} = « le pays a besoin de vous » (\esp{necesitar} + pronom personnel, à rendre par « avoir besoin de »). Barème indicatif : 2 pts par phrase.
\end{corrige}

\begin{corrige}[Exercice 11 — Expression écrite type Bac]
\textbf{Proposition de rédaction modèle} (environ 95 mots) :

\esp{Yaoundé, 15 de mayo

Querido hermanito:

Estoy muy preocupado por ti. No fumes con tus amigos: el tabaco destruye tu salud. No dejes los estudios: la educación es la llave del futuro. Debes ser fuerte y pensar en tu porvenir. Hay que trabajar duro para tener éxito en la vida. Es importante que hables con papá y mamá de tus problemas: ellos pueden ayudarte. Ven a verme este fin de semana y hablaremos juntos. Escucha mis consejos y sé responsable.

Te quiere mucho,

Tu hermano}

\textbf{Traduction :} « Yaoundé, le 15 mai. Cher petit frère, je suis très inquiet pour toi. Ne fume pas avec tes amis : le tabac détruit ta santé. N'abandonne pas les études : l'éducation est la clé de l'avenir. Tu dois être fort et penser à ton avenir. Il faut travailler dur pour réussir dans la vie. Il est important que tu parles de tes problèmes à papa et maman : ils peuvent t'aider. Viens me voir ce week-end et nous parlerons ensemble. Écoute mes conseils et sois responsable. Ton frère qui t'aime beaucoup. »

\textbf{Vérification des consignes :}
\begin{itemize}
\item Forme de la lettre : lieu et date, \esp{Querido hermanito}, formule finale \esp{Te quiere mucho}.
\item Impératifs affirmatifs : \esp{ven}, \esp{escucha}, \esp{sé}, \esp{habla} (dans le subjonctif).
\item Impératifs négatifs : \esp{no fumes}, \esp{no dejes}.
\item Structures imposées : \esp{debes ser} (\esp{deber + infinitivo}), \esp{hay que trabajar}, \esp{es importante que hables} (subjonctif).
\end{itemize}

\textbf{Barème indicatif (20 points) :} respect de la forme épistolaire 3 pts ; impératifs correctement employés 6 pts ; les trois structures du conseil 6 pts ; richesse du vocabulaire et cohérence 3 pts ; orthographe et accents 2 pts.

\textbf{Points de vigilance :} ne pas écrire \esp{no fumas} (indicatif) mais \esp{no fumes} ; ne pas oublier le subjonctif après \esp{es importante que} ; soigner les accents (\esp{educación}, \esp{atención}) ; éviter le calque « je suis très inquiet \emph{à propos de} toi » : on dit \esp{estoy preocupado por ti}.
\end{corrige}

\newpage

\coursec{Fiche bilan}

\begin{popfigure}
\begin{center}
\begin{tikzpicture}[
  mindmap,
  grow cyclic,
  every node/.style={concept, font=\small, align=center},
  root concept/.append style={concept color=popblue, font=\small\bfseries, text=white, minimum size=3.2cm},
  level 1/.append style={level distance=3.6cm, sibling angle=60},
  level 1 concept/.append style={font=\footnotesize, minimum size=2.2cm},
]
\node[root concept, align=center]{Conductas de riesgo\\Comportamientos responsables}
  child[concept color=poporange]{ node[align=center]{Lexique\\riesgo, prevención, salud} }
  child[concept color=popgreen]{ node[align=center]{Impératif\\affirmatif\\tú, vosotros, usted} }
  child[concept color=poppurple]{ node[align=center]{Impératif\\négatif\\no + subjonctif} }
  child[concept color=popgold]{ node[align=center]{Conseil\\deber, hay que,\\es importante que} }
  child[concept color=poppink]{ node[align=center]{Méthode\\commentaire\\de texte} }
  child[concept color=popblue!70]{ node[align=center]{Production\\message de\\prévention} };
\end{tikzpicture}
\end{center}
\legende{Carte mentale de la leçon : sensibiliser pour des comportements responsables}
\end{popfigure}

\begin{bilan}
\textbf{1. Lexique clé (espagnol -- français)}
\begin{itemize}
  \item \esp{la conducta de riesgo} : la conduite à risque ; \esp{el comportamiento responsable} : le comportement responsable
  \item \esp{la prevención} : la prévention ; \esp{prevenir} : prévenir ; \esp{evitar} : éviter
  \item \esp{la salud} : la santé ; \esp{la droga / el alcohol / el tabaco} : la drogue / l'alcool / le tabac
  \item \esp{el respeto} : le respect ; \esp{respetar a los demás} : respecter les autres
  \item \esp{la responsabilidad} : la responsabilité ; \esp{la solidaridad} : la solidarité
  \item \esp{proteger(se)} : (se) protéger ; \esp{cuidar} : prendre soin de ; \esp{ayudar} : aider
  \item \esp{la campaña de sensibilización} : la campagne de sensibilisation ; \esp{el cartel / el eslogan} : l'affiche / le slogan
\end{itemize}

\textbf{2. Grammaire à connaître par c\oe ur}
\begin{center}
\begin{tabularx}{\linewidth}{|l|X|X|}
\hline
\rowcolor{popblueL} \textbf{Personne} & \textbf{Impératif affirmatif} & \textbf{Impératif négatif} \\
\hline
tú & \esp{habla, come, escribe} & \esp{no hables, no comas, no escribas} \\
\hline
vosotros & \esp{hablad, comed, escribid} & \esp{no habléis, no comáis, no escribáis} \\
\hline
usted & \esp{hable, coma, escriba} & \esp{no hable, no coma, no escriba} \\
\hline
\end{tabularx}
\end{center}
\begin{itemize}
  \item Impératif \esp{tú} = 3\textsuperscript{e} personne du singulier du présent ; irréguliers : \esp{di, haz, ve, pon, sal, sé, ten, ven}.
  \item Impératif négatif = \esp{no} + subjonctif présent.
  \item Pronoms accolés à l'impératif affirmatif : \esp{levántate, hazlo, cuídate} ; avant le verbe au négatif : \esp{no te levantes, no lo hagas}.
  \item Conseil et obligation : \esp{deber + infinitif} (\esp{Debes respetar a los demás.}), \esp{hay que + infinitif} (\esp{Hay que proteger la salud.}), \esp{es importante que + subjonctif} (\esp{Es importante que cuides tu cuerpo.}).
\end{itemize}

\textbf{3. Méthodes express}
\begin{itemize}
  \item \textbf{Commentaire de texte} : présenter (auteur, source, thème) $\rightarrow$ résumer les idées $\rightarrow$ analyser la langue et les procédés $\rightarrow$ donner son avis argumenté.
  \item \textbf{Message de prévention} : un slogan court à l'impératif + une phrase de conseil (\esp{deber}, \esp{hay que}) + une image mentale forte. Exemple : \esp{Di no a las drogas. Tu futuro te espera.}
\end{itemize}
\end{bilan}

\begin{aretenir}[Checklist avant le Bac]
\begin{itemize}
  \item $\square$ Je connais le vocabulaire des conduites à risque et de la prévention.
  \item $\square$ Je sais former l'impératif affirmatif de \esp{tú}, \esp{vosotros} et \esp{usted}.
  \item $\square$ Je connais par c\oe ur les huit impératifs irréguliers de \esp{tú} : \esp{di, haz, ve, pon, sal, sé, ten, ven}.
  \item $\square$ Je sais que l'impératif négatif se construit avec \esp{no} + subjonctif présent.
  \item $\square$ Je place correctement les pronoms : accolés à l'affirmatif, avant le verbe au négatif.
  \item $\square$ Je sais donner un conseil avec \esp{deber + infinitif} et \esp{hay que + infinitif}.
  \item $\square$ Je sais utiliser \esp{es importante que + subjonctif} sans oublier le subjonctif.
  \item $\square$ Je ne confonds pas \esp{deber} (devoir) et \esp{deber de} (probabilité).
  \item $\square$ Je sais structurer un commentaire de texte en quatre étapes.
  \item $\square$ Je sais rédiger un slogan de prévention court et percutant à l'impératif.
  \item $\square$ Je vérifie les accents : \esp{prevención, salud, respeto, sensibilización}.
  \item $\square$ Je peux jouer un court dialogue de sensibilisation avec un camarade.
\end{itemize}
\end{aretenir}

\findecours

\end{document}
